%%%%%%%%%%%%%%%%%%%%%%%%%%%%%%%%%%%%%%%%%%%%%%%%%%%%%%%%%%%%%%%%%%%%%%
% Lớp 11 | Các số đặc trưng đo xu thế trung tâm cho mẫu số liệu | Số trung bình và mốt của mẫu số liệu ghép nhóm
% Dạng 1D5?1-2: Mẫu số liệu ghép nhóm
% Số câu: 175  (sinh tự động bởi tach_cau.py; từ nay thầy sửa trực tiếp file này)
%%%%%%%%%%%%%%%%%%%%%%%%%%%%%%%%%%%%%%%%%%%%%%%%%%%%%%%%%%%%%%%%%%%%%%

%%%%%%%%%%%%%%%% TRẮC NGHIỆM 4 PHƯƠNG ÁN %%%%%%%%%%%%%%%%

% nguồn: Ok_Toan11_TracNghiem_Lan2
\begin{ex}%[1D5N1-2]%[ma:cc5f312c]
	Độ dài của mỗi nhóm trong mẫu số liệu ghép nhóm được tính  thế nào?
	\choice 
		 {$a-b$}
		 {$a+b$}
		{$ab$} 
		 {\True $b-a$}
	\loigiai{Độ dài của nhóm $[a;b)$ là  $b-a$.}
\end{ex}

% nguồn: Ok_Toan11_TracNghiem_Lan2
\begin{ex}%[1D5N1-2]%[ma:1e7c7497]
	Khảo sát thời gian đọc sách trong ngày của một số học sinh khối $11$  thu được mẫu số liệu ghép nhóm sau
	\begin{center}
		\begin{tabular}{|c|c|c|c|c|c|}
			 \hline
			Thời gian (phút) & $[0;30)$ & $[30;60)$ & $[60;90)$ & $[90;120)$ & $[120;150)$  \\
			\hline 
			Số học sinh & $2$ & $5$ & $12$ & $15$ & $6$ \\ 
			 \hline
		\end{tabular}
	\end{center}
	Giá  trị đại diện của nhóm $[30;60)$ là
	\choice 
		 {60}
		 {40}
		{\True  45}
		{30} 
	\loigiai{Giá trị đại diện của nhóm $[30;60)$ là $\dfrac{30+60}{2}=45$.} 
\end{ex}

% nguồn: Ok_Toan11_TracNghiem_Lan2
\begin{ex}%[1D5N1-2]%[ma:7b5136a0]
	Để chuyển mẫu số liệu không ghép nhóm sang mẫu số liệu ghép nhóm ta  làm thế nào?
	\choice 
		{\True Bước 1: Chia miền giá trị của mẫu số liệu thành một số nhóm theo tiêu  chí cho trước.
		\\ Bước 2: Đếm số giá trị của  mẫu số liệu thuộc mỗi nhóm}
		{Bước 1: Đếm số giá trị giống nhau của mẫu số liệu. \\ Bước 2: Lập bảng thống  kê cho mẫu số liệu ghép nhóm}
		{Cả  A, B đều sai}
		{Cả A,  B đều đúng}
	\loigiai{Để chuyển mẫu số liệu không ghép nhóm sang mẫu số liệu ghép nhóm ta làm  như sau
		\begin{itemize}
			\item Bước 1:  Chia miền giá trị của mẫu số liệu thành một số nhóm theo tiêu chí cho trước.
			\item Bước 2:  Đếm số giá trị của mẫu số liệu thuộc mỗi nhóm và lập bảng thống kê cho mẫu số liệu ghép nhóm.
	\end{itemize}}
\end{ex}

% nguồn: Ok_Toan11_TracNghiem_Lan2
\begin{ex}%[1D5N1-2]%[ma:27e104d6]
	Cho mẫu số liệu ghép nhóm về thống kê nhiệt độ tại một địa điểm trong $40$ ngày, ta có bảng số liệu sau 
	\begin{center}
		\begin{tabular}{|c|c|c|c|c|c|}
			 \hline
			Nhiệt độ ($^\circ C$) &  $[19;22)$ & $[22;25)$ & $[25;28)$ & $[28;31)$ \\
			 \hline
			Số ngày & $7$ & $15$ &  $12$ & $6$ \\ \hline
		\end{tabular}
	\end{center}
	Có bao nhiêu ngày có nhiệt  độ từ $28^\circ$C đến dưới $31^\circ$C?
	 \choice
		{$4$} 
		{$5$} 
		{\True  $6$}
		 {$7$}
	\loigiai{Có $6$ ngày có nhiệt độ từ $28^\circ$C đến dưới  $31^\circ$C.}
\end{ex}

% nguồn: Ok_Toan11_TracNghiem_Lan2
\begin{ex}%[1D5N1-2]%[ma:54a80c26]
	Điều tra về chiều cao của học sinh khối lớp $11$, ta được mẫu số liệu  sau
	\begin{center}
		\begin{tabular}{|c|c|}
			 \hline
			Chiều cao (cm) &  Số học sinh \\
			\hline 
			$[150;152)$ & $5$ \\ 
			 \hline
			$[152;154)$ &  $18$ \\
			 \hline
			$[154;156)$ & $40$  \\
			 \hline
			 $[156;158)$ & $26$ \\
			\hline 
			$[158;160)$ & $8$ \\ 
			\hline 
			$[160;162)$ & $3$ \\ 
			 \hline
			Tổng  & $N=100$ \\
			 \hline
		\end{tabular}
	\end{center}
	Hỏi có bao nhiêu học sinh có  chiều cao từ $156$ cm trở lên?
	 \choice
		{$5$} 
		 {\True $37$}
		{$7$} 
		 {$12$}
	\loigiai{Số học sinh có  chiều cao từ $156$ cm trở lên là $26+8+3=37$ học sinh.}
\end{ex}

% nguồn: Ok_Toan11_TracNghiem_Lan2
\begin{ex}%[1D5N1-2]%[ma:83a7db00]
	 Điều tra về chiều cao của học sinh khối lớp $11$, ta có kết quả sau
	\begin{center}
		\begin{tabular}{|c|c|c|}
			 \hline
			Nhóm & Chiều cao  (cm) & Số học sinh \\
			 \hline
			$1$ & $[150;152)$ & $5$  \\
			\hline 
			$2$ & $[152;154)$  & $18$ \\
			 \hline
			$3$ & $[154;156)$  & $40$ \\
			\hline 
			$4$  & $[156;158)$ & $26$ \\
			\hline 
			$5$ & $[158;160)$ & $8$  \\
			\hline 
			$6$ & $[160;162)$ & $3$  \\
			\hline 
			\multicolumn{2}{|c|}{} & $N=100$  \\
			\hline 
		\end{tabular}
	\end{center}
	 Giá trị đại diện của nhóm thứ tư là?
	\choice 
		{$156{,}5$} 
		{\True$157$} 
		 {$157{,}5$}
		{$158$} 
	\loigiai{Giá trị đại diện của nhóm thứ tư là  $\dfrac{156+158}{2}=157$.}
\end{ex}

% nguồn: Ok_Toan11_TracNghiem_Lan2
\begin{ex}%[1D5N1-2]%[ma:e78278a2]
	Mẫu số liệu $(T)$ được mô tả dưới dạng bảng thống kê  sau
	\begin{center}
		\begin{tabular}{|c|c|c|c|c|c|c|}
			\hline 
			Tổng điểm & $< 6$ &  $[6;7)$ & $[7;8)$ & $\cdots$ & $[28;29)$ & $[29;30]$ \\
			 \hline
			Số thí sinh & $23$ & $69$ & $192$ & $\cdots$ & $126$ & $12$  \\ \hline
		\end{tabular}
	\end{center}
	Tổng số giá trị của mẫu  số liệu là
	 \choice
		{$2589863$} 
		 {$14586632$}
		{$125863$} 
		{\True  $1379008$}
	\loigiai{Tổng số giá trị của mẫu số liệu là $344752 \times 4 = 1379008$ giá trị.} 
\end{ex}

% nguồn: Ok_Toan11_TracNghiem_Lan2
\begin{ex}%[1D5N1-2]%[ma:d053b1b4]
	Đo chiều cao (tính bằng: cm) của $500$ học sinh trong một trường THPT  ta thu được kết quả như sau
	\begin{center}
		\begin{tabular}{|c|c|c|c|c|c|}
			\hline 
			Chiều  cao & $[150;154)$ & $[154;158)$ & $[158;162)$ & $[162;166)$ & $[166;170)$ \\
			\hline 
			 Số học sinh & $25$ & $50$ & $200$ & $175$ & $50$ \\
			\hline 
		\end{tabular}
	\end{center}
	Mẫu số liệu ghép nhóm đã  cho có tất cả bao nhiêu nhóm?
	 \choice
		{\True  $5$}
		{$6$} 
		 {$7$}
		{$12$} 
	\loigiai{Mẫu số liệu ghép nhóm đã cho  có tất cả $5$ nhóm.}
\end{ex}

% nguồn: Ok_Toan11_TracNghiem_Lan2
\begin{ex}%[1D5N1-2]%[ma:6aa21c39]
	Mẫu số liệu $(T)$ được mô  tả dưới dạng bảng thống kê sau
	\begin{center}
		\begin{tabular}{|c|c|c|c|c|c|c|}
			\hline 
			Tổng điểm & $< 6$ & $[6;7)$ & $[7;8)$ &  $\cdots$ & $[28;29)$ & $[29;30]$ \\
			\hline 
			Số thí sinh & $23$ & $69$ & $192$ & $\cdots$ & $\cdots$ & $12$ \\  \hline
		\end{tabular}
	\end{center}
	Số lượng thí sinh  có ít nhất một môn học có điểm từ $6$ đến dưới $7$ là
	 \choice
		 {$23$}
		{$192$} 
		{$56$} 
		{\True $69$} 
	\loigiai{Số lượng thí sinh có ít nhất một môn học có điểm từ $6$ đến dưới $7$ là  $69$.}
\end{ex}

% nguồn: Ok_Toan11_TracNghiem_Lan2
\begin{ex}%[1D5N1-2]%[ma:ed84c9a0]
	Thời gian truy cập Internet mỗi buổi  tối của một số học sinh được cho trong bảng sau
	\begin{center}
		\begin{tabular}{|c|c|c|c|c|c|c|}
			 \hline
			Thời gian (phút) & $[9{,}5;12{,}5)$ & $[12{,}5;15{,}5)$ & $[15{,}5;18{,}5)$  & $[18{,}5;21{,}5)$ & $[21{,}5;24{,}5)$ \\
			 \hline
			 Số học sinh & $3$ & $12$ & $15$ & $24$ & $2$\\ \hline
		\end{tabular}
	\end{center}
	Có bao nhiêu học sinh truy cập Internet mỗi buổi tối có thời gian từ $18{,}5$ phút đến dưới $21{,}5$ phút? 
	\choice 
		{\True  $24$}
		 {$15$}
		 {$2$}
		 {$20$}
	\loigiai{Số  học sinh truy cập Internet mỗi buổi tối có thời gian từ $18{,}5$ phút đến dưới $21{,}5$ phút là $24$.
	} 
\end{ex}

% nguồn: Ok_Toan11_TracNghiem_Lan2
\begin{ex}%[1D5N1-2]%[ma:bf195b1c]
	Khảo sát  tổng thời gian truy cập internet mỗi tối (đơn vị: phút) của một số học sinh thu được mẫu số liệu ghép nhóm sau
	\begin{center}
		\begin{tabular}{|c|c|c|c|c|c|c|}
			 \hline
			Thời gian (phút) & $[0;30)$ & $[30;60)$ & $[60;90)$ & $[90;120)$ &  $[120;150)$ & $[150;180)$ \\
			\hline 
			Số học sinh & $11$  & $25$ & $36$ & $15$ & $8$ & $5$ \\
			 \hline
		\end{tabular}
	\end{center}
	Xác định cỡ mẫu của  mẫu số liệu trên
	 \choice
		{90} 
		{\True 100} 
		{105} 
		 {110}
	\loigiai{Cỡ mẫu  của mẫu số liệu là $11+25+36+15+8+5=100$.}
\end{ex}

% nguồn: Ok_Toan11_TracNghiem_Lan2
\begin{ex}%[1D5N1-2]%[ma:050a490a]
	Khảo sát nhiệt độ  ($^\circ \text{C}$) tại một địa điểm trong $40$ ngày thu được mẫu số liệu ghép nhóm sau
	\begin{center}
		\begin{tabular}{|c|c|c|c|c|c|}
			\hline 
			Nhiệt độ ($^\circ \text{C}$)  & $[20;23)$ & $[23;26)$ & $[26;29)$ & $[29;32)$ & $[32;35)$ \\
			\hline 
			Số ngày & $7$ & $6$ & $10$ & $8$ &  $9$ \\
			 \hline
		\end{tabular}
	\end{center}
	Số ngày có nhiệt độ từ $29^\circ \text{C}$ đến dưới $32^\circ \text{C}$  là
	 \choice
		 {\True 8}
		 {40}
		{10} 
		{9} 
	\loigiai{Số ngày có nhiệt độ từ $29^\circ \text{C}$  đến dưới $32^\circ \text{C}$ là $8$ ngày.}
\end{ex}

% nguồn: Ok_Toan11_TracNghiem_Lan2
\begin{ex}%[1D5N1-2]%[ma:7c8719dc]
	Đo cân nặng của một số học sinh lớp $11\text{D}$ cho trong bảng sau 
	\begin{center}
		\begin{tabular}{|c|c|c|c|c|c|c|}
			 \hline
			Cân nặng (kg) & $[40{,}5; 45{,}5)$ & $[45{,}5; 50{,}5)$  & $[50{,}5; 55{,}5)$ & $[55{,}5; 60{,}5)$ & $[60{,}5; 65{,}5)$ & $[65,5; 70,5)$ \\
			\hline 
			Số học sinh & $10$ & $7$ & $16$  & $4$ & $2$ & $3$ \\
			\hline 
		\end{tabular}
	\end{center}
	Giá trị đại  diện của nhóm $[60{,}5; 65{,}5)$ là
	\choice 
		 {$55{,}5$}
		{$58$} 
		 {$60{,}5$}
		 {$5$}
	\loigiai{Trong mỗi khoảng cân nặng, giá trị đại diện trung bình cộng của giá trị hai đầu mút nên  ta có bảng sau
		\begin{center}
			\begin{tabular}{|c|c|c|c|c|c|c|}
				 \hline
				Cân nặng (kg) & $[40{,}5; 45{,}5)$ & $[45{,}5; 50{,}5)$ & $[50{,}5; 55{,}5)$ & $[55{,}5; 60{,}5)$ & $[60{,}5; 65{,}5)$ & $[65{,}5; 70{,}5)$ \\ 
				\hline 
				Giá trị đại diện & $43$ & $48$ & $53$ & $58$ &  $63$ & $68$ \\
				 \hline
				Số họ sinh  & $10$ & $7$ & $16$ & $4$ & $2$ & $3$ \\
				\hline 
			\end{tabular}
		\end{center}
	} 
\end{ex}

% nguồn: Ok_Toan11_TracNghiem_Lan2
\begin{ex}%[1D5N1-2]%[ma:cd8c90fe]
	Một công ty may quần áo  đồng phục học sinh cho biết cỡ áo theo chiều cao của học sinh được tính như sau
	\begin{center}
		\begin{tabular}{|c|c|c|c|c|c|}
			 \hline
			Chiều cao (cm) & $[150;160)$ & $[160;167)$ & $[167;170)$ & $[170;175)$  & $[175;180)$ \\
			 \hline
			Cỡ áo & S & M & L & XL & XXL \\  \hline
		\end{tabular}
	\end{center}
	Công ty muốn ước lượng tỷ lệ các cỡ áo khi may cho học sinh lớp $11$ đã đo chiều cao của $36$ học sinh khối $11$  của một trường và thu được mẫu số liệu sau (đơn vị là cm)
	\begin{center}
		160;  161; 161; 162; 162; 162; 163; 163; 163 \\
		164; 164; 164; 164;  165; 165; 165; 165; 165 \\
		166; 166;  166; 166; 167; 167; 168; 168; 168 \\
		168; 169; 169; 170; 171; 171; 172;  172; 174.
	\end{center}
	Số học sinh có chiều cao vừa vặn với size  L là bao nhiêu?
	\choice 
		{\True  $8$}
		{$7$} 
		{$6$} 
		 {$5$}
	\loigiai{Số học sinh có chiều cao vừa vặn  với size L là $8$ vì $167; 167; 168; 168; 168; 169; 169 \in [167;170)$.}
\end{ex}

% nguồn: Ok_Toan11_TracNghiem_Lan2
\begin{ex}%[1D5N1-2]%[ma:a35d8eed]
	\immini[thm]{Điều tra về chiều cao của  học sinh khối lớp $11$, ta được mẫu số liệu. Mẫu số liệu ghép nhóm đã cho có bao nhiêu nhóm?
	\choice 
		{$5$} 
		 {\True $6$}
		 {$7$}
		{$12$}} 
		{\begin{tabular}{|c|c|}
		\hline 
		Chiều  cao & Số học sinh \\
		 \hline
		$[150;152)$  & $5$ \\
		 \hline
		$[152;154)$ & $18$  \\
		\hline 
		$[154;156)$ & $40$ \\ 
		 \hline
		 $[156;158)$ & $26$ \\
		 \hline
		$[158;160)$ &  $8$ \\
		\hline 
		$[160;162)$ & $3$ \\  \hline
\end{tabular}}
\loigiai{Có  $6$ nhóm số liệu ở bảng trên là $[150;152)$, $[152;154)$, $[154;156)$, $[156;158)$, $[158;160)$, $[160;162)$.}
\end{ex}

% nguồn: Ok_Toan11_TracNghiem_Lan2
\begin{ex}%[1D5N1-2]%[ma:c284a626]
	Điền từ thích hợp và chỗ trống: Mẫu số liệu $\cdots$ là mẫu số liệu  cho dưới dạng bảng tần số của các nhóm số liệu.
	 \choice
		{Không ghép nhóm} 
		{\True Ghép nhóm} 
		{Ghép nhóm và không  ghép nhóm}
		{Cả ba câu trên đều  sai}
	\loigiai{Mẫu số  liệu ghép nhóm là mẫu số liệu cho dữ liệu dạng bảng tần số của các nhóm số liệu.}
\end{ex}

% nguồn: Ok_Toan11_TracNghiem_Lan2
\begin{ex}%[1D5N1-2]%[ma:0fde87f4]
	Mỗi nhóm  số liệu ghép nhóm là tập hợp gồm
	 \choice
		{Các giá trị của số liệu được  ghép nhóm theo nhiều tiêu chí xác định}
		 {Các giá trị của số liệu được ghép nhóm theo hai tiêu chí xác định}
		{\True Các giá trị số liệu được ghép  nhóm theo một tiêu chí xác định}
		{Các giá trị số liệu được ghép nhóm theo ba tiêu  chí xác định}
	\loigiai{Mỗi nhóm số  liệu là tập hợp các giá trị của số liệu được ghép nhóm theo một tiêu chí xác định.}
\end{ex}

% nguồn: Ok_Toan11_TracNghiem_Lan2
\begin{ex}%[1D5N1-2]%[ma:5ae4cfd4]
	Khảo sát thời gian xem ti vi trong một ngày của một số học  sinh khối $11$ thu được mẫu số liệu ghép nhóm sau
	\begin{center}
		\begin{tabular}{|c|c|c|c|c|c|}
			 \hline
			 Thời gian (phút) & $[0;20)$ & $[20;40)$ & $[40;60)$ & $[60;80)$ & $[80;100)$ \\
			 \hline
			Sô học sinh & $5$ & $9$ & $12$ &  $10$ & $6$ \\
			 \hline
		\end{tabular}
	\end{center}
	 Số học sinh xem ti vi từ $60$ phút đến dưới $80$ phút là
	 \choice
		 {9}
		{\True 10} 
		 {12}
		 {5}
	\loigiai{ 
	Số  học sinh xem ti vi từ $60$ phút đến dưới $80$ phút là $10$.}
\end{ex}

% nguồn: Ok_Toan11_TracNghiem_Lan2
\begin{ex}%[1D5N1-2]%[ma:ca749754]
	Tìm hiểu thời gian xem tivi trong tuần trước (đơn vị: giờ)  của một số học sinh thu được kết quả sau.
	\begin{table}[h]
		 \centering
		\begin{tabular}{|c|c|c|c|c|c|c|}
			\hline 
			Thời gian  (giờ) & $[0;5)$ & $[5;10)$ & $[10;15)$ & $[15;20)$ & $[20;25)$ \\
			 \hline
			Số học sinh & $8$ & $16$  & $4$ & $2$ & $2$ \\ \hline
		\end{tabular}
	\end{table}\\
	Có bao nhiêu học sinh có thời gian xem tivi từ $20$ giờ đến dưới $25$  giờ trong tuần trước?
	 \choice
		 {$5$}
		 {$4$}
		{$3$} 
		{\True $2$} 
	\loigiai{Từ bảng số liệu ta thấy có 2 học sinh có thời gian xem  tivi từ $20$ giờ đến dưới $25$ giờ trong tuần trước.}
\end{ex}

% nguồn: Ok_Toan11_TracNghiem_Lan2
\begin{ex}%[1D5N1-2]%[ma:9c02cc62]
	Nhóm số liệu ghép nhóm thường được  cho dưới dạng
	\choice 
		{$(a;b]$, trong đó $a$ là đầu mút trái, $b$ là đầu mút phải} 
		{$(a;b)$, trong đó $a$ là đầu mút trái,  $b$ là đầu mút phải}
		{$(a;b]$, trong đó $a$ là đầu mút trái, $b$ là đầu mút phải} 
		{\True $[a;b)$, trong đó $a$ là  đầu mút trái, $b$ là đầu mút phải}
	\loigiai{Nhóm số liệu thường được  cho dưới dạng $[a;b)$, trong đó $a$ là đầu mút trái, $b$ là đầu mút phải.}
\end{ex}

% nguồn: Ok_Toan11_TracNghiem_Lan2
\begin{ex}%[1D5N1-2]%[ma:af5c6b2e]
	Chọn đáp án sai trong các đáp  án sau
	\choice 
		 {\True Nên chia thành nhiều nhóm}
		{Mẫu  số liệu ghép nhóm được dùng khi ta không thể thu thập được số liệu chính xác}
		{Trong một số trường hợp,  nhóm số liệu cuối cùng có thể lấy đầu mút bên phải}
		{Độ dài của  nhóm là $b-a$}
	\loigiai{Ta không nên chia thành quá nhiều nhóm hoặc quá  ít nhóm. Câu A sai.}
\end{ex}

% nguồn: Ok_Toan11_TracNghiem_Lan2
\begin{ex}%[1D5N1-2]%[ma:5a2786c5]
	Cho mẫu số liệu ghép nhóm về số tiền mà sinh viên chi cho thanh  toán cước điện thoại trong tháng
	\begin{center}
		\begin{tabular}{|c|c|c|c|c|c|c|}
			 \hline
			Số tiền  (nghìn đồng) & $[0;50)$ & $[50;100)$ & $[100;150)$ & $[150;200)$ & $[200;250)$ \\
			 \hline
			Số sinh viên & $5$ & $12$ & $23$  & $17$ & $3$ \\ \hline
		\end{tabular}
	\end{center}
	Có bao nhiêu sinh viên chi từ $100$  đến dưới $150$ nghìn đồng cho việc thanh toán cước điện thoại trong tháng?
	 \choice
		 {$5$}
		 {\True $23$}
		{$12$} 
		 {$17$}
	 \loigiai{Có $23$ sinh viên chi từ $100$ đến dưới $150$ nghìn đồng cho việc thanh toán cước điện thoại trong tháng.}
\end{ex}

% nguồn: Ok_Toan11_TracNghiem_Lan2
\begin{ex}%[1D5N1-2]%[ma:a9f00f0a]
	Tìm hiểu thời gian xem tivi trong tuần trước  (đơn vị: giờ) của một số học sinh thu được kết quả sau
	\begin{center}
		\begin{tabular}{|c|c|c|c|c|c|}
			 \hline
			Thời gian (giờ) &  $[0; 5)$ & $[5; 10)$ & $[10; 15)$ & $[15; 20)$ & $[20; 25)$ \\
			\hline 
			Số học  sinh & $8$ & $16$ & $4$ & $2$ & $2$ \\
			\hline 
		\end{tabular}
	\end{center}
	Giá trị đại diện của  nhóm $[20; 25)$ là
	\choice 
		{\True $22{,}5$} 
		{$23$} 
		{$0$} 
		{$5$} 
	\loigiai{Giá trị đại diện  của nhóm $[20;25)$ là $\dfrac{20+25}{2}=22{,}5$.}
\end{ex}

% nguồn: Ok_Toan11_TracNghiem_Lan2
\begin{ex}%[1D5N1-2]%[ma:59b92939]
	Đo chiều cao (tính bằng cm) của $500$ học sinh trong một trường THPT ta thu được kết quả  như sau
	\begin{center}
		\begin{tabular}{|c|c|c|c|c|c|}
			 \hline
			Chiều cao&$[150;154)$  & $[154;158)$ & $[158;162)$ & $[162;166)$ & $[166;170)$\\
			\hline 
			Số học sinh &$25$ &$50$ &  $200$ & $175$ & $50$\\ \hline
		\end{tabular}
	\end{center}
	Giá trị đại diện của  nhóm $[162;166)$ là
	\choice 
		 {$162$}
		{\True $164$} 
		{$166$} 
		 {$4$}
	\loigiai{  Ta có bảng sau
		\begin{center}
			\begin{tabular}{|c|c|c|}
				 \hline
				Lớp chiều cao& Giá trị đại diện  & Số học sinh \\
				\hline 
				$[150;154)$ & $152$ & $25$\\ 
				\hline 
				$[154;158)$ & $156$ &  $50$\\
				 \hline
				 $[158;162)$ & $160$ & $200$\\
				 \hline
				$[162;166)$ & $164$ & $175$\\ 
				 \hline
				$[166;170)$  & $168$ & $50$\\ \hline
			\end{tabular}
		\end{center}
	} 
\end{ex}

% nguồn: Ok_Toan11_TracNghiem_Lan2
\begin{ex}%[1D5N1-2]%[ma:ba80111c]
	Cho mẫu số liệu ghép nhóm về thời gian (phút) đi từ nhà đến nơi làm việc  của các nhân viên một công ty như sau
	\begin{center}
		\begin{tabular}{|c|c|c|c|c|c|c|c|}
			\hline 
			Thời gian & $[15;20)$ & $[20;25)$ & $[25;30)$ & $[30;35)$  & $[35;40)$ & $[40;45)$ & $[45;50)$ \\
			 \hline
			Số nhân viên & $6$ & $14$ & $25$ &  $37$ & $21$ & $13$ & $9$ \\ \hline
		\end{tabular}
	\end{center}
	Mẫu số liệu được  chia thành bao nhiêu nhóm?
	\choice 
		{$6$ nhóm} 
		 {$5$ nhóm}
		{\True  $7$ nhóm}
		{$8$  nhóm}
	\loigiai{Mẫu số liệu được chia  thành $7$ nhóm.}
\end{ex}

% nguồn: Ok_Toan11_TracNghiem_Lan2
\begin{ex}%[1D5N1-2]%[ma:54460d54]
	Cho bảng khảo sát về cân nặng học sinh trong  lớp.
	\begin{table}[h]
		 \centering
		\begin{tabular}{|c|c|c|c|c|c|c|}
			 \hline
			Cân nặng (kg) & [45;50) & [50;55) &  [55;60) & [60;65) & [65;70) \\
			 \hline
			Số học sinh & $2$ & $14$ & $11$ &  $10$ & $3$ \\ \hline
		\end{tabular}
	\end{table}\\
	Khoảng cân nặng mà số học sinh chiếm nhiều nhất  là
	 \choice
		 {$[60;65)$}
		{$[55;60)$} 
		{\True  $[50;55)$}
		{$[60;65)$} 
	\loigiai{Có  $2$ học sinh có cân nặng từ $45$ kg đến dưới $50$ kg.\\
		Có $14$ học sinh có cân nặng từ $50$ kg đến dưới $55$  kg.\\
		Có $11$ học sinh có  cân nặng từ $55$ kg đến dưới $60$ kg.\\
		Có $10$ học sinh  có cân nặng từ $60$ kg đến dưới $654$ kg.\\
		Có $3$  học sinh có cân nặng từ $65$ kg đến $70$ kg.\\
	Vậy khoảng cân nặng từ $50$ kg đến dưới $55$ kg chiếm nhiều học  sinh nhất.}
\end{ex}

% nguồn: Ok_Toan11_TracNghiem_Lan2
\begin{ex}%[1D5N1-2]%[ma:3baa8389]
	Mẫu số liệu ghép nhóm được dùng khi nào? 
	\choice 
		{Khi ta có thể thu thập được  số liệu chính xác hoặc do yêu cầu của bài toán mà ta phải biểu diễn mẫu số liệu dưới dạng ghép nhóm để thuận lợi cho việc tổ chức, đọc và phân tích số liệu}
		{Khi ta không thể thu thập được số liệu chính xác do yêu cầu của bài toán mà ta  phải biểu diễn mẫu số liệu dưới dạng ghép nhóm để thuận lợi cho việc phân tích số liệu}
		{\True Khi ta không thể thu thập được số liệu chính xác hoặc do yêu cầu của bài toán mà ta phải  biểu diễn mẫu số liệu dưới dạng ghép nhóm để thuận lợi cho việc tổ chức, đọc và phân tích số liệu}
		{Cả  ba câu trên đều sai}
	\loigiai{Mẫu  số liệu ghép nhóm được dùng khi ta không thể thu thập được số liệu chính xác hoặc do yêu cầu của bài toán mà ta phải biểu diễn mẫu số liệu dưới dạng ghép nhóm để thuận lợi cho việc tổ chức, đọc và phân tích số liệu}
\end{ex}

% nguồn: Ok_Toan11_TracNghiem_Lan2
\begin{ex}%[1D5N1-2]%[ma:b1de4dee]
	Chiều cao (đơn vị: cm) của số học sinh lớp $11\text{A}$ được cho trong bảng dưới đây: 
	\begin{center}
		\begin{tabular}{|c|c|c|c|c|c|c|}
			 \hline
			Chiều cao (cm) & $[150;155)$ & $[155;160)$ & $[160;165)$ & $[165;170)$ & $[170;175)$ & $[175;180)$  \\
			 \hline
			Số học sinh  & $4$ & $8$ & $12$ & $10$ & $6$ & $4$ \\
			 \hline
		\end{tabular}
	\end{center}
	Có bao nhiêu học sinh lớp $11\text{A}$  có chiều cao từ $165$ cm đến dưới $175$ cm?
	 \choice
		{22} 
		{\True 16} 
		{6} 
		 {10}
	 \loigiai{
	Số  học sinh lớp $11\text{A}$ có chiều cao từ $165$ cm đến dưới $165$ cm là $10+6=16$ (em).}
\end{ex}

% nguồn: Ok_TracNghiemLop11_Lan1
\begin{ex}%[1D5N1-2]%[ma:57dbdf8e]
	Chiều cao  của $35$ cây bạch đàn (đơn vị: $m$)
	\begin{center}
		\begin{tabular}{|c|c|c|c|c|c|c|c|c|c|}
			\hline 
			$6{,}6$ & $7{,}5$ & $8{,}2$ & $8{,}2$ & $7{,}8$ & $7{,}9$ & $9{,}0$ &  $8{,}9$ & $8{,}2$ & $7{,}2$ \\
			\hline 
			$7{,}5$ & $8{,}3$ & $7{,}4$ & $8{,}7$ & $7{,}7$ &  $7{,}0$ & $9{,}4$ & $8{,}7$ & $8{,}0$ & $7{,}7$ \\
			\hline 
			$7{,}8$ & $8{,}3$  & $8{,}6$ & $8{,}1$ & $8{,}1$ & $8{,}5$ & $6{,}9$ & $8{,}0$ & $7{,}6$ & $7{,}9$ \\
			\hline 
			$7{,}3$  & $8{,}5$ & $8{,}4$ & $8{,}0$ & $8{,}8$ & & & & & \\
			 \hline
		\end{tabular}
	\end{center}
	Từ mẫu số liệu không ghép nhóm trên, ghép các số liệu thành $6$ nhóm theo các nửa khoảng có độ dài bằng nhau.  Nhóm chiếm tỉ lệ cao nhất là
	 \choice
		{$[7{,}0; 7{,}5)$} 
		 {$[7{,}5; 8{,}0)$}
		{\True  $[8{,}0; 8{,}5)$}
		 {$[8{,}5; 9{,}0)$}
	\loigiai{ 
		Khoảng biến thiên  là $9{,}4-6{,}6=2{,}8$\\
		Ta chia thành các nhóm sau: $[6{,}5; 7{,}0),[7{,}0; 7{,}5),[7{,}5; 8{,}0),[8{,}0; 8{,}5),[8{,}5; 9{,}0),[9{,}0;  9{,}5)$.\\
		Đếm số giá trị của mỗi nhóm, ta có bảng ghép nhóm sau: 
		\begin{center}
			\begin{tabular}{|c|c|}
				 \hline
				Chiều cao $(\mathrm{m})$ & Số  cây \\
				\hline 
				$[6{,}5; 7{,}0)$ &  $2$ \\
				 \hline
				$[7{,}0;  7{,}5)$ & $4$ \\
				 \hline
				$[7{,}5; 8{,}0)$  & $9$ \\
				\hline 
				$[8{,}0;  8{,}5)$ & $11$ \\
				\hline 
				$[8{,}5; 9{,}0)$ & $7$  \\
				\hline 
				$[9{,}0;  9{,}5)$ & $2$ \\
				 \hline
			\end{tabular}
		\end{center}
	Từ bảng số liệu ta thấy nhóm chiếm tỉ lệ cao nhất là $[8{,}0;  8{,}5)$.}
\end{ex}

% nguồn: Ok_TracNghiemLop11_Lan1
\begin{ex}%[1D5N1-2]%[ma:8578c20e]
	Cho  mẫu số liệu như bảng bên dưới.
	\begin{center}
		\begin{tabular}{|c|c|c|c|c|c|}
			 \hline
			Giá  trị & $[u_1;u_2)$ & $[u_2;u_3)$ & $[u_3;u_4)$ & $[u_4;u_5)$ & $[u_5;u_6)$ \\
			 \hline
			Tần số & 6 & 1 & 3 &  9 & 7 \\
			\hline 
		\end{tabular}
	\end{center}
	Nhóm $[u_1;u_2)$ có giá trị đại diện  là
	 \choice
		 {$u_1+u_2$}
		 {\True $\dfrac{1}{2}\cdot(u_1+u_2)$}
		{$\dfrac{1}{2}\cdot(u_1-u_2)$} 
		{$\dfrac{1}{2}\cdot(u_2-u_1)$} 
	 \loigiai{
		 Giá trị đại diện của nhóm $[u_1;u_2)=\dfrac{1}{2}\cdot(u_1+u_2)$.
	} 
\end{ex}

% nguồn: Ok_TracNghiemLop11_Lan1
\begin{ex}%[1D5N1-2]%[ma:7d43aa4d]
	Điều tra về chiều cao của học sinh khối lớp  11, ta được mẫu số liệu sau:
	\begin{center}
		\begin{tabular}{|c|c|}
			\hline Chiều cao $(\mathrm{cm})$ & Số học  sinh \\
			\hline$[150; 152)$ & $5$ \\ 
			 \hline$[152; 154)$ & $18$ \\
			\hline$[154;  156)$ & $40$ \\
			\hline$[156;  158)$ & $26$ \\
			\hline$[158; 160)$  & $8$ \\
			\hline$[160; 162)$ & $3$ \\ 
			\hline Tổng  & $N=100$ \\
			 \hline
		\end{tabular}
	\end{center}
	Mẫu số liệu  ghép nhóm đã cho có tất cả bao nhiêu nhóm?
	\choice 
		 {$5$}
		{\True $6$} 
		 {$7$}
		 {$12$}
	\loigiai{ 
		Mẫu số liệu ghép nhóm đã cho có  tất cả $6$ nhóm.
	 }
\end{ex}

% nguồn: Ok_TracNghiemLop11_Lan1
\begin{ex}%[1D5N1-2]%[ma:16be59f4]
	Cho bảng số liệu ghép nhóm  về độ dài của 60 lá dương xỉ trưởng thành như sau
	\begin{center}
		\begin{tabular}{|c|c|c|c|c|}
			\hline Độ dài (cm) & { $[10; 20)$ } & { $[20; 30)$ } & { $[30; 40)$ } & { $[40;  50)$ } \\
			\hline Số lá dương xỉ & 8 & 18 &  24 & 10 \\
			\hline 
		\end{tabular}
	\end{center}
	Số lá có  chiều dài từ $30$ cm đến dưới $40$ cm là
	 \choice
		{$18$} 
		{$34$} 
		{$10$} 
		{\True $24$} 
	\loigiai{ 
		Số lá có chiều  dài từ $30$ cm đến dưới $40$ cm là $24$.
	 }
\end{ex}

% nguồn: Ok_TracNghiemLop11_Lan1
\begin{ex}%[1D5N1-2]%[ma:58c27893]
	Cho độ dài của $60$ là dương xỉ trưởng thành được cho bằng mẫu số liệu  ghép nhóm sau:
	\begin{center}
		\begin{tabular}{|c|c|c|c|c|}
			\hline 
			Độ dài của lá (cm) & $[10;20)$  & $[20;30)$ & $[30;40)$ & $[40;50)$ \\
			\hline 
			Số lá &  $9$ & $17$ & $24$ & $10$ \\
			\hline 
		\end{tabular}
	\end{center}
	Tần  số của nhóm $\left[30;40\right)$
	 \choice
		{$9$} 
		 {$17$}
		 {\True $24$}
		{$10$} 
	 \loigiai{
		Tần số của nhóm $\left[30;40\right)$ là 24. 
	} 
\end{ex}

% nguồn: Ok_TracNghiemLop11_Lan1
\begin{ex}%[1D5N1-2]%[ma:83dfeda2]
	Cho bảng mẫu số liệu ghép nhóm là chiều cao của học sinh lớp  $5$ tuổi
	\begin{center}
		\begin{tabular}{|c|c|}
			\hline 
			Nhóm chiều  cao & Tần số \\
			 \hline
			$\left[85;90\right)$ & $1$ \\ 
			 \hline
			$\left[90;95\right)$ & $4$  \\
			\hline 
			$\left[95;100\right)$ &  $8$ \\
			\hline 
			$\left[100;105\right)$  & $12$ \\
			 \hline
			$\left[105;110\right)$  & $3$ \\
			 \hline
			$\left[110;115\right)$ & $2$ \\ 
			 \hline
			& $n=30$ \\ 
			\hline 
		\end{tabular}
	\end{center}
	Số liệu và nhóm của  bảng trên là
	\choice 
		{$30$  và $5$}
		{$115$ và $30$} 
		{$115$  và $6$}
		{\True $30$  và $6$}
	\loigiai{ 
	Mẫu số  liệu đó gồm $30$ số liệu và $6$ nhóm.}
\end{ex}

% nguồn: Ok_TracNghiemLop11_Lan1
\begin{ex}%[1D5N1-2]%[ma:284eee7a]
	\immini[thm]{Thống kê chiều cao của học sinh lớp $11 \mathrm{A}$ ta có bảng số liệu như hình bên. Hỏi lớp có bao nhiêu học sinh có chiều cao từ $168 \mathrm{~cm}$  trở lên?
	\choice 
		 {\True$11$}
		{$20$} 
		 {$31$}
		{$8$}}{\begin{tabular}{|c|c|}
		\hline 
		 Chiều cao (cm) & Số học sinh\\
		\hline 
		$[150; 156)$  & $8$\\
		$[156; 162)$ & $12$\\ 
		 $[162; 168)$ & $11$\\
		$[168;  174)$ & $8$\\
		 $[174; 180)$ & $3$\\
		\hline 
\end{tabular}}
\loigiai{ 
	Số học sinh có chiều cao từ $168 \mathrm{~cm}$ trở lên là $8+3=11$. 
} 
\end{ex}

% nguồn: Ok_TracNghiemLop11_Lan1
\begin{ex}%[1D5N1-2]%[ma:2c3cc58c]
	Khảo sát thời gian chạy bộ trong một ngày của một số học sinh khối $11$ thu được mẫu số liệu  ghép nhóm sau
	\begin{center}
		{\tabcolsep = 2mm 
		\begin{tabular}{|c|c|c|c|c|c|}
			 \hline
			Thời gian (phút) & $[0; 20)$ & $[20;  40)$ & $[40; 60)$ & $[60; 80)$ & $[80; 100)$ \\
			 \hline
			Số học sinh & $5$ & $9$ & $12$ &  $10$ & $6$\\
			 \hline
		\end{tabular}
		 }
	\end{center}
	Giá trị đại diện  của nhóm $[0; 20)$ là
	\choice 
		 {\True $10$}
		 {$20$}
		{$30$} 
		 {$40$}
	 \loigiai{
		Giá trị đại diện của nhóm $[0; 20)$  là $\dfrac{0+20}{2}=10$.
	} 
\end{ex}

% nguồn: Ok_TracNghiemLop11_Lan1
\begin{ex}%[1D5N1-2]%[ma:ea45279d]
	Đo chiều  cao (tính bằng cm) của $500$ học sinh trong một trường THPT, ta thu được kết quả như bảng sau
	\begin{center}
		\renewcommand{\arraystretch}{1.5} 
		\begin{tabular}{|c|c|c|c|c|c|}
			 \hline
			Chiều cao & {$[150; 154)$} & {$[154; 158)$} & {$[158;  162)$} & {$[162; 166)$} & {$[166; 170)$}\\
			 \hline
			Số học sinh & 25 & 50 & 200 & 175 & 50  \\
			\hline 
		\end{tabular}
	\end{center}
	Giá trị đại  diện của nhóm $[150;154)$ là
	 \choice
		{$153$} 
		 {$25$}
		{\True $152$} 
		{$164$} 
	\loigiai{ 
		Giá trị đại diện  của nhóm $[150;154)$ là $\dfrac{150+154}{2}=152$.
	} 
\end{ex}

% nguồn: Ok_TracNghiemLop11_Lan1
\begin{ex}%[1D5N1-2]%[ma:9a92bc24]
	 Cho bảng số liệu ghép nhóm về độ tuổi và số lượng khách hàng của một cửa hàng như sau
	\begin{center}\begin{tabular}{|c|c|c|c|c|c|c|}
		 \hline
		Khoảng tuổi & $[23; 32)$ & $[32; 41)$ & $[41; 50)$ &  $[50; 59)$ & $[59; 68)$\\
		\hline 
		Số người & $15$ & $12$ & $7$ &  $11$ & $12$ \\
		 \hline
	\end{tabular}
\end{center}
Tính giá trị đại diện của nhóm  $[59;68)$.
\choice 
	 { $68$}
	 { $31{,}75$}
	{ \True $63{,}5$} 
	 { $59$}
 \loigiai{
	Giá trị đại diện của nhóm $[59;68)$ là $\dfrac{59+68}{2}=63{,}5$. 
} 
\end{ex}

% nguồn: Ok_TracNghiemLop11_Lan1
\begin{ex}%[1D5N1-2]%[ma:c05d98eb]
	 Cho mẫu số liệu ghép nhóm về thống kê điểm số (thang điểm $20$) của $100$ học sinh tham dự kỳ thi học sinh giỏi toán, ta có bảng số liệu sau
	\begin{center}
		 \renewcommand{\arraystretch}{1.5}
		\begin{tabular}{|c|c|c|c|c|c|c|}
			 \hline
			Điểm & {$[8; 10)$} &  {$[10; 12)$} & {$[12; 14)$} & {$[14; 16)$} & {$[16; 18)$}& {$[18; 20)$}\\
			 \hline
			Số học sinh &  6 & 21 & 30 & 25 & 14&4 \\
			\hline 
		\end{tabular}
	\end{center}
	 Nhóm $[12; 14)$ có tần số là
	\choice 
		{\True  $30$}
		 {$6$}
		{$21$} 
		{$25$} 
	 \loigiai{
		Dựa vào bảng số liệu, ta có tần số của  nhóm $[12; 14)$ là $30$.
	} 
\end{ex}

% nguồn: Ok_TracNghiemLop11_Lan1
\begin{ex}%[1D5N1-2]%[ma:a8c36d27]
	Cho mẫu số liệu về chiều cao $(\mathrm{cm})$ của các học $\sinh$  nữ trong khối 11 của một trường như sau:
	\begin{center}
		\begin{tabular}{|l|c|c|c|c|c|c|}
			\hline Chiều cao & {$[145; 150)$}  & {$[150; 155)$} & {$[155; 160)$} & {$[160; 165)$} & {$[165; 170)$} & {$[170; 175)$} \\
			\hline Số học sinh & $20$ & $45$ & $34$  & $27$ & $15$ & $4$ \\
			\hline 
		\end{tabular}
	\end{center}
	Mẫu số  liệu trên có bao nhiêu số liệu, bao nhiêu nhóm?
	\choice 
		{$145$ số liệu; $6$ nhóm} 
		 {$30$ số liệu; $5$ nhóm}
		{\True $6$ số liệu; $145$ nhóm} 
		{$5$ số liệu;  $30$ nhóm}
	\loigiai{ 
		 Mẫu số liệu $(T)$ có
		$$20+45+34+27+15+4=145\text{  (số liệu).}$$
		$6$ nhóm là $[145; 150)$; $[150; 155)$; $[155; 160)$;  $[160; 165)$; $[165; 170)$; $[170; 175)$.
	 }
\end{ex}

% nguồn: Ok_TracNghiemLop11_Lan1
\begin{ex}%[1D5N1-2]%[ma:f7d36e2b]
	Điều tra về chiều cao của học sinh khối lớp $10$,  ta có kết quả sau
	\begin{center}
		\begin{tabular}{|c|c|c|}
			\hline Nhóm & Chiều cao $(\mathrm{cm})$ & Số học  sinh \\
			\hline 1 & {$[150;  152)$} & $5$ \\
			 \hline 2 & {$[152; 154)$} & $18$ \\
			\hline 3 & {$[154; 156)$}  & $40$ \\
			\hline 4 & {$[156; 158)$} & $26$ \\ 
			\hline 5 &  {$[158; 160)$} & $8$ \\
			\hline 6 &  {$[160; 162)$} & $3$ \\
			\hline \multicolumn{2}{|r|}{} &  $N=100$ \\
			\hline 
		\end{tabular}
	\end{center}
	 Giá trị đại diện của nhóm thứ tư là
	 \choice
		 {$156{,}5$}
		{\True $157$} 
		 {$157{,}5$}
		 {$158$}
	 \loigiai{
		Giá trị đại diện của nhóm thứ tư  là $\dfrac{156+158}{2}=157$.
	} 
\end{ex}

% nguồn: Ok_TracNghiemLop11_Lan1
\begin{ex}%[1D5N1-2]%[ma:e7d40969]
	Cho số liệu về chiều cao của $45$ học sinh nữ lớp 11 ở một trường THPT (đơn vị$\colon cm$) như sau 
	\begin{center}
		\begin{tabular}{|c|c|c|c|c|c|c|c|c|}
			 \hline
			$150$ & $160$ & $151$ & $156$ &  $154$ & $158$ & $155$ & $157$ & $157$ \\
			\hline 
			$155$ & $153$ & $159$ & $163$ & $161$ & $151$ & $163$ &  $164$ & $160$ \\
			\hline 
			$156$ & $161$ &  $158$ & $152$ & $153$ & $161$ & $157$ & $159$ & $154$ \\
			\hline 
			$158$ & $157$ & $155$ & $160$ & $164$  & $160$ & $160$ & $161$ & $161$ \\
			 \hline
			$162$ & $159$ & $162$ & $156$ & $159$ & $164$ & $154$ &  $150$ & $157$\\
			\hline 
		\end{tabular}
	\end{center}
	 Bạn $C$ quyết định lập bảng tần số ghép nhóm theo mẫu sau
	\begin{center}
		\begin{tabular}{|c|c|c|c|c|c|}
			\hline 
			Chiều cao & $[150;152]$ & $[153;155]$ &  $[156;158]$ & $[159;161]$ & $[162;164]$ \\
			 \hline
			Tần số  &&&&&\\
			\hline 
		\end{tabular}
	\end{center}
	Lúc này tần  số của nhóm $[159;161]$ là bao nhiêu?
	 \choice
		 {$12$}
		 {$13$}
		{\True $14$} 
		 {$15$}
	 \loigiai{
		Nhóm $[159;161]$  gồm các phần tử
		\begin{center}
			\begin{tabular}{|c|c|c|c|c|c|c|c|c|}
				\hline 
				 $150$ & $\textbf{160}$ & $151$ & $156$ & $154$ & $158$ & $155$ & $157$ & $157$ \\
				\hline 
				$155$ & $153$ & $\textbf{\textbf{159}}$ & $163$ & $\textbf{161}$ & $151$ & $163$ & $164$ & $\textbf{160}$ \\ 
				\hline 
				$156$ & $\textbf{161}$ & $158$ & $152$ & $153$ & $\textbf{161}$ & $157$ & $\textbf{159}$ &  $154$ \\
				\hline 
				$158$ & $157$ & $155$ & $\textbf{160}$ & $164$ & $\textbf{160}$ & $\textbf{160}$ & $\textbf{161}$ &  $\textbf{161}$ \\
				 \hline
				$162$ &  $\textbf{159}$ & $162$ & $156$ & $\textbf{159}$ & $164$ & $154$ & $150$ & $157$\\
				\hline 
			\end{tabular}
		\end{center}
		Do đó tần số của  nhóm $[159;161]$ là $14$.
	 }
\end{ex}

% nguồn: Ok_TracNghiemLop11_Lan1
\begin{ex}%[1D5N1-2]%[ma:545d5dac]
	Khảo sát thời gian tập thể dục mỗi ngày của học sinh lớp 11C thu được bảng số liệu  ghép nhóm sau:
	\begin{center}
		\begin{tabular}{|c|c|c|c|c|c|}
			 \hline
			\textbf{Thời gian (phút)} &  $ \left[0;20\right) $ & $ \left[20;40\right) $ & $ \left[40;60 \right) $ & $ \left[60;80\right) $&$ \left[80;100 \right) $ \\
			\hline 
			\textbf{Số học sinh} & $ 5 $ &  $ 9 $ & $ 12 $ & $ 10 $&$6$ \\
			 \hline
		\end{tabular}
	\end{center}
	Giá trị đại diện của nhóm $\left[20;40\right)$ là 
	 \choice
		 { $10$}
		{$20$} 
		 {\True $30$}
		 {$40$}
	 \loigiai{
		Giá trị  đại diện của nhóm $\left[20;40\right)$ là $c_2=\dfrac{20+40}{2}=30$.
	} 
\end{ex}

% nguồn: Ok_TracNghiemLop11_Lan1
\begin{ex}%[1D5N1-2]%[ma:94b3857d]
	Tìm hiểu thời gia xem tivi trong tuần trước (đơn vị: giờ) của một số học sinh thu được kết quả sau: 
	\begin{center}
		\begin{tabular}{|l|c|c|c|c|c|}
			\hline Thời  gian (giờ) & {$[0; 5)$} & {$[5; 10)$} & {$[10; 15)$} & {$[15; 20)$} & {$[20; 25)$} \\
			\hline Số học sinh & $8$  & $16$ & $4$ & $2$ & $2$ \\
			 \hline
		\end{tabular}
	\end{center}
	Giá trị đại diện của nhóm $[20; 25)$ là 
	 \choice
		 {\True $22{,}5$}
		{$23$} 
		 {$20$}
		{$5$} 
	 \loigiai{
		Giá trị đại diện của nhóm  $[20; 25)$ là $\dfrac{20+25}{2}=22{,}5$.
	 }
\end{ex}

% nguồn: Ok_TracNghiemLop11_Lan1
\begin{ex}%[1D5N1-2]%[ma:e0a43f70]
	Một trường trung học cơ sở chọn $36$ học sinh nam của khối $9$ để  đo chiều cao của các bạn học sinh đó và thu được mẫu số liệu theo bảng sau (đơn vị là centimet)
	\begin{center}
		\begin{tabular}{|l|l|l|l|l|l|l|l|l|l|}
			 \hline
			$160$ & $161$ & $161$ & $162$  & $162$ & $162$ & $163$ & $163$ & $163$ & $164$ \\
			 \hline
			$164$ & $164$ & $164$ & $165$  & $165$ & $165$ & $165$ & $165$ & $166$ & $166$ \\
			 \hline
			 $166$ & $166$ & $167$ & $167$ & $168$ & $168$ & $168$ & $168$ & $169$ & $169$ \\
			 \hline
			$170$ & $171$ & $171$ & $172$ & $172$ & $174$ & & & &  \\
			\hline 
		\end{tabular}
	\end{center}
	Từ mẫu số liệu không ghép nhóm trên, hãy ghép các số liệu thành $5$ nhóm theo các nửa khoảng có độ dài bằng nhau.  Các nhóm đó là
	 \choice
		 {$[160;163);[163;165);[165;170);[170;172);[172;175)$}
		{$[160; 165);[165; 168);[168; 170);[170;  172);[172; 175)$}
		{\True $[160; 163);[163; 166);[166; 169);[169; 172);[172; 175)$} 
		{$[160; 162);[162; 166);[166; 168);[168; 170);[172;  175)$}
	\loigiai{ 
		 Khoảng biến thiên là $174-160=14$.\\
		Để số liệu thành $5$ nhóm theo các nửa khoảng có  độ dài bằng nhau, ta chia thành các nhóm có độ dài là $3$.\\
		Ta sẽ chọn  đầu mút phải của nhóm cuối cùng là $175$.\\
		Ta được các nhóm $[160; 163);[163; 166);[166; 169);[169;  172);[172; 175)$.
	} 
\end{ex}

% nguồn: Ok_TracNghiemLop11_Lan1
\begin{ex}%[1D5N1-2]%[ma:dddae409]
	Tuổi thọ của $35$ bóng đèn (đơn vị: giờ) 
	\begin{center}
		\begin{tabular}{|l|l|l|l|l|l|l|}
			 \hline
			$1120$ &  $1150$ & $1121$ & $1170$ & $1136$ & $1150$ & $1140$ \\
			 \hline
			 $1130$ & $1165$ & $1142$ & $1133$ & $1157$ & $1115$ & $1132$ \\
			 \hline
			$1162$ & $1179$ & $1109$ & $1131$ & $1147$ & $1168$ & $1152$ \\ 
			 \hline
			$1134$ & $1116$ & $1177$ & $1145$ & $1164$ & $1111$  & $1125$ \\
			 \hline
			$1144$ & $1160$ & $1155$ & $1103$ & $1127$ & $1166$  & $1101$ \\
			 \hline
		\end{tabular}
	\end{center}
	Từ mẫu số liệu không ghép nhóm trên, ghép các số liệu thành  $4$ nhóm theo các nửa khoảng có độ dài bằng nhau. Độ dài của mỗi nhóm là
	\choice 
		{\True  $20$}
		 {$22$}
		 {$23$}
		 {$25$}
	\loigiai{ 
		Khoảng biến  thiên là $1179-1101=78$.\\
		Để số  liệu thành $4$ nhóm theo các nửa khoảng có độ dài bằng nhau, ta chia thành các nhóm có độ dài là $20$.\\
	Ta chia thành các nhóm  sau: $[1100; 1120),[1120; 1140),[1140; 1160),[1160; 1180)$.}
\end{ex}

% nguồn: Ok_TracNghiemLop11_Lan1
\begin{ex}%[1D5N1-2]%[ma:04d12c33]
	Một vườn  thú ghi lại tuổi thọ (đơn vị năm) của $20$ con hổ và thu được mẫu số liệu ghép nhóm được cho ở bảng sau
	\begin{center}
		\begin{tabular}{|c|c|c|c|c|c|}
			\hline 
			Nhóm & $\left[14;15\right)$ & $\left[15;16\right)$ & $\left[16;17\right)$ & $\left[17;18\right)$  & $\left[18;19\right)$ \\
			\hline 
			Tần số & 1 &  2 & 7 & 5 & 5 \\
			 \hline
		\end{tabular}
	\end{center}
	Giá trị đại diện  của nhóm $\left[16;17\right)$ là
	 \choice
		 {$16$}
		 {$17$}
		{\True  $16{,}5$}
		 {$1$}
	 \loigiai{
		Giá trị đại diện của nhóm $\left[16;17\right)$ là  $\dfrac{16+17}{2} = 16{,}5$.
	} 
\end{ex}

% nguồn: Ok_TracNghiemLop11_Lan1
\begin{ex}%[1D5N1-2]%[ma:5054783b]
	Công việc nào sau đây không phụ  thuộc vào các công việc của môn thống kê?
	\choice 
		{Thu  thập số liệu}
		{Trình  bày số liệu}
		 {Phân tích và xử lí số liệu}
		 {\True Ra quyết định dựa trên số liệu}
	 \loigiai{
		Công việc nào sau đây không phụ thuộc vào các công việc của môn thống kê là ra quyết định  dựa trên số liệu.
	} 
\end{ex}

% nguồn: Ok_TracNghiemLop11_Lan1
\begin{ex}%[1D5N1-2]%[ma:dd16a137]
	Khi khảo sát về cân nặng của $510$ học sinh nam khối $11$ trường  $A$, người ta thu được bảng tần số ghép nhóm như sau
	\begin{center}
		\begin{tabular}{|c|c|c|c|c|c|c|}
			\hline Nhóm & {$[45; 55)$} & {$[55; 65)$} & {$[65; 75)$} & {$[75; 85)$} & {$[85;  95)$} & {$[95; 105)$} \\
			\hline Tần số &  74 & 245 & 112 & 49 & 23 & 7 \\
			 \hline
		\end{tabular}
	\end{center}
	Tìm giá trị đại diện của nhóm $5$. 
	\choice 
		{\True  $90$}
		{$80$} 
		 {$60$}
		 {$70$}
	 \loigiai{
		Giá  trị đại diện của nhóm $5$ là $90$.
	 }
\end{ex}

% nguồn: Ok_Toan11_TracNghiem_Lan2
\begin{ex}%[1D5H1-2]%[ma:8b3da089]
	Kết quả khảo sát cân nặng của $25$ quả cam ở lô hàng A được cho  ở bảng sau:
	\begin{center}
		\begin{tabular}{|l|c|c|c|c|c|}
			\hline 
			Cân nặng (g) & {$[150; 155)$} &  {$[155; 160)$} & {$[160; 165)$} & {$[165; 170)$} & {$[170; 175)$} \\
			 \hline
			Số quả cam  ở lô hàng A & $1$ & $3$ & $7$ & $10$ & $4$ \\
			 \hline
		\end{tabular}
	\end{center}
	Nhóm chứa mốt là  nhóm nào?
	 \choice
		{$[150;  155)$}
		 {$[155; 160)$}
		 {\True $[165; 170)$}
		 {$[170; 175)$}
	 \loigiai{
		Nhóm chứa mốt của mẫu số  liệu ghép nhóm là nhóm có tần số lớn nhất.\\
		Nhóm chứa mốt của  mẫu số liệu ghép nhóm là nhóm $[165; 170)$.
	 }
\end{ex}

% nguồn: Ok_Toan11_TracNghiem_Lan2
\begin{ex}%[1D5H1-2]%[ma:902316f2]
	Một thư viện  thống kê số lượng sách được mượn mỗi ngày trong $3$ tháng ở bảng sau:
	\begin{center}
		\begin{tabular}{|c|c|}
			 \hline
			 Số sách & Số ngày\\
			 \hline
			$[16; 20]$ &  $3$\\
			\hline 
			$[21; 25]$  & $6$\\
			\hline 
			$[26; 30]$ & $15$\\ 
			\hline 
			$[31; 35]$ & $27$\\ 
			\hline 
			$[36; 40]$ &  $22$\\
			\hline 
			$[41; 45]$ &  $14$\\
			\hline 
			$[46; 50]$ & $5$\\ 
			 \hline
		\end{tabular}
	\end{center}
	Trong $3$  tháng trên số sách được mượn mỗi ngày cao nhất gần bằng giá trị nào sau đây?
	\choice 
		 {$32$}
		{\True $34$} 
		 {$33$}
		 {$35$}
	\loigiai{ 
		Do số sách là số nguyên nên ta hiệu chỉnh lại  như sau:
		\begin{center}
			\begin{tabular}{|c|c|}
				\hline 
				 Số sách & Số ngày\\
				\hline 
				$[15{,}5; 20{,}5)$  & $3$\\
				\hline 
				$[20{,}5; 25{,}5)$ & $6$\\ 
				\hline 
				$[25{,}5; 30{,}5)$ & $15$\\ 
				\hline 
				$[30{,}5;  35{,}5)$ & $27$\\
				\hline 
				 $[35{,}5; 40{,}5]$ & $22$\\
				 \hline
				$[40{,}5; 45{,}5]$ &  $14$\\
				 \hline
				 $[45{,}5; 50{,}5)$ & $5$\\
				 \hline
			\end{tabular}
		\end{center}
		Ta có nhóm chứa mốt của mẫu  số liệu trên là nhóm $[30,5; 35,5)$.\\
		Do đó  $u_m=30{,}5$, $n_{m-1}=15$, $n_m=27$, $n_{m+1}=22$, $u_{m+1}-u_m=35{,}5-30{,}5=5$.\\
		Mốt của mẫu số liệu ghép nhóm là $M_0=30{,}5+\dfrac{27-15}{(27-15)+(27-22)} \cdot 5=\dfrac{1157}{34} \approx 34$. 
	 }
\end{ex}

% nguồn: Ok_Toan11_TracNghiem_Lan2
\begin{ex}%[1D5H1-2]%[ma:a7851df9]
	Khảo sát thời gian tập  thể dục của một số học sinh khối $11$ thu được mẫu số liệu ghép nhóm sau
	\begin{center}
		\begin{tabular}{|c|c|c|c|c|c|}
			\hline 
			Thời gian & $[0;20)$ & $[20;40)$ & $[40;60)$  & $[60;80)$ & $[80;100)$ \\
			 \hline
			Số học sinh & $5$ & $9$ & $12$  & $10$ & $6$ \\
			 \hline
		\end{tabular}
	\end{center}
	Giá trị đại  diện của nhóm $[40;60)$ là
	 \choice
		{$0$} 
		{$42$} 
		{\True  $50$}
		{$100$} 
	 \loigiai{
		Giá trị  đại diện của nhóm $[40;60)$ là $\dfrac{40+60}{2}=50$.
	 }
\end{ex}

% nguồn: Ok_Toan11_TracNghiem_Lan2
\begin{ex}%[1D5H1-2]%[ma:724267b2]
	Khi thống kê chiều cao của học sinh khối $12$ trong một trường trung học, ta thu được  mẫu số liệu ghép nhóm sau
	\begin{center}
		\begin{tabular}{|c|c|c|c|c|c|c|}
			\hline 
			Chiều cao & $[150;156)$ & $[156;162)$ & $[162;168)$ & $[168;174)$ & $[174;180)$ & $[180;186)$  \\
			 \hline
			Số học  sinh & $5$ & $18$ & $40$ & $26$ & $8$ & $3$ \\
			\hline 
		\end{tabular}
	\end{center}
	Giá trị đại diện của nhóm  $[162;168)$ là
	\choice 
		{$5$} 
		{$40$} 
		 {$130$}
		{\True $165$} 
	\loigiai{ 
		 Giá trị đại diện của nhóm $[162;168)$ là $\dfrac{162+168}{2}=165$.
	} 
\end{ex}

% nguồn: Ok_Toan11_TracNghiem_Lan2
\begin{ex}%[1D5H1-2]%[ma:edadf4bf]
	Thời gian truy cập Internet mỗi buổi tối của một  số học sinh được cho trong bảng sau:
	\begin{center}
		\begin{tabular}{|l|c|c|c|c|c|}
			 \hline
			Thời gian (phút) & {$[9{,}5; 12{,}5)$} & {$[12{,}5; 15{,}5)$} & {$[15{,},5; 18{,}5)$}  & {$[18{,}5; 21{,}5)$} & {$[21{,}5; 24{,}5)$} \\
			 \hline
			Só học sinh & $3$ & $12$ & $15$  & $24$ & $2$ \\
			 \hline
		\end{tabular}
	\end{center}
	Tìm mốt của mẫu số liệu ghép  nhóm này.
	\choice 
		 {$18{,}37$}
		 {\True $19{,}37$}
		 {$19{,}73$}
		 {$18{,}73$}
	\loigiai{ 
		Tần số lớn  nhất là $24$ nên nhóm chứa mốt là nhóm $[18{,}5; 21{,}5]$.\\
		Ta có $ 
		u_m=18{,}5$, $u_{m+1}=21{,}5$, $n_m=24$, $n_{m+1}=15$;  $n_{m-1}=2$, $u_{m+1}-u_m=3$.\\
		Do đó: $M_0=18{,}5+\dfrac{24-15}{(24-15)+(24-2)}  \cdot 3=18{,}5+\dfrac{27}{31} \approx 19{,}37$
	} 
\end{ex}

% nguồn: Ok_Toan11_TracNghiem_Lan2
\begin{ex}%[1D5H1-2]%[ma:cb336ff7]
	Khảo sát thời gian tập  thể dục của một số học sinh khối $11$ thu được mẫu số liệu ghép nhóm sau:
	\begin{center}
		\begin{tabular}{|l|c|c|c|c|c|}
			\hline 
			Thời gian (phút) & {$[0; 20)$} & {$[20; 40)$} & {$[40; 60)$}  & {$[60; 80)$} & {$[80; 100)$} \\
			 \hline
			Số học  sinh & $5$ & $9$ & $12$ & $10$ & $6$ \\
			 \hline
		\end{tabular}
	\end{center}
	Mốt của mẫu số liệu  trên là
	\choice 
		 {$42$}
		 {\True $52$}
		{$53$} 
		 {$54$}
	\loigiai{ 
		Mốt chứa  trong nhóm $[40; 60)$.\\
		Do đó: $u_m=40; u_{m+1}=60\Rightarrow u_{m+1}-u_m=60-40=20, n_{m-1}=9;  n_m=12; n_{m+1}=10$.\\
		Và  $M_0=40+\dfrac{12-9}{(12-9)+(12-10)}(60-20)=52$.
	 }
\end{ex}

% nguồn: Ok_Toan11_TracNghiem_Lan2
\begin{ex}%[1D5H1-2]%[ma:77908c34]
	Mức thưởng tết cho  các nhân viên của một công ty được thống kê trong bảng sau
	\begin{center}
		\begin{tabular}{|c|c|c|c|c|c|}
			 \hline
			Mức  thưởng tết & $[5;10)$ & $[10;15)$ & $[15;20)$ & $[20;25)$ & $[25;30)$ \\
			 \hline
			Số nhân viên & $13$ & $35$ & $47$  & $25$ & $10$ \\
			\hline 
		\end{tabular}
	\end{center}
	 Giá trị đại diện của nhóm $[15;20)$ là
	\choice 
		{$5$} 
		 {\True $17{,}5$}
		 {$30$}
		{$130$} 
	\loigiai{ 
		 Giá trị đại diện của nhóm $[15;20)$ là $\dfrac{15+20}{2}=17{,}5$.
	 }
\end{ex}

% nguồn: Ok_Toan11_TracNghiem_Lan2
\begin{ex}%[1D5H1-2]%[ma:3d6e3d2c]
	Cho  mẫu số liệu ghép nhóm về chiều cao của $25$ cây dừa giống như sau:
	\begin{center}
		\begin{tabular}{|l|c|c|c|c|c|}
			\hline 
			Chiều cao (cm) &  {$[0; 10)$} & {$[10; 20)$} & {$[20; 30)$} & {$[30; 40)$} & {$[40; 50)$} \\
			\hline 
			 Số cây & $4$ & $6$ & $7$ & $5$ & $3$ \\
			\hline 
		\end{tabular}
	\end{center}
	Mốt  của mẫu số liệu ghép nhóm này là
	\choice 
		 {\True $M_0=\dfrac{70}{3}$}
		 {$M_o=\dfrac{50}{3}$}
		{$M_o=\dfrac{70}{2}$} 
		{$M_o=\dfrac{80}{3}$} 
	 \loigiai{
		Tấn số  lớn nhất là $7$ nên nhóm chưa mốt là $[20; 30)$.\\
		Ta có  $j=3$, $a_3=20$, $m_3=7$, $m_2=6$, $m_4=5$, $h=10$.\\
		Do đó 
		$M_0=20+\dfrac{7-6}{(7-6)+(7-5)} \cdot  10=\dfrac{70}{3}$.
	} 
\end{ex}

% nguồn: Ok_Toan11_TracNghiem_Lan2
\begin{ex}%[1D5H1-2]%[ma:f77f502e]
	Điểm kiểm tra giữa kỳ I của một lớp  được cô giáo chủ nhiệm ghi lại theo bảng sau
	\begin{center}
		\begin{tabular}{|l|c|c|c|c|c|}
			\hline 
			Điểm & {$[4{,}5; 5{,}5)$} & {$[5{,}5; 6{,}5)$} &  {$[6{,}5; 7{,}5)$} & {$[7{,}5; 8{,}5)$} & {$[8{,}5; 9{,}5)$} \\
			 \hline
			Số học sinh  & $3$ & $7$ & $9$ & $15$ & $6$ \\
			 \hline
		\end{tabular}
	\end{center}
	Mốt của số liệu ghép nhóm  trên là
	\choice 
		{$8{,}4$} 
		 {$8{,}9$}
		{\True  $7{,}9$}
		 {$11{,}6$}
	\loigiai{ 
		Mốt của số liệu ghép nhóm  là $7{,}5+\dfrac{15-9}{15-9+15-6} \cdot 1=7{,}9$.
	 }
\end{ex}

% nguồn: Ok_Toan11_TracNghiem_Lan2
\begin{ex}%[1D5H1-2]%[ma:da862f1c]
	Khảo sát về số giờ mượn sách thư viện của học sinh khối $11$ trường  Y ta được một mẫu số liệu ghép nhóm như sau:
	\begin{center}
		\begin{tabular}{|l|c|c|c|c|c|}
			\hline 
			Số giờ mượn &  {$[0; 4)$} & {$[4; 8)$} & {$[8; 12)$} & {$[12; 16)$} & {$[16; 20)$} \\
			 \hline
			Số học sinh &  $54$ & $78$ & $120$ & $45$ & $12$ \\
			\hline 
		\end{tabular}
	\end{center}
	Mốt của mẫu số liệu ghép nhóm trên gần nhất với giá trị nào sau  đây?
	\choice 
		{$120$} 
		 {$12$}
		{$8$} 
		{\True $9$} 
	\loigiai{ 
		Nhóm chứa  Mốt là $[8; 12)$.\\
		Do  đó $u_m=8; n_{m-1}=78; n_m=120; n_{m+1}=45; u_{m+1}-u_m=12-8=4.$.\\
		Mốt của mẫu số liệu ghép nhóm là $M_0=8+\dfrac{120-78}{(120-78)+(120-45)} \cdot 4=\dfrac{368}{39} \approx 9{,}4$. 
	} 
\end{ex}

% nguồn: Ok_Toan11_TracNghiem_Lan2
\begin{ex}%[1D5H1-2]%[ma:66cd8a60]
	Bảng số liệu  ghép nhóm sau cho biết chiều cao (cm) của $50$ học sinh lớp $11A$
	\begin{center}
		\begin{tabular}{|c|c|c|c|c|c|}
			 \hline
			Khoảng chiều cao (cm) & $[145;150)$ & $[150;155)$ & $[155;160)$ & $[160;165)$ &  $[165;170)$ \\
			 \hline
			Số học sinh & $7$ & $14$ & $10$ & $10$  & $9$ \\
			 \hline
		\end{tabular}
	\end{center}
	Tính mốt của mẫu số liệu ghép nhóm này (làm  tròn đến hàng phần trăm)
	 \choice
		 {\True $153{,}18$}
		 {$153{,}81$}
		 {$154{,}18$}
		 {$153{,}28$}
	\loigiai{ 
		Tần số  lớn nhất là 14 nên nhóm chứa mốt là nhóm $[150;155)$.
		 Ta có $j=2$; $a_2=150$; $m_2=14$; $m_1=7$; $m_3=10$; $h=5$.
		Do  đó $M_0=150+\dfrac{14-7}{(14-7)+(14-10)}\cdot 5\approx 153{,}18$.
	} 
\end{ex}

% nguồn: Ok_Toan11_TracNghiem_Lan2
\begin{ex}%[1D5H1-2]%[ma:651ad8e8]
	Khảo  sát thời gian xem ti vi trong một ngày của một số học sinh khối $11$ thu được mẫu số liệu ghép nhóm sau:
	\begin{center}
		\begin{tabular}{|c|c|c|c|c|c|c|}
			 \hline
			Thời gian  (phút) & $[0;20)$ & $[20;40)$ & $[40;60)$ & $[60;80)$ & $[80;100)$ \\
			 \hline
			Số học sinh & $5$ & $9$ & $12$ & $10$ & $6$ \\ 
			\hline 
		\end{tabular}
	\end{center}
	Nhóm chứa trung vị  là
	 \choice
		{$\left[0;20\right)$} 
		 {$\left[20;40\right)$}
		 {\True $\left[40;60\right)$}
		{$\left[60;80\right)$} 
	 \loigiai{
		Cỡ mẫu là: $n=5+9+12+10+6=42$. 
		Trung vị là $\dfrac{x_{21}+x_{22}}{2}$. Do hai giá trị $x_{21}$, $x_{22}$ thuộc nhóm $\left[40;60\right)$  nên nhóm chứa trung vị là $\left[40;60\right)$.
	} 
\end{ex}

% nguồn: Ok_Toan11_TracNghiem_Lan2
\begin{ex}%[1D5H1-2]%[ma:c909d14d]
	Người ta tiến hành phỏng vấn $40$ người về một mẫu áo khoác. Người điều tra yêu cầu cho điểm mẫu áo đó theo thang điểm là $100$. Kết quả được trình  bày trong bảng ghép nhóm sau:
	\begin{center}
		\begin{tabular}{|c|c|c|c|c|c|c|}
			 \hline
			Nhóm & $[50;60)$ & $[60;70)$  & $[70;80)$ & $[80;90)$ & $[90;100)$ & \\
			 \hline
			Tần số & $4$ & $5$ & $23$ & $6$ &  $2$ &$N = 40$ \\
			\hline 
		\end{tabular}
	\end{center}
	Mốt của mẫu số liệu  ghép nhóm trên (làm tròn kết quả đến hàng đơn vị) là
	 \choice
		 {$73$}
		{\True $74$} 
		{$75$} 
		{$76$} 
	\loigiai{ 
		\[M_0=70+\dfrac{18}{18+27}\cdot  10=74.\]
	} 
\end{ex}

% nguồn: Ok_Toan11_TracNghiem_Lan2
\begin{ex}%[1D5H1-2]%[ma:95751c36]
	 Nhiệt độ trung bình của $12$ tháng tại thành phố Hà Nội từ năm $1961$ đến hết năm $1990$ được cho trong bảng sau
	\begin{center}
		\begin{tabular}{|l|c|c|c|c|c|}
			 \hline
			Nhiệt độ $\left(^{\circ} C\right)$ & {$[12;  14)$} & {$[14; 16)$} & {$[16; 18)$} & {$[18; 20)$} & {$[20; 22)$} \\
			 \hline
			Tần số & $1$ & $3$  & $12$ & $9$ & $5$ \\
			\hline 
		\end{tabular}
	\end{center}
	 Mốt của mẫu số liệu trên là
	\choice 
		 {$15{,}7$}
		{$16{,}0$} 
		 {\True $17{,}5$}
		 {$17{,}0$}
	\loigiai{ 
		Nhóm chứa mốt của mẫu số liệu trên  là nhóm $[16; 18)$.\\
		Nên ta có $u_m=16, n_{m-1}=3, n_m=12, n_{m+1}=9,  u_{m+1}-u_m=18-16=2$.\\
		Mốt của mẫu số liệu  trên là $M_0=16+\dfrac{12-3}{(12-3)+(12-9)} \cdot 2=17{,}5$.
	 }
\end{ex}

% nguồn: Ok_Toan11_TracNghiem_Lan2
\begin{ex}%[1D5H1-2]%[ma:0fc4c49e]
	Khảo sát thời gian tự học bài ở nhà của một số em học sinh lớp $11$ thu được  mẫu ghép nhóm số lượng như sau:
	\begin{center}
		\begin{tabular}{|c|c|c|c|c|c|}
			\hline 
			Thời gian & {$[0; 30)$}  & {$[30; 60)$} & {$[60; 90)$} & {$[90; 120)$} & {$[120; 150)$} \\
			\hline 
			Số học  sinh & $9$ & $10$ & $9$ & $15$ & $7$ \\
			 \hline
		\end{tabular}
	\end{center}
	Mốt của số liệu ghép nhóm  trên là
	 \choice
		{\True  $102{,}86$}
		{$102{,}85$} 
		 {$102{,}8$}
		 {$102{,}9$}
	\loigiai{ 
		Mốt của số liệu ghép nhóm là $90+\dfrac{15-9}{15-9+15-7} \cdot  30 \approx 102{,}86$.
	 }
\end{ex}

% nguồn: Ok_Toan11_TracNghiem_Lan2
\begin{ex}%[1D5H1-2]%[ma:157cc1a9]
	Tìm hiểu thời gian làm bài tập về nhà (đơn vị: giờ) trong tuần trước của một số học sinh thu được kết quả như  sau
	\begin{center}
		\begin{tabular}{|c|c|c|c|}
			 \hline
			Thời gian  & $[18;20)$ & $[20;22)$ & $[22;24)$ \\
			\hline 
			Số học sinh & $16$ & $35$ & $8$ \\ 
			\hline 
		\end{tabular}
	\end{center}
	Tính thời gian làm bài tập về  nhà trung bình trong tuần trước của các bạn học sinh này.
	 \choice
		 {\True $20{,}73$ giờ}
		 {$21{,}73$ giờ}
		{$20{,}72$  giờ}
		{$21{,}72$ giờ} 
	\loigiai{ 
		Trong mỗi khoảng thời gian,  giá trị đại diện là trung bình cộng của giá trị hai đầu mút nên ta có bảng sau
		\begin{center}
			\begin{tabular}{|c|c|c|c|}
				 \hline
				 Thời gian & $19$ & $21$ & $23$ \\
				 \hline
				Số học sinh & $16$ & $35$ & $8$  \\
				\hline 
			\end{tabular}
		\end{center}
		Thời gian  làm bài tập về nhà trung bình là
		\[\overline{x}=\dfrac{19\cdot16 +  21\cdot35 + 23\cdot8}{59} \approx 20{,}73 \text{ giờ}.\]
	} 
\end{ex}

% nguồn: Ok_Toan11_TracNghiem_Lan2
\begin{ex}%[1D5H1-2]%[ma:fc47da4d]
	Cho đồ thị thể hiện điểm thi đánh giá năng lực của một trường đại học vào năm $2020$ dưới  đây.
	\begin{center}
		\begin{tikzpicture}[scale=0.9,font=\footnotesize,line join=round,line cap=round,>=stealth]
				\draw[-stealth] (0,0)  -- (17,0) node[above]{\textbf{Điểm}};
				\draw[-stealth] (0,0) -- (0,8) node[below  right]{\textbf{Số học sinh}};
				\foreach  \y in {0,20,...,140}
				\draw[dotted] (17,\y/20)  -- (0,\y/20) node[left] {$\y$};
				\foreach \x/\label in  {
					1/{[0; 350]},2/{[351;  400]},3/{[401; 450]},
					4/{[451; 500]},5/{[501; 550]},6/{[551; 600]},7/{[601; 650]},8/{[651; 700]},9/{[701; 750]},10/{[751; 800]},11/{[801; 850]},12/{[851; 900]},13/{[901; 950]},14/{[951; 1000]},15/{[1001; 1050]},16/{[1051; 1200]} 
				} 
				\draw (\x,0)--(\x,-0.2) 
					(\x-0.5,-0.3)  node[below,rotate=30] {\tiny $\label$};
				\foreach \x/\y/\h  in {
					1/0/0, 
					 2/0.05/1,
					 3/0.4/8,
					4/1.2/24, 
					5/2.7/54, 
					 6/4.75/95,
					 7/4.75/95,
					 8/6.65/133,
					 9/6.1/122,
					 10/5.2/104,
					 11/3.1/62,
					 12/2.75/55,
					 13/1.05/21,
					14/0.6/12, 
					 15/0.05/1,
					16/0/0 
				 }
					{ 
				 \fill[orange] (\x-0.4,0) rectangle (\x+0.4,\y);
				\draw[above] (\x,\y)  node{$\h$};
				 }
				\node[align=center,font=\bfseries] at (7.5,9) {Điểm thi đánh giá năng  lực của một trường đại học năm 2020};
			\end{tikzpicture}
	\end{center}
	Giá trị  đại diện cho nhóm chứa mốt của mẫu số liệu ghép nhóm trên là
	 \choice
		{$625{,}5$} 
		 {\True $675{,}5$}
		 {$725{,}5$}
		 {$775{,}5$}
	 \loigiai{
		Nhóm chứa mốt của mẫu số liệu ghép nhóm trên là nhóm $[651;700]$ và giá trị đại diện  của nhóm là $\dfrac{1}{2}(651+700)=675{,}5$.
	 }
\end{ex}

% nguồn: Ok_Toan11_TracNghiem_Lan2
\begin{ex}%[1D5H1-2]%[ma:23e2d9eb]
	Khảo sát thời gian tập thể dục của một số học sinh khối $11$ thu được mẫu số liệu ghép nhóm  sau
	\begin{center}
		\begin{tabular}{|c|c|c|c|c|c|}
			\hline 
			Thời gian (phút) & $[0;20)$ & $[20;40)$ &  $[40;60)$ & $[60;80)$ & $[80;100)$ \\
			\hline 
			Số học sinh & $5$ &  $9$ & $12$ & $10$ & $6$ \\
			 \hline
		\end{tabular}
	\end{center}
	Giá trị đại diện của nhóm $[20;40)$  là
	 \choice
		 {$10$}
		 {$20$}
		{\True  $30$}
		 {$40$}
	\loigiai{ 
		Giá trị đại  diện của nhóm $[20;40)$ là
		\[c=\dfrac{20+40}{2}=30.\] 
	} 
\end{ex}

% nguồn: Ok_Toan11_TracNghiem_Lan2
\begin{ex}%[1D5H1-2]%[ma:82c17a13]
	Cho mẫu số liệu ghép nhóm về thống kê nhiệt  độ tại một địa điểm trong $30$ ngày, ta có bảng số liệu sau:
	\begin{center}
		\begin{tabular}{|l|c|c|c|c|}
			\hline 
			Nhiệt độ $\left(^{\circ} C\right)$ & {$[18;  21)$} & {$[21; 24)$} & {$[24; 27)$} & {$[27; 30)$} \\
			\hline 
			Số ngày & $6$ &  $12$ & $9$ & $3$ \\
			\hline 
		\end{tabular}
	\end{center}
	Mốt  của mẫu số liệu ghép nhóm trên là
	\choice 
		 {$24$}
		{$25$} 
		 {$22$}
		{\True $23$} 
	\loigiai{ 
		Nhóm $2$ ứng với nửa  khoảng $\left[21; 24\right)$ là nhóm có tần số lớn nhất với $u=21; g=3; n_2=12$.\\
		Nhóm $1$ có tần  số $n_1=6$; Nhóm $3$ có tần số $n_3=9$.\\
		Ta có mốt  của mẫu số liệu là $M_0=21+\left(\dfrac{12-6}{2\cdot 12-6-9}\right) \cdot 3=23$.
	 }
\end{ex}

% nguồn: Ok_Toan11_TracNghiem_Lan2
\begin{ex}%[1D5H1-2]%[ma:3eaaa05a]
	Cho bảng tần số ghép nhóm số  liệu thống kê chiều cao của $40$ mẫu cây ở một vườn thực vật. Mốt của mẫu số liệu ghép nhóm này là bao nhiêu?
	\begin{center}
		\begin{tabular}{|c|c|}
			\hline 
			Nhóm & Tần số  \\
			\hline 
			$[30; 40)$  & $4$ \\
			\hline 
			$[40; 50)$ & $10$  \\
			\hline 
			$[50;  60)$ & $14$ \\
			 \hline
			$[60; 70)$  & $6$ \\
			 \hline
			$[70; 80)$ & $4$  \\
			\hline 
			$[80; 90)$ & $2$ \\ 
			\hline 
			& $n=40$ \\ 
			\hline 
		\end{tabular}
	\end{center}
	 \choice
		 {$56{,}6$}
		{\True  $53{,}3$}
		 {$55$}
		 {$50{,}6$}
	\loigiai{ 
		Tần số  lớn nhất là $14$ nên nhóm chứa mốt là nhóm $[50; 60)$.\\
		Ta có $j=3$,  $a_3=50$, $m_2=10$, $m_3=14$, $m_4=6$, $h=10$.\\
		Do đó $M_0=50+\dfrac{14-10}{(14-10)+(14-6)} \cdot 10 \approx 53{,}3$. 
	 }
\end{ex}

% nguồn: Ok_Toan11_TracNghiem_Lan2
\begin{ex}%[1D5H1-2]%[ma:25b09f18]
	Cho mẫu số liệu ghép nhóm về thống kê chiều cao của $35$ cây bạch đàn  trong rừng, ta có bảng số liệu sau:
	\begin{center}
		\begin{tabular}{|l|c|c|c|c|}
			\hline 
			Khoảng chiều cao $(m)$ & {$[6{,}5; 7{,}0)$} & {$[7{,}0; 7{,}5)$} &  {$[7{,}5; 8{,}0)$} & {$[8{,}0; 8{,}5)$} \\
			 \hline
			Số cây & $6$ & $15$ & $11$ & $3$ \\ 
			\hline 
		\end{tabular}
	\end{center}
	Mốt của mẫu số liệu ghép nhóm  trên là
	\choice 
		{\True $7{,}35$} 
		{$7{,}34$} 
		 {$3{,}75$}
		 {$7{,}43$}
	 \loigiai{
		Nhóm $2$ ứng với nửa khoảng $[7{,}0; 7{,}5)$ là nhóm có tần số lớn  nhất với $u=7{,}0; g=0{,}5$; $n_2=15$.\\
		Nhóm $1$ có tần số $n_1=6$; Nhóm $3$ có tần số  $n_3=11$.\\
		Ta có mốt  của mẫu số liệu là $M_0=7{,}0+\left(\dfrac{15-6}{2\cdot 15-6-11}\right) \cdot 0{,}5=7{,}35$.
	} 
\end{ex}

% nguồn: Ok_Toan11_TracNghiem_Lan2
\begin{ex}%[1D5H1-2]%[ma:7fa7dd4d]
	Cho mẫu số liệu ghép nhóm về thống kê thời gian hoàn thành một bài kiểm tra trực tuyến  của $100$ học sinh, ta có bảng số liệu sau:
	\begin{center}
		\begin{tabular}{|l|c|c|c|c|c|c|}
			 \hline
			Thời gian (phút) & {$[33; 35)$} & {$[35; 37)$} & {$[37; 39)$} & {$[39; 41)$}  & {$[41; 43)$} & {$[43; 45)$} \\
			 \hline
			Số học sinh & $4$ & $13$ & $38$  & $27$ & $14$ & $4$ \\
			\hline 
		\end{tabular}
	\end{center}
	Tính  mốt của mẫu số liệu trên.
	 \choice
		 {$38{,}93$}
		 {\True $38{,}39$}
		 {$39{,}28$}
		 {$39{,}82$}
	\loigiai{ 
		Nhóm $3$ ứng với nửa khoảng $[37; 39)$ là nhóm có tần số lớn nhất với $u=37; g=2;  n_3=38$.\\
		Nhóm $2$ có tần số $n_2=13$; Nhóm $4$ có tần số  $n_4=27$.\\
		Ta có mốt của mẫu số liệu là  $M_0=37+\left(\dfrac{38-13}{2\cdot 38-13-27}\right) \cdot 2=38{,}39$.
	} 
\end{ex}

% nguồn: Ok_Toan11_TracNghiem_Lan2
\begin{ex}%[1D5H1-2]%[ma:5e2375b4]
	Doanh thu bán hàng trong $20$ ngày được lựa chọn ngẫu nhiên của một của  hàng được ghi lại ở bảng sau:
	\begin{center}
		\begin{tabular}{|l|c|c|c|c|c|}
			 \hline
			Doanh thu & {$[5; 7)$} & {$[7; 9)$} & {$[9; 11)$} & {$[11; 13)$} & {$[13;  15)$} \\
			\hline 
			Số ngày &  $2$ & $7$ & $7$ & $3$ & $1$ \\
			 \hline
		\end{tabular}
	\end{center}
	Mốt của mẫu số liệu trên  thuộc khoảng nào trong các khoảng dưới đây?
	\choice 
		{$[7;9)$} 
		{\True  $[9;11)$}
		{$[11;13)$} 
		 {$[13;15)$}
	\loigiai{ 
		Có $2$ nhóm chứa mốt của mẫu số liệu trên  đó là $[7; 9)$ và $[9; 11)$, do đó:\\
		Xét nhóm $[7;  9)$ ta có $M_0=7+\dfrac{7-2}{(7-2)+(7-7)}\cdot (9-7)=9$.\\
		Xét nhóm $[9; 11)$  ta có $M_0'=9+\dfrac{7-7}{(7-7)+(7-3)}\cdot (11-9)=9$.\\
		Vậy mốt của mẫu số liệu là $9$. 
	} 
\end{ex}

% nguồn: Ok_Toan11_TracNghiem_Lan2
\begin{ex}%[1D5H1-2]%[ma:8be0fc4a]
	Kiểm tra cân nặng $25$ quả dưa của hai nhà vườn A và B cho kết  quả như bảng sau:
	\begin{center}
		\begin{tabular}{|c|c|c|c|c|c|}
			 \hline
			\begin{tabular}{c} Cân nặng \\
				 (kg)
			\end{tabular} & {$[1{,},5; 2{,}5)$} & {$[2{,}5; 3{,}5)$} & {$[3{,}5; 4{,}5)$} & {$[4{,}5; 5{,}5)$} & {$[5{,}5; 6{,}5)$} \\
			\hline 
			\begin{tabular}{c} Số quả của \\
				vườn  A
			\end{tabular} & $2$ & $6$ & $4$ & $10$ & $3$ \\
			 \hline
			\begin{tabular}{c} Số quả của \\
				vườn  $B$
			\end{tabular} & $3$ & $7$ & $8$ & $5$ & $2$ \\
			 \hline
		\end{tabular}
	\end{center}
	Nếu  lấy một quả dưa bất kỳ của vườn A thì khả năng cao nhất là lấy được quả dưa nặng bao nhiêu kg?
	 \choice
		{$3{,}84$ kg} 
		{$4{,}84$ kg} 
		 {$4{,}9$ kg}
		{\True  $4{,}96$ kg}
	\loigiai{ 
		Mốt của cân nặng của các quả dưa vườn A là $ 
		 M_0=4{,}5+\dfrac{10-4}{10-4+10-3} \cdot 1=4{,}96.$\\
		Dựa vào kết quả trên, khả năng lấy được quả dưa nặng $4{,}96$ kg là cao nhất. 
	 }
\end{ex}

% nguồn: Ok_Toan11_TracNghiem_Lan2
\begin{ex}%[1D5H1-2]%[ma:e8589ce2]
	Khi thống kê điểm kiểm tra môn Toán của lớp $10\text{A}$, ta thu được mẫu số liệu ghép nhóm được cho  trong bảng sau
	\begin{center}
		\begin{tabular}{|c|c|}
			 \hline
			Nhóm & Tần số  \\
			 \hline
			$[2;4)$ & $4$ \\ 
			 \hline
			 $[4;6)$ & $15$ \\
			\hline 
			$[6;8)$ &  $12$ \\
			 \hline
			$[8;10]$  & $14$ \\
			 \hline
			\multicolumn{1}{|c|}{}& $n=45$  \\
			\hline 
		\end{tabular}
	\end{center}
	Độ dài mỗi nhóm của mẫu số liệu  bằng
	 \choice
		{\True  2}
		{4} 
		{1} 
		 {3}
	\loigiai{Xét nhóm $[2;4)$ thì độ dài của  nhóm là $L=4-2=2$.}
\end{ex}

% nguồn: Ok_Toan11_TracNghiem_Lan2
\begin{ex}%[1D5H1-2]%[ma:8cedc003]
	Điểm kiểm tra $15$ phút của $36$ học sinh lớp  $11$A được cho bởi bảng tần số ghép nhóm sau:
	\begin{center}
		\begin{tabular}{|c|c|}
			\hline 
			Nhóm điểm  & Tần số \\
			 \hline
			$[1; 3)$ & $3$  \\
			\hline 
			 $[3; 5)$ & $2$ \\
			\hline 
			$[5;  7)$ & $10$ \\
			\hline 
			$[7; 9)$ &  $14$ \\
			 \hline
			$[9; 11)$ &  $7$ \\
			 \hline
			& $n=36$  \\
			\hline 
		\end{tabular}
	\end{center}
	Mốt của  bảng ghép lớp trên là giá trị nào sau?
	 \choice
		{\True  $7{,}73$}
		{$6{,}12$} 
		{$5{,}09$} 
		 {$7{,}03$}
	 \loigiai{
		Áp dụng công thức  tính mốt $M_0=u+\left(\dfrac{n_i-n_{i-1}}{2n_i-n_{i-1}-n_{i+1}}\right) \cdot g$
		Với $i=4, u=7, g=2,  n_4=14, n_3=10, n_5=7$. Ta có
		$M_0=7+\left(\dfrac{14-10}{2\cdot  14-10-7}\right) \cdot 2=7{,}73$.
	 }
\end{ex}

% nguồn: Ok_Toan11_TracNghiem_Lan2
\begin{ex}%[1D5H1-2]%[ma:b4002ede]
	Cho mẫu số liệu  ghép nhóm về thống kê thời gian chạy $50$ m của $20$ học sinh, ta có bảng số liệu sau:
	\begin{center}
		\begin{tabular}{|l|c|c|c|c|c|}
			 \hline
			Thời gian (giây) & {$[7{,}8;  8{,}0)$} & {$[8{,}0; 8{,}2)$} & {$[8{,}2; 8{,}4)$} & {$[8{,}4; 8{,}6)$} & {$[8{,}6; 8{,}8)$} \\
			 \hline
			Số học sinh & $9$ & $3$ & $5$ & $2$ & $1$  \\
			\hline 
		\end{tabular}
	\end{center}
	Mốt của mẫu số liệu ghép nhóm  trên là
	 \choice
		{$9{,}72$} 
		 {$9{,}27$}
		{\True  $7{,}92$}
		 {$7{,}29$}
	 \loigiai{
		Nhóm $1$ ứng với nửa  khoảng $[7{,}8; 8{,}0)$ là nhóm có tần số lớn nhất với $u=7{,}8; g=0,2; n_1=9$.\\
		Nhóm  $0$ có tần số $n_0=0$; Nhóm $2$ có tần số $n_2=3$.\\
		Ta có mốt của mẫu số liệu là $M_0=7{,}8+\left(\dfrac{9-0}{2\cdot 9-0-3}\right) \cdot  0{,}2=7{,}92$.
	} 
\end{ex}

% nguồn: Ok_Toan11_TracNghiem_Lan2
\begin{ex}%[1D5H1-2]%[ma:8127c46f]
	Cho các số liệu thống kê ghi trong bảng sau là thành tích chạy $50$  m của học sinh lớp $11\text{A}$ ở trường THPT (đơn vị: giây)
	\begin{center}
		\begin{tabular}{|c|c|c|c|c|c|c|c|c|c|c|c|}
			 \hline
			$6{,}3$ & $6{,}2$ & $6{,}5$ & $6{,}8$ & $6{,}9$ & $8{,}2$ &  $8{,}6$ & $7{,}1$ & $7{,}3$ & $7{,}5$ & $7{,}5$ \\
			\hline 
			$6{,}6$ & $6{,}7$ & $7{,}0$ & $7{,}1$ & $7{,}2$ & $8{,}3$  & $8{,}5$ & $7{,}6$ & $8{,}7$ & $7{,}6$ & $7{,}7$\\
			\hline 
			$7{,}4$ & $7{,}3$ & $7{,}2$ & $7{,}1$ & $7{,}0$ & $8{,}4$ & $8{,}1$ & $7{,}8$ &  $7{,}5$ & $7{,}7$ & $7,8$ \\
			 \hline
		\end{tabular}
	\end{center}
	 Trong lớp $11\text{A}$, số học sinh chạy $50$ m hết từ $7$ giây đến dưới $8{,}5$ giây chiếm bao nhiêu phần trăm?
	\choice 
		 {$20{,}25\%$}
		 {$40{,}46\%$}
		{\True $69{,}69\%$} 
		 {$80{,}89\%$}
	 \loigiai{
		Tần số của các  lớp
		\begin{center}
			\begin{tabular}{c c c}
				$n_1=2;$ & $n_2=5;$  & $n_3=10;$\\
				$n_4=9;$  & $n_5=4;$ & $n_6=3$.
			\end{tabular}
		\end{center}
		Tần suất của  các lớp
		\begin{center}
			\begin{tabular}{c c c}
				$f_1 \approx 6{,}06\%$; & $f_2 \approx15{,}15\%$; & $f_3 \approx 30{,}30\%$; \\ 
				$f_4 \approx 27{,}27\%$; & $f_5 \approx12{,}12\%$; &  $f_6 \approx 9{,}10\%$.
			\end{tabular}
		\end{center}
		Từ đó ta có bảng phân bố tần số ghép nhóm thành tích chạy $50$ m của học sinh lớp $11\text{A}$ ở trường  THPT
		\begin{center}
			\begin{tabular}{|c|c|}
				 \hline
				Lớp  thời gian chạy (giây) & Tần suất (\%) \\
				 \hline
				$[6{,}0;6{,}5)$ & $6{,}06$  \\
				\hline 
				$[6{,}5;7{,}0)$ & $15{,}15$ \\ 
				\hline 
				$[7{,}0;7{,}5)$ &  $30{,}30$ \\
				\hline 
				$[7{,}5;8{,}0)$ & $27{,}27$  \\
				 \hline
				$[8{,}0;8{,}5)$ & $12{,}12$  \\
				 \hline
				 $[8{,}5;9{,}0)$ & $9{,}10$ \\
				\hline 
				Cộng & $100(\%)$  \\
				 \hline
			\end{tabular}
		\end{center}
		Số học sinh chạy $50$ m  hết từ $7$ giây đến dưới $8{,}5$ giây chiếm số phần trăm là
	$30{,}30\%+27{,}27\%+12{,}12\%=69{,}69\%$.} 
\end{ex}

% nguồn: Ok_Toan11_TracNghiem_Lan2
\begin{ex}%[1D5H1-2]%[ma:a9248439]
	Điểm kiểm tra giữa kỳ I của 1 lớp được cô giáo chủ nhiệm ghi  lại theo bảng sau
	\begin{center}
		\begin{tabular}{|c|c|c|c|c|c|}
			 \hline
			Điểm & $[4{,}5;5{,}5)$ & $[5{,}5;6{,}5)$ & $[6{,}5;7{,}5)$ & $[7{,}5;8{,}5)$ &  $[8{,}5;9{,}5)$ \\
			 \hline
			Số học sinh &  $3$ & $7$ & $9$ & $15$ & $6$ \\
			 \hline
		\end{tabular}
	\end{center}
	Giá trị đại diện của nhóm $[7{,}5;8{,}5)$ là 
	 \choice
		 {$7$}
		 {$7{,}5$}
		 {\True $8$}
		{$8{,}5$} 
	\loigiai{ 
		 Giá trị đại diện của nhóm $[7{,}5;8{,}5)$ là
		\[\dfrac{7{,}5+8{,}5}{2}=8.\] 
	} 
\end{ex}

% nguồn: Ok_Toan11_TracNghiem_Lan2
\begin{ex}%[1D5H1-2]%[ma:0cda8392]
	Thống  kê về thời lượng mỗi trận đấu bi-a trong vòng tứ kết giải đấu European Open người ta được mẫu số liệu ghép nhóm như sau:
	\begin{center}
		\begin{tabular}{|l|c|c|c|c|c|}
			 \hline
			Thời gian & {$[9{,}5; 12{,}5)$} &  {$[12{,}5; 15{,}5)$} & {$[15{,}5; 18{,}5)$} & {$[18{,}5; 21{,}5)$} & {$[21{,}5; 24{,}5)$} \\
			 \hline
			Số  trận & $3$ & $12$ & $15$ & $24$ & $2$ \\
			\hline 
		\end{tabular}
	\end{center}
	Mốt của  mẫu số liệu trên gần nhất với giá trị nào sau đây?
	\choice 
		{$18$} 
		 {$18{,}5$}
		 {$19$}
		{\True $19{,}5$} 
	 \loigiai{
		Ta có nhóm chứa Mốt là $[18{,}5;  21{,}5)$.\\
		Do đó $u_m=18{,}5$; $n_{m-1}=15$;  $n_m=24$; $n_{m+1}=2$; $u_{m+1}-u_m=3$.\\
		Mốt của mẫu số liệu ghép nhóm là $M_0=18{,}5+\dfrac{24-15}{(24-15)+(24-2)} \cdot  3=\dfrac{1201}{62} \approx 19{,}4$.
	} 
\end{ex}

% nguồn: Ok_TracNghiemLop11_Lan1
\begin{ex}%[1D5H1-2]%[ma:86eb2e09]
	Cho bảng phân bố tần số, tần suất ghép lớp  chiều cao của các học sinh trong một lớp học như sau
	\begin{center}
		\begin{tabular}{|c|c|c|}
			\hline Lớp & Tần  số & Tần suất \\
			\hline$[150; 155)$ & 8 & $20 \%$  \\
			\hline$[155; 160)$ & 5 &  $12{,}5 \%$ \\
			\hline$[160; 165)$ & $a$ &  $30 \%$ \\
			\hline$[165; 170)$ & $b$ & $c \%$ \\ 
			\hline$[170; 175]$ & 6 &  $15 \%$ \\
			\hline Tổng & $n$ &  $100 \%$ \\
			\hline 
		\end{tabular}
	\end{center}
	Tìm giá  trị của $a+2b+10c$.
	\choice 
		 {$225$}
		 {$158$}
		{\True $255$} 
		 {$202$}
	 \loigiai{
		Ta có  $c=22,5$,
		vì lớp $\left[155;159\right)$ có tần số là $5$ và  tần suất là $12,5\%$ tức là chiếm $\dfrac{1}{8}$ mẫu số liệu nên $n=5\cdot8=40$.\\
		Suy  ra $\dfrac{a\cdot 100}{40}=30\% \Rightarrow a=12$.\\
		Đồng thời, $\dfrac{b\cdot 100}{40}=22{,}5\%  \Rightarrow b=9$.\\
		Vậy  $a+2b+10c=255$.
	} 
\end{ex}

% nguồn: Ok_TracNghiemLop11_Lan1
\begin{ex}%[1D5H1-2]%[ma:bc2a7139]
	Khảo sát thời gian tập thể dục trong ngày của một số học sinh khối  11 thu được mẫu số liệu ghép nhóm sau
	\begin{center}
		\begin{tabular}{|c|c|c|c|c|c|}
			\hline 
			Thời gian  (phút)&$[0;20)$&$[20;40)$&$[40;60)$&$[60;80)$&$[80;100)$\\
			 \hline
			Số học sinh  &$5$&$9$&$12$&$10$&$6$\\
			\hline 
		\end{tabular}
	\end{center}
	Mẫu số liệu ghép nhóm này  có số mốt là
	 \choice
		 {$50$}
		 {\True $52$}
		{$29$} 
		{$35$} 
	\loigiai{ 
		Vì mẫu số liệu  chỉ có một lớp có tần số lớn nhất, nên mẫu này có $1$ mốt.\\
		Tần số lớn nhất  là $12$ nên nhóm chứa mốt là nhóm thứ ba $[40; 60)$.\\
		Ta có  $j=3$; $a_3=40$; $m_3=12$; $m_2=9$; $m_4=10$; $ h=20$.\\
		 Khi đó $M_O=a_j+\dfrac{m_j-m_{j-1}}{\left(m_j-m_{j-1}\right)+\left(m_j-m_{j+1}\right)}\cdot h=40+\dfrac{12-9}{(12-9)+(12-10)}\cdot 20=52$.
	} 
\end{ex}

% nguồn: Ok_TracNghiemLop11_Lan1
\begin{ex}%[1D5H1-2]%[ma:1be41720]
	Cho bảng số liệu bên dưới. Số trung bình cộng của mẫu số liệu trên bằng bao nhiêu (làm tròn kết quả đến hàng phần trăm)?\begin{center}\begin{tabular}{|c|c|c|c|c|c|c|}
	 \hline
	 Nhóm & $[57; 58)$ & $[58; 59)$ & $[59; 60)$ & $[60; 61)$ & $[61; 62)$ & Tổng tần số \\
	\hline 
	Tần số & $6$ & $8$ & $8$ & $9$  & $9$ & $40$ \\
	 \hline
\end{tabular}
\end{center}
\choice 
	{$69$} 
	{$66$} 
	{$-36$} 
	 {\True $59{,}67$}
\loigiai{ 
	Ta  có các giá trị đại diện của các lớp là
	$$x_1 = 57{,}5; x_2 =  58{,}5; x_3 = 59{,}5; x_4 = 60{,}5; x_5 = 61{,}5.$$
	Áp dụng công thức tính số trung  bình cộng ta có
	$$\overline{x} = \dfrac{x_1n_1+x_2n_2+x_3n_3+x_4n_4+x_5n_5}{40} =  59{,}67.$$
} 
\end{ex}

% nguồn: Ok_TracNghiemLop11_Lan1
\begin{ex}%[1D5H1-2]%[ma:d84c7cb9]
	Khảo sát thời gian tập thể dục trong ngày của một số học sinh  khối $11$ thu được mẫu số liệu ghép nhóm sau
	\begin{center}
		\begin{tabular}{ |c|c|c|c|c|c| }
			\hline 
			Thời  gian (phút)& $[0;20)$ & $[20;40)$ & $[40;60)$ & $[60;80)$ & $[80;100)$ \\
			\hline 
			Số học sinh  & $5$ & $9$ & $12$ & $10$ & $6$ \\
			\hline 
		\end{tabular}
	\end{center}
	Giá  trị đại diện của nhóm $[20;40)$ là
	\choice 
		{$10$} 
		{\True  $30$}
		{$20$} 
		 {$40$}
	 \loigiai{
		Giá trị đại diện của nhóm $[20;40)$ là $\dfrac{20+40}{2}=30$. 
	} 
\end{ex}

% nguồn: Ok_TracNghiemLop11_Lan1
\begin{ex}%[1D5H1-2]%[ma:07c32f5d]
	Một người thống kê thời gian thực hiện các cuộc gọi điện thoại của mình  (đơn vị: phút) trong một tuần ở bảng sau:
	\begin{center}
		\begin{tabular}{|l|c|c|c|c|c|}
			\hline 
			Thời gian & $\left[0;60\right)$ & $\left[60;120\right)$  & $\left[120;180\right)$ &
			$\left[180;240\right)$ & $\left[240;300\right)$  \\
			\hline 
			Số  cuộc gọi & $7$ & $12$ & $6$ & $5$ & $2$\\
			 \hline
		\end{tabular}
	\end{center}
	Giá  trị đại diện của nhóm $\left[180;240\right)$ là
	 \choice
		 {$205$}
		{\True $210$} 
		{$200$} 
		{$220$} 
	\loigiai{ 
		Giá trị đại diện của nhóm  $\left[180;240\right)$ là $\dfrac{180+240}{2}=210$.
	 }
\end{ex}

% nguồn: Ok_TracNghiemLop11_Lan1
\begin{ex}%[1D5H1-2]%[ma:1555e059]
	Cho mẫy số liệu ghép nhóm về tuổi thọ (đơn vị tính  là năm) của một loại bóng đèn mới như sau.
	\begin{center}
		\begin{longtable}{|c|c|c|c|c|}
			\hline 
			Tuổi thọ (năm) & $[2;3{,}5)$ & $[3{,}5;5)$ & $[5;6{,}5)$  & $[6{,}5;8)$\\
			 \hline
			Số bóng đèn & $8$ & $22$ & $35$ & $15$\\ 
			\hline 
		\end{longtable}
	\end{center}
	Số mốt của mẫu số liệu  ghép nhóm này là
	 \choice
		{$0$} 
		{\True $1$} 
		 {$2$}
		{$3$} 
	 \loigiai{
		Tần số lớn  nhất là $35$ nên nhóm chứa mốt là nhóm $[3{,}5;5)$. Ta có $a_3=3{,}5$, $m_3=35$, $m_2=22$, $m_4=15$, $h=1{,}5$.\\
		Do đó mốt của mẫu số liệu ghép nhóm này là 
		 $$M_0=3{,}5+\dfrac{22-35}{(22-35)+(22-15)}\cdot 1{,}5=6{,}75.$$
		Vậy mẫu số  liệu ghép nhóm này có $1$ mốt.
	 }
\end{ex}

% nguồn: Ok_TracNghiemLop11_Lan1
\begin{ex}%[1D5H1-2]%[ma:2b7b9b23]
	Quãng đường (km) từ nhà đến nơi làm việc của $40$ công nhân một nhà  máy được ghi lại như sau:
	\begin{center}
		\begin{tabular}{|c|c|c|c|c|c|c|c|c|c|}
			\hline 
			$5$ &  $3$ & $10$ & $20$ & $25$ & $11$ & $13$ & $7$ & $12$ & $31$ \\
			 \hline
			$19$ & $10$ & $12$ & $17$ & $18$ & $11$  & $32$ & $17$ & $16$ & $2$ \\
			\hline 
			 $7$ & $9$ & $7$ & $8$ & $3$ & $5$ & $12$ & $15$ & $18$ & $3$ \\
			\hline 
			$12$  & $14$ & $2$ & $9$ & $6$ & $15$ & $15$ & $7$ & $6$ & $12$ \\
			\hline 
		\end{tabular}
	\end{center}
	Ghép nhóm dãy số liệu trên thành các khoảng có độ  rộng bằng nhau, khoảng đầu tiên là $[0;5)$. Tần suất nhóm $[10; 15)$ là
	\choice 
		{$15 \%$} 
		{$22{,}5 \%$} 
		{$25  \%$}
		 {\True $27{,}5 \%$}
	\loigiai{ 
		Ta chia thành các nhóm có độ dài là $5$. Ta sẽ chọn đầu mút phải của  nhóm cuối cùng là $35$. Ta có bảng ghép nhóm sau:
		\begin{center}
			\begin{tabular}{|l|c|c|c|c|c|c|c|}
				 \hline
				Quảng đường $(\mathrm{km})$ & $[0; 5)$ &  $[5; 10)$ & $[10; 15)$ & $[15; 20)$ & $[20; 25)$ & $[25; 30)$ & $[30; 35)$ \\
				 \hline
				Số công nhân & 6 & 10 & 11 & 9 & 1  & 1 & 2 \\
				\hline 
			\end{tabular}
		\end{center}
	Ta có tần  suất của lớp $[10; 15)$ là $\dfrac{11 \times 100}{40}=27{,}5 \%$. }
\end{ex}

% nguồn: Ok_TracNghiemLop11_Lan1
\begin{ex}%[1D5H1-2]%[ma:faf6f3f7]
	Cho bảng mẫu số liệu ghép nhóm là chiều cao  của học sinh lớp $5$ tuổi
	\begin{center}
		\begin{tabular}{|c|c|}
			\hline 
			Nhóm chiều cao & Tần số \\ 
			\hline 
			$\left[85;90\right)$ & $1$ \\ 
			\hline 
			$\left[90;95\right)$ &  $x-4$ \\
			 \hline
			$\left[95;100\right)$ & $x^2-3x+11$ \\ 
			\hline 
			$\left[100;105\right)$ & $12$ \\ 
			\hline 
			$\left[105;110\right)$ &  $3$ \\
			\hline 
			 $\left[110;115\right)$ & $2$ \\
			 \hline
		\end{tabular}
	\end{center}
	Tìm $x$ biết tần số  tích lũy của nhóm $3$ là $23$
	\choice 
		 {\True $5$}
		 {$3$}
		 {$8$}
		{$10$} 
	 \loigiai{
		 Tần số tích lũy của nhóm $3$ là $23$ nên ta có\\
		\begin{eqnarray*}
			&  &23=1+x-4+x^2-3x+11\\
			&\Leftrightarrow& x^2-2x-15=0\\ 
			&\Leftrightarrow&  \hoac{&{x=5} \\&{x=-3}.}
		\end{eqnarray*}
		 Do $x > 0$ nên $x=5$.
	} 
\end{ex}

% nguồn: Ok_TracNghiemLop11_Lan1
\begin{ex}%[1D5H1-2]%[ma:f9448e7d]
	Cho mẫy số liệu ghép nhóm về tuổi thọ (đơn vị tính là  năm) của một loại bóng đèn mới như sau.
	\begin{center}
		\begin{longtable}{|c|c|c|c|c|}
			\hline 
			Tuổi thọ (năm) & $[2;3{,}5)$ & $[3{,}5;5)$ & $[5;6{,}5)$ & $[6{,}5;8)$\\ 
			 \hline
			Số  bóng đèn & $8$ & $22$ & $35$ & $15$\\
			\hline 
		\end{longtable}
	\end{center}
	Nhóm chứa mốt của  mẫu số liệu là
	\choice 
		{$[2;3{,}5)$} 
		 {$[3{,}5;5)$}
		 {\True $[5;6{,}5)$}
		{$[6{,}5;8)$} 
	 \loigiai{
		Tần số lớn nhất là $35$ nên nhóm chứa  mốt là nhóm $[5;6{,}5)$.
	 }
\end{ex}

% nguồn: Ok_TracNghiemLop11_Lan1
\begin{ex}%[1D5H1-2]%[ma:efcc7641]
	Thống kê chỉ số chất lượng không khí $(AQI)$ tại một địa điểm vào các ngày trong tháng $6 / 2022$  được cho trong bảng sau:
	\begin{center}
		\begin{tabular}{|c|c|c|c|c|c|}
			\hline Chỉ số $AQI$ & {$[0;  50)$} & {$[50; 100)$} & {$[100; 150)$} & {$[150; 200)$} & Trên 200 \\
			\hline  Số ngày & $5$ & $11$ & $7$ & $4$ & $3$ \\
			 \hline
		\end{tabular}
	\end{center}
	 Chất lượng không khí được xem là tốt nếu $AQI$ nhỏ hơn $50$, là trung bình nếu $AQI$ từ $50$ đến dưới $100$. Trong tháng $6/2022$ tại địa điểm này có bao nhiêu ngày chất lượng không khí dưới mức trung bình?
	 \choice
		 {$15$}
		 {\True $14$}
		 {$17$}
		{$12$} 
	 \loigiai{
		\begin{enumerate}
			\item Trong tháng $6/2022$ có $5$ ngày chỉ số $AQI$  dưới $50$; $11$ ngày chỉ số $AQI$ từ $50$ đến dưới $100$; $7$ ngày chỉ số $AQI$ từ $100$ đến dưới $150$; $4$ ngày chỉ số $AQI$ từ $150$ đến dưới $200$; $3$ ngày chỉ số $AQI$ trên $200$.
			\item Số ngày chất lượng không khí  dưới mức trung bình là $7+4+3=14$.
		\end{enumerate}
	 }
\end{ex}

% nguồn: Ok_TracNghiemLop11_Lan1
\begin{ex}%[1D5H1-2]%[ma:1271dce2]
	Cho bảng phân bố tần  số ghép lớp
	\begin{center}
		\begin{tabular}{|c|c|c|c|c|c|c|}
			\hline Các lớp giá trị của X & {$[50; 52)$} & {$[52; 54)$} & {$[54; 56)$} & {$[56; 58)$}  & {$[58; 60)$} & Cộng \\
			\hline Tần số $n_i$ & $15$  & $20$ & $45$ & $15$ & $5$ & $100$ \\
			\hline 
		\end{tabular}
	\end{center}
	 Mệnh đề đúng là
	\choice 
		{Giá trị  trung tâm của lớp $\left[50;52\right)$ là $53$}
		{Tần số của lớp $\left[58;60\right)$ là  $95$}
		{Tần số của lớp $\left[52;54\right)$  là $35$}
		{\True Giá trị $50$ không phụ thuộc lớp $\left[54;56\right)$} 
	\loigiai{Vì $50\in  \left[54;56\right)$.
	} 
\end{ex}


%%%%%%%%%%%%%%%% ĐÚNG/SAI %%%%%%%%%%%%%%%%

% nguồn: Ok_dungsai_lop11_Lan1
\begin{ex}%[1D5N1-2]%[ma:c7c7eeb8]
 Thống kê tuổi thọ của các bóng đèn do một nhà máy sản xuất ta có bảng  số liệu sau
 \begin{center}
  \begin{tabular}{|c|c|c|c|c|c|}
   \hline 
   Tuổi thọ (giờ) & $[1200;1300)$ & $[1300;1400)$ & $[1400;1500)$ & $[1500;1600)$ & $[1600;1700)$  \\
   \hline 
   Số bóng & $15$ & $20$ &  $48$ & $42$ & $25$ \\
    \hline
  \end{tabular}
 \end{center}
  \choiceTF
  {\True Bảng số liệu trên gồm  $5$ nhóm}
  {Số lượng bóng đèn  là $120$ bóng đèn}
  {\True Giá trị đại diện của nhóm  $[1200;1300)$ là $1250$}
  {\True Độ dài của mỗi nhóm trong  mẫu số liệu ở trên là $100$}
  \loigiai{
  Ta  có bảng số liệu sau
  \begin{center}
   \begin{tabular}{|c|c|c|c|c|c|}
     \hline
    Tuổi thọ (giờ) & $[1200;1300)$ & $[1300;1400)$ & $[1400;1500)$  & $[1500;1600)$ & $[1600;1700)$ \\
    \hline 
    Giá trị đại diện & $1250$ & $1350$ & $1450$ &  $1550$ & $1650$ \\
    \hline 
    Độ dài của nhóm & $100$ &  $100$ & $100$ & $100$ & $100$ \\
     \hline
   \end{tabular}
  \end{center}
  \begin{itemchoice}
    \itemch Đúng. Bảng số liệu trên gồm $5$ nhóm là $[1200;1300)$, $[1300;1400)$, $[1400;1500)$, $[1500;1600)$, $[1600;1700)$.
   \itemch Sai. Số lượng bóng đèn là  $15+20+48+42+25=150$ bóng đèn.
   \itemch Đúng. Giá trị đại diện của nhóm  $[1200;1300)$ là $1250$.
   \itemch Đúng. Độ dài của  mỗi nhóm trong mẫu số liệu ở trên là $100$.
  \end{itemchoice}
  }
\end{ex}

% nguồn: Ok_dungsai_lop11_Lan1
\begin{ex}%[1D5N1-2]%[ma:5e36ced5]
 Một mẫu số liệu đươc cho ở dạng bảng tần số ghép  nhóm như sau
 \begin{center}
  \begin{tabular}{|c|c|c|c|c|c|}
   \hline 
   Nhóm & $[0{,}5;2{,}5)$ & $[2{,}5;4{,}5)$ & $[4{,}5;6{,}5)$ & $[6{,}5;8{,}5)$ &  $[8{,}5;10{,}5)$ \\
    \hline
   Tần số & $4$  & $7$ & $16$ & $8$ & $5$ \\
    \hline
  \end{tabular}
 \end{center}
 \choiceTF 
  {\True Cỡ mẫu của mẫu số liệu là  $40$}
  {\True Nhóm $[0{,}5;2{,}5)$ có giá trị đại diện  là $1{,}5$}
  {\True Nhóm $[4{,}5;6{,}5)$ có  giá trị đại diện là $5{,}5$}
  {Nhóm $[8{,}5;10{,}5)$ có độ dài lớn nhất  trong các nhóm còn lại}
  \loigiai{
  Bảng sau cho biết giá trị đại diện và  độ dài của mỗi nhóm
  \begin{center}
   \begin{tabular}{|c|c|c|c|c|c|}
    \hline 
    Nhóm & $[0{,}5;2{,}5)$ & $[2{,}5;4{,}5)$  & $[4{,}5;6{,}5)$ & $[6{,}5;8{,}5)$ & $[8{,}5;10{,}5)$ \\
     \hline
    Giá trị đại diện & $1{,}5$ & $3{,}5$ &  $5{,}5$ & $7{,}5$ & $9{,}5$ \\
    \hline 
    Độ  dài nhóm & $2$ & $2$ & $2$ & $2$ & $2$ \\
     \hline
   \end{tabular}
  \end{center}
  \begin{itemchoice}
   \itemch Đúng. Cỡ mẫu của  mẫu số liệu là $n=4+7+15+8+5=40$.
   \itemch Đúng. Nhóm $[0{,}5;2{,}5)$ có giá trị đại diện  là $1{,}5$.
   \itemch Đúng.  Nhóm $[4{,}5;6{,}5)$ có giá trị đại diện là $5{,}5$.
   \itemch  Sai. Độ dài các nhóm đều giống nhau.
  \end{itemchoice}
 } 
\end{ex}

% nguồn: Ok_dungsai_lop11_Lan2
\begin{ex}%[1D5N1-2]%[ma:0cfc6a9a]
	Số tiền thu từ heo đất của một trường THPT phát động ở 40 lớp thu được  trong năm học được tổng hợp qua bảng sau
	\begin{center}
		\begin{tabular}{|c|c|c|c|c|c|}
			 \hline
			$\textbf{Số tiền  (triệu đồng)}$ & $[1{,}5;1{,}7)$ & $[1{,}7;1{,}9)$ & $[1{,}9;2{,}1)$ & $[2{,}1;2{,}3)$ & $[2{,}3;2{,}5)$ \\
			\hline 
			$\textbf{Số lớp}$ & $5$ & $6$ & $13$ &  $7$ & $9$\\
			\hline 
		\end{tabular}
	\end{center}
	 \choiceTF
		 {\True Bảng số liệu trên gồm $5$ nhóm}
		{\True  Giá trị đại diện của nhóm $[1{,}9;2{,}1)$ là $2{,}0$}
		{\True Giá trị  đại diện của nhóm $[1{,}5;1{,}7)$ là $1{,}6$}
		{Độ dài của mỗi nhóm trong mẫu số liệu ở trên là  $0{,}3$}
	 \loigiai{
		\begin{itemchoice}
			\itemch \textbf{Đúng.}  \\
			Từ bảng dữ liệu, ta  thấy có $5$ nhóm số tiền thu từ heo đất.
			\itemch  \textbf{Đúng.} \\
			 Giá trị đại diện của nhóm $[1{,}9;2{,}1)$ là $\dfrac{1{,}9+2{,}1}{2}=2{,}0$.
			\itemch \textbf{Đúng.}\\ 
			Giá trị đại diện của nhóm $[1{,}5;1{,}7)$ là $\dfrac{1{,}5+1{,}7}{2}=1{,}6$. 
			 \itemch \textbf{Sai.} \\
			Độ dài của mỗi nhóm trong mẫu số liệu ở  trên là $1{,}7-1{,}5=0{,}2$.
		\end{itemchoice}
	} 
\end{ex}

% nguồn: Ok_dungsai_lop11_Lan2
\begin{ex}%[1D5N1-2]%[ma:39f803b4]
	Một mẫu số liệu được cho ở dạng  bảng tần số ghép nhóm như sau
	\begin{center}
		\begin{tabular}{|c|c|c|c|c|c|}
			\hline 
			$\textbf{Nhóm}$ & $[0;5)$  & $[5;10)$ & $[10;15)$ & $[15;20)$ & $[20;25)$ \\
			\hline 
			 $\textbf{Tần số}$ & $11$ & $31$ & $45$ & $21$ & $12$\\
			\hline 
		\end{tabular}
	\end{center}
	 \choiceTF
		{Mẫu  trên có $110$ số liệu}
		{\True Mẫu  trên chia thành $5$ nhóm}
		{\True Tần số của nhóm $[0;5)$  bằng $11$}
		{Tần số của nhóm $[20;25)$ là  cao nhất}
	 \loigiai{
		\begin{itemchoice}
			\itemch \textbf{Sai.} \\ 
			Mẫu trên có  $11+31+45+21+12=120$ số liệu (cỡ mẫu).
			 \itemch \textbf{Đúng.} \\
			Từ bảng dữ liệu, ta thấy mẫu  được chia thành $5$ nhóm với độ dài là $5$.
			\itemch  \textbf{Đúng.}\\
			Từ bảng dữ liệu, ta  thấy tần số của nhóm $[0;5)$ bằng $11$
			\itemch \textbf{Sai.}  \\
			Nhóm $[10;15)$ có tần số cao nhất là $45$. 
		\end{itemchoice}
	 }
\end{ex}

% nguồn: Ok_dungsai_lop11_Lan2
\begin{ex}%[1D5N1-2]%[ma:d42edeb2]
	Thống kê tuổi thọ của các bóng đèn do một nhà máy sản xuất  ta có bảng số liệu sau
	\begin{center}
		\begin{tabular}{|c|c|c|c|c|c|}
			 \hline
			$\textbf{Tuổi thọ  (giờ)}$ & $[1200;1300)$ & $[1300;1400)$ & $[1400;1500)$ & $[1500;1600)$ & $[1600;1700)$ \\
			 \hline
			$\textbf{Số bóng}$ &  $15$ & $20$ & $48$ & $42$ & $25$\\
			 \hline
		\end{tabular}
	\end{center}
	 \choiceTF
		{\True Bảng số liệu  trên gồm $5$ nhóm}
		 {Số lượng bóng đèn là $120$ bóng đèn}
		{\True Giá trị đại diện  của nhóm $[1200;1300)$ là $1250$}
		{\True Độ dài của  mỗi nhóm trong mẫu số liệu ở trên là $100$}
	 \loigiai{
		\begin{itemchoice}
			\itemch \textbf{Đúng.}  \\
			Từ bảng  dữ liệu, ta thấy có $5$ nhóm tuổi thọ.
			\itemch \textbf{Sai.}  \\
			Số lượng bóng đèn là $15+20+48+42+25=150$ bóng đèn. 
			\itemch \textbf{Đúng.}\\ 
			Giá trị đại diện của nhóm $[1200;1300)$ là $\dfrac{1200+1300}{2}=1250$. 
			\itemch \textbf{Đúng.}  \\
			Độ dài của  mỗi nhóm trong mẫu số liệu ở trên là $1700-1600=100$.
		\end{itemchoice}
	} 
\end{ex}

% nguồn: Ok_dungsai_lop11_Lan2
\begin{ex}%[1D5N1-2]%[ma:7649ab1b]
	Một trường trung học phổ thông có $36$ học sinh nam của khối $11$,  đo chiều cao của các bạn học sinh đó, người ta thu được mẫu số liệu sau (đơn vị: centimét).
	\begin{center}
		\begin{tabular}{|c|c|c|c|c|c|c|c|c|c|c|c|}
			\hline 
			$160$ & $161$ & $161$ & $162$  & $162$ & $162$ & $163$ & $163$ & $163$ & $164$ & $164$ & $164$ \\
			 \hline
			$164$ & $165$ & $165$ & $165$ & $165$ & $165$ & $166$ & $166$ & $166$ & $166$ & $167$ & $167$\\ 
			 \hline
			$168$ & $168$ & $168$ &  $168$ & $169$ & $169$ & $170$ & $171$ & $171$ & $172$ & $172$ & $174$\\
			\hline 
		\end{tabular}
	\end{center}
	 \choiceTF
		{\True  Giá trị lớn nhất $x_{\max}=174$}
		{\True Giá trị nhỏ nhất $x_{\min}=160$} 
		{Khoảng biến thiên $15$} 
		{Ta  có bảng tần số ghép nhóm
		\begin{center}
			\begin{tabular}{|c|c|c|c|c|c|}
				\hline 
				 $\textbf{Chiều cao}$ & $[160;163)$ & $[163;166)$ & $[166;169)$ & $[169;172)$ & $[172;175)$ \\
				\hline 
				$\textbf{Số học sinh}$ & $6$ &  $10$ & $12$ & $3$ & $3$\\
				\hline 
		\end{tabular}
	\end{center}}
	 \loigiai{
		\begin{itemchoice}
			\itemch  \textbf{Đúng.} \\
			Giá trị lớn nhất $x_{\max}=174$. 
			\itemch \textbf{Đúng.} \\ 
			 Giá trị nhỏ nhất $x_{\min}=160$.
			\itemch  \textbf{Sai.}\\
			Khoảng biến thiên là $R=x_{\max}-x_{\min}=174-160=14$. 
			\itemch \textbf{Sai.} \\ 
			 Số nhóm $k=5$, độ dài mỗi nhóm $L>\dfrac{R}{k}=\dfrac{14}{5}=2{,}8$.\\
			Chọn độ  dài mỗi nhóm $L=3$, ta có bảng tần số ghép nhóm sau
			\begin{center}
				\begin{tabular}{|c|c|c|c|c|c|}
					 \hline
					$\textbf{Chiều cao}$ & $[160;163)$  & $[163;166)$ & $[166;169)$ & $[169;172)$ & $[172;175)$ \\
					 \hline
					$\textbf{Số học sinh}$ & $6$ &  $12$ & $10$ & $5$ & $3$\\
					 \hline
				\end{tabular}
			\end{center}
		\end{itemchoice}
	 }
\end{ex}

% nguồn: Ok_dungsai_lop11_Lan2
\begin{ex}%[1D5N1-2]%[ma:c6a378f9]
	Một mẫu số liệu được cho ở dạng bảng tần số ghép  nhóm như sau
	\begin{center}
		\begin{tabular}{|c|c|c|c|c|c|}
			 \hline
			$\textbf{Nhóm}$ & $[0{,}5;2{,}5)$ & $[2{,}5;4{,}5)$ & $[4{,}5;6{,}5)$ & $[6{,}5;8{,}5)$ & $[8{,}5;10{,}5)$ \\ 
			 \hline
			$\textbf{Tần số}$ & $4$ & $7$ & $16$ & $8$ &  $5$\\
			 \hline
		\end{tabular}
	\end{center}
	 \choiceTF
		 {\True Cỡ mẫu của số liệu là $40$}
		{\True Nhóm $[0{,}5;2{,}5)$  có giá trị đại diện là $1{,}5$}
		{\True Nhóm $[4{,}5;5{,}5)$ có  giá trị đại diện là $5{,}5$}
		{Nhóm $[8{,}5;10{,}5)$ có  độ dài lớn nhất trong các nhóm}
	 \loigiai{
		\begin{itemchoice}
			\itemch \textbf{Đúng.}  \\
			Cỡ mẫu của mẫu số liệu  là $n=4+7+15+8+5=40$.
			\itemch \textbf{Đúng.}  \\
			Nhóm $[0{,}5;2{,}5)$ có giá trị  đại diện là $\dfrac{0{,}5+2{,}5}{2}=1{,}5$.
			\itemch \textbf{Đúng.}\\ 
			Nhóm $[4{,}5;6{,}5)$ có  giá trị đại diện là $\dfrac{4{,}5+6{,}5}{2}=5{,}5$.
			\itemch  \textbf{Sai.} \\
			Các nhóm  của mẫu số liệu có độ dài như nhau và bằng $2{,}5-0{,}5=2$.
		\end{itemchoice}
	} 
\end{ex}

% nguồn: Ok_dungsai_lop11_Lan2
\begin{ex}%[1D5N1-2]%[ma:475f7241]
	Cho mẫu số liệu  ghép nhóm sau
	\begin{center}
		\begin{tabular}{|c|c|c|c|c|c|}
			\hline 
			$\textbf{Điểm số  môn Toán}$ & $[0;2)$ & $[2;4)$ & $[4;6)$ & $[6;8)$ & $[8;10)$ \\
			 \hline
			$\textbf{Số học sinh đạt được}$ & $1$ & $6$ & $12$ & $14$  & $8$\\
			\hline 
		\end{tabular}
	\end{center}
	 \choiceTF
		{Cỡ mẫu của mẫu số liệu bằng $40$} 
		{\True Giá trị đại diện  của nhóm $[2;4)$ bằng $3$}
		{Độ dài của nhóm $[6;8)$ bằng $3$} 
		{\True Độ  dài của nhóm $[8;10)$ bằng $2$}
	\loigiai{ 
		\begin{itemchoice}
			 \itemch \textbf{Sai.}\\
			Cỡ mẫu của  mẫu số liệu là $1+6+12+14+8=41$.
			\itemch \textbf{Đúng.}\\ 
			Giá trị đại diện của nhóm $[2;4)$  bằng $\dfrac{2+4}{2}=3$.
			 \itemch \textbf{Sai.}\\
			Độ dài của nhóm $[6;8)$  bằng $8-6=2$.
			\itemch  \textbf{Đúng.}\\
			Độ  dài của nhóm $[8;10)$ bằng $10-8=2$.
		\end{itemchoice}
	 }
\end{ex}

% nguồn: Ok_dungsai_lop11_Lan2
\begin{ex}%[1D5N1-2]%[ma:a0ed7684]
	Bảng thống kê sau cho biết thời gian chạy (phút) của  $30$ vận động viên (VĐV) trong một giải chạy marathon
	\begin{center}
		\begin{tabular}{|c|c|c|c|c|c|c|c|c|c|c|c|c|}
			\hline 
			$\textbf{Thời gian}$ & $129$ & $130$ & $133$ & $134$ & $135$  & $136$ & $138$ & $141$ & $142$ & $143$ & $144$ & $145$\\
			\hline 
			$\textbf{Số VĐV}$ & $1$ & $2$  & $1$ & $1$ & $1$ & $2$ & $3$ & $3$ & $4$ & $5$ & $2$ & $5$\\
			 \hline
		\end{tabular}
	\end{center}
	\choiceTF 
		{\True Giá trị lớn  nhất $x_{\max}=145$}
		{\True Giá trị nhỏ nhất  $x_{\min}=129$}
		{Khoảng biến  thiên $15$}
		{Mẫu số liệu ghép  nhóm
		\begin{center}
			\begin{tabular}{|c|c|}
				 \hline
				$\textbf{Thời gian}$ & $\textbf{Số  VĐV}$ \\
				\hline  $[127{,}5;130{,}5)$ & $3$\\
				 \hline $[130{,}5;133{,}5)$ & $1$\\
				\hline $[133{,}5;136{,}5)$ & $4$\\ 
				\hline $[136{,}5;139{,}5)$ & $3$\\ 
				 \hline $[139{,}5;142{,}5)$ & $9$\\
				\hline $[142{,}5;145{,}5)$  & $10$\\
				 \hline
		\end{tabular}
	\end{center}}
	\loigiai{ 
		\begin{itemchoice}
			\itemch \textbf{Đúng.}  \\
			Giá trị lớn nhất  $x_{\max}=145$.
			\itemch \textbf{Đúng.}  \\
			Giá trị  nhỏ nhất $x_{\min}=129$.
			 \itemch \textbf{Sai.}\\
			Khoảng  biến thiên bằng $145-129=16$.
			\itemch \textbf{Sai.}  \\
			Ta có mẫu số liệu ghép  nhóm sau
			\begin{center}
				\begin{tabular}{|c|c|}
					 \hline
					 $\textbf{Thời gian}$ & $\textbf{Số VĐV}$ \\
					\hline  $[127{,}5;130{,}5)$ & $3$\\
					\hline  $[130{,}5;133{,}5)$ & $1$\\
					 \hline $[133{,}5;136{,}5)$ & $4$\\
					\hline  $[136{,}5;139{,}5)$ & $3$\\
					\hline $[139{,}5;142{,}5)$ & $7$\\ 
					\hline $[142{,}5;145{,}5)$ & $12$\\ 
					\hline 
				\end{tabular}
			\end{center}
		\end{itemchoice}
	} 
\end{ex}

% nguồn: Ok_dungsai_lop11_Lan1
\begin{ex}%[1D5H1-2]%[ma:35beaf7c]
 Một trường trung học phổ thông có $36$ học sinh  nam của khối $11$, đo chiều cao của các bạn học sinh đó, người ta thu được mẫu số liệu sau (đơn vị: centimét).
 \begin{center}
  \begin{tabular}{|l|l|l|l|l|l|l|l|l|l|l|l|}
   \hline 
   $160$ & $161$ & $161$ & $162$ & $162$ &  $162$ & $163$ & $163$ & $163$ & $164$ & $164$ & $164$ \\
    \hline
   $164$ & $165$ & $165$ & $165$ & $165$ & $165$ & $166$ & $166$  & $166$ & $166$ & $167$ & $167$ \\
    \hline
   $168$ & $168$ & $168$ &  $168$ & $169$ & $169$ & $170$ & $171$ & $171$ & $172$ & $172$ & $174$ \\
    \hline
  \end{tabular}
 \end{center}
  Khi đó
  \choiceTF
  {\True Giá trị lớn nhất $x_{\max}=174$} 
  {\True Giá trị nhỏ nhất $x_{\min}=160$} 
  {Khoảng biến thiên  là $15$}
  {Ta có bảng tần số ghép  nhóm sau:\\
  \begin{tabular}{|c|c|c|c|c|c|}
   \hline 
   Chiều cao & $[160; 163)$ & $[163; 166)$ & $[166; 169)$ &  $[169; 172)$ & $[172; 175)$ \\
   \hline 
   Số học sinh & $6$ & $10$  & $12$ & $3$ & $3$ \\
    \hline
 \end{tabular}
  }
  \loigiai{
  \begin{itemchoice}
   \itemch \textbf{Đúng}.\\ 
    Giá trị lớn nhất $x_{\max}=174$.
    \itemch \textbf{Đúng}.\\
   Giá trị  nhỏ nhất $x_{\min}=160$.
   \itemch  \textbf{Sai}.\\
   Khoảng biến thiên là $R=x_{\max}-x_{\min}=174-160=14$,  số nhóm $k=5$.
   \itemch \textbf{Sai}.\\ 
   Ta có bảng tần số ghép  nhóm sau:
   Độ dài  mỗi nhóm: $L>\dfrac{R}{k}=\dfrac{14}{5}=2{,}8$.\\
   Chọn độ dài mỗi nhóm $L=3$, ta có bảng tần số ghép nhóm sau: 
   \begin{center}
    \begin{tabular}{|c|c|c|c|c|c|}
      \hline
     Chiều cao & $[160; 163)$ & $[163; 166)$ & $[166; 169)$ & $[169; 172)$ & $[172;  175)$ \\
      \hline
     Số học sinh & $6$ & $12$ & $10$ & $5$ &  $3$ \\
      \hline
    \end{tabular}
   \end{center}
  \end{itemchoice}
 } 
\end{ex}

% nguồn: Ok_dungsai_lop11_Lan1
\begin{ex}%[1D5H1-2]%[ma:9546b886]
 Số tiền thu từ heo đất của một trường THPT phát động ở $40$ lớp thu được trong năm học được tổng hợp  qua bảng sau
 \begin{center}
  \begin{tabular}{|c|c|c|c|c|c|}
    \hline
   Số tiền (triệu đồng) & $[1{,}5;1{,}7)$ & $[1{,}7;1{,}9)$  & $[1{,}9;2{,}1)$ & $[2{,}1;2{,}3)$ & $[2{,}3;2{,}5)$ \\
   \hline 
   Số lớp & $5$ & $6$ & $13$ & $7$ & $9$ \\ 
    \hline
  \end{tabular}
 \end{center}
 \choiceTF 
  {\True Bảng số liệu trên gồm  $5$ nhóm}
  {\True  Giá trị đại diện nhóm $[1{,}9;2{,}1)$ là $2{,}0$}
  {\True  Giá trị đại diện nhóm $[1{,}5;1{,}7)$ là $1{,}6$}
  {Độ dài của  mỗi nhóm trong mẫu số liệu ở trên là $0{,}3$}
  \loigiai{
  Ta  có bảng sau
  \begin{center}
   \begin{tabular}{|c|c|c|c|c|c|}
     \hline
    Số tiền (triệu đồng) & $[1{,}5;1{,}7)$ & $[1{,}7;1{,}9)$ & $[1{,}9;2{,}1)$ & $[2{,}1;2{,}3)$ & $[2{,}3;2{,}5)$ \\ 
    \hline 
    Giá trị đại diện & $1{,}6$ & $1{,}8$ & $2{,}0$ & $2{,}2$ &  $2{,}4$ \\
    \hline 
    Độ dài của nhóm & $0{,}2$ & $0{,}2$ & $0{,}2$ &  $0{,}2$ & $0{,}2$ \\
     \hline
   \end{tabular}
  \end{center}
  \begin{itemchoice}
   \itemch Đúng. Bảng số liệu trên gồm $5$ nhóm là  $[1{,}5;1{,}7)$, $[1{,}7;1{,}9)$, $[1{,}9;2{,}1)$, $[2{,}1;2{,}3)$, $[2{,}3;2{,}5)$.
   \itemch Đúng. Giá trị đại diện nhóm $[1{,}9;2{,}1)$ là $2{,}0$. 
   \itemch Đúng. Giá trị đại diện nhóm $[1{,}5;1{,}7)$  là $1{,}6$.
    \itemch Sai. Độ dài của mỗi nhóm trong mẫu số liệu ở trên là $0{,}2$.
  \end{itemchoice}
 } 
\end{ex}

% nguồn: Ok_dungsai_lop11_Lan1
\begin{ex}%[1D5H1-2]%[ma:b3e88163]
 Bảng thống kê sau cho biết thời gian chạy (phút) của $30$ vận động viên (VĐV) trong  một giải chạy marathon:
 \begin{center}
  \begin{tabular}{|c|c|c|c|c|c|c|c|c|c|c|c|c|}
   \hline 
   Thời  gian & $129$ & $130$ & $133$ & $134$ & $135$ & $136$ & $138$ & $141$ & $142$ & $143$ & $144$ & $145$ \\
   \hline 
   Số VĐV & $1$ & $2$ & $1$ & $1$ & $1$  & $2$ & $3$ & $3$ & $4$ & $5$ & $2$ & $5$ \\
   \hline 
  \end{tabular}
 \end{center}
  Khi đó
  \choiceTF
  {\True Giá  trị lớn nhất $x_{\max}=145$}
  {\True Giá trị nhỏ nhất  $x_{\text{min}}=129$}
  {Khoảng biến thiên bằng  $15$}
   {Mẫu số liệu ghép nhóm:\\
  \begin{tabular}{|c|c|}
   \hline 
   Thời gian &  Số VĐV \\
    \hline
   $[127{,}5; 130{,}5)$ & $3$  \\
    \hline
   $[130{,}5; 133{,}5)$ &  $1$ \\
   \hline 
   $[133{,}5; 136{,}5)$ &  $4$ \\
    \hline
   $[136{,}5;  139{,}5)$ & $3$ \\
   \hline 
   $[139{,}5; 142{,}5)$ &  $9$ \\
   \hline 
   $[142{,}5; 145{,}5)$ & $10$ \\ 
    \hline
 \end{tabular}
 } 
  \loigiai{
  \begin{itemchoice}
   \itemch \textbf{Đúng}.\\ 
   Giá trị  lớn nhất $x_{\max}=145$.
   \itemch \textbf{Đúng}.\\ 
   Giá trị nhỏ nhất $x_{\min}=129$. 
   \itemch \textbf{Sai}.\\ 
   Khoảng biến  thiên: $R=x_{\max}-x_{\text{min}}=16$.\\
   Tổng độ dài sáu nhóm:  $6.3=18$.
   \itemch  \textbf{Sai}.\\
   Chọn đầu mút trái nhóm đầu tiên là $127{,}5$, đầu mút phải  nhóm cuối là $145{,}5$.\\
   Ta có  mẫu số liệu ghép nhóm sau:
   \begin{center}
    \begin{tabular}{|c|c|}
      \hline
     Thời  gian & Số VĐV \\
     \hline 
     $[127{,}5; 130{,}5)$ & $3$  \\
      \hline
     $[130{,}5; 133{,}5)$ & $1$ \\ 
     \hline 
     $[133{,}5;  136{,}5)$ & $4$ \\
     \hline 
     $[136{,}5; 139{,}5)$ & $3$  \\
     \hline 
     $[139{,}5;  142{,}5)$ & $7$ \\
      \hline
     $[142{,}5; 145{,}5)$ & $12$  \\
     \hline 
    \end{tabular}
   \end{center}
  \end{itemchoice}
  }
\end{ex}

% nguồn: Ok_dungsai_lop11_Lan1
\begin{ex}%[1D5H1-2]%[ma:f0b6d38c]
 Cho mẫu số liệu  ghép nhóm sau:
 \begin{center}
  \begin{tabular}{|c|c|c|c|c|c|}
   \hline 
   Điểm số môn Toán & $[0; 2)$ & $[2; 4)$ & $[4; 6)$ & $[6; 8)$ &  $[8; 10)$ \\
   \hline 
   Số học sinh  đạt được & $1$ & $6$ & $12$ & $14$ & $8$ \\
    \hline
  \end{tabular}
 \end{center}
 Khi đó 
  \choiceTF
  {Cỡ mẫu của mẫu số  liệu bằng $40$}
  {\True Giá trị đại  diện nhóm $[2; 4)$ bằng $3$}
  { Độ dài nhóm $[6; 8)$ bằng  $3$}
  {\True Độ  dài nhóm $[8; 10)$ bằng $2$}
 \loigiai{ 
  Giá trị đại  diện và độ dài mỗi nhóm được thể hiện ở bảng sau:
  \begin{center}
   \begin{tabular}{|c|c|c|c|c|c|}
    \hline 
    Điểm  số môn Toán & $[0; 2)$ & $[2; 4)$ & $[4; 6)$ & $[6; 8)$ & $[8; 10)$ \\
    \hline 
    Giá trị đại diện & $1$ & $3$ & $5$ & $7$  & $9$ \\
    \hline 
    Độ dài mồi nhóm & $2$ & $2$ & $2$ &  $2$ & $2$ \\
    \hline 
   \end{tabular}
  \end{center}
  \begin{itemchoice}
    \itemch \textbf{Sai}.\\
   Cỡ mẫu của mẫu số liệu là  $n=1+6+12+14+8=41$.
   \itemch  \textbf{Đúng}.\\
   Giá trị đại diện nhóm  $[2; 4)$ bằng $3$.
   \itemch  \textbf{Sai}.\\
   Độ  dài nhóm $[6; 8)$ bằng $3$.
   \itemch \textbf{Đúng}.\\ 
   Độ  dài nhóm $[8; 10)$ bằng $2$.
  \end{itemchoice}
 } 
\end{ex}

% nguồn: Ok_dungsai_lop11_Lan2
\begin{ex}%[1D5H1-2]%[ma:69e84ddf]
	Mẫu số liệu sau cho biết phân bố theo độ tuổi của dân số Việt  Nam năm $2019$.
	\begin{center}
		\begin{tabular}{|c|c|c|c|}
			 \hline
			Độ tuổi & Dưới $15$ & Từ $15$ đến dưới $65$ & Từ $65$ tuổi trở lên  \\
			\hline 
			Số người & $23\;371\;882$ &  $65\;429\;451$ & $7\;416\;651$ \\
			\hline 
		\end{tabular}
	\end{center}
	\choiceTF 
		{Mẫu số liệu đã cho là mẫu số liệu không ghép nhóm} 
		{\True Mẫu số liệu đã  cho có ba nhóm và ba tần số tương ứng}
		{Dân số của Việt Nam từ $15$ tuổi trở lên là $65\,429\,451$  người}
		{\True Dân số Việt  Nam năm $2019$ là $96\,208\,984$ người}
	 \loigiai{
		\begin{itemchoice}
			\itemch Mẫu số liệu  đã cho là mẫu số liệu ghép nhóm.
			\itemch Có ba nhóm: Dưới $15$ tuổi; Từ $15$ đến dưới $65$  tuổi; Từ $65$ tuổi trở lên.
			Có $23\,371\,882$ người dưới $15$ tuổi; $65\,420\,451$ người từ $15$  đến dưới $65$ tuổi và $7\,416\,651$ người từ $65$ tuổi trở lên.
			\itemch Dân số  của Việt Nam từ $15$ tuổi trở lên là $65\,429\,451+7\,416\,151=72\,845\,602$ người.
			\itemch Dân số Việt Nam  năm $2019$ là $23\;371\;882+65\;420\;451+7\;416\;651=96\;208\;984$
			người. 
	\end{itemchoice}}
\end{ex}

% nguồn: Ok_dungsai_lop11_Lan2
\begin{ex}%[1D5H1-2]%[ma:6abe58f0]
	Theo Tổ chức Y tế Thế giới (WHO), thiếu máu là tình trạng giảm lượng huyết sắc tố (Hb) dẫn tới sự thiếu cung cấp oxygen cho các mô trong cơ thể. Đối với nam giới trên $15$ tuổi, chỉ số Hb (đơn vị tính là gam/lít) lớn hơn hoặc bằng $130$ được xem là không bị thiếu máu, từ $110$ đến dưới $130$ là thiếu máu mức nhẹ, từ $80$ đến  dưới $110$ là thiếu máu mức vừa, dưới $80$ là mức nặng. Đo chỉ số Hb của một số học sinh nam lớp $11$ cho kết quả như bảng dưới đây. Xét tính đúng sai của
	các khẳng định sau 
	\begin{center}
		\begin{tabular}{|c|c|c|c|c|c|c|c|c|c|c|c|c|}
			\hline 
			$132$ & $135$ & $137$ & $131$ & $129$ &  $125$ & $140$ & $147$ & $138$ & $137$ & $128$ & $112$ \\
			\hline 
			$127$ & $129$ & $125$ & $98$ & $139$ & $138$ & $139$ & $141$ & $149$  & $105$ & $136$ & $133$ \\
			\hline 
			$137$ & $138$ & $108$ & $133$ & $136$ & $141$  & $144$ & $134$ & $136$ & $137$ & $142$ & \\
			 \hline
		\end{tabular}
	\end{center}
	 \choiceTF
		{\True Có $25$ học sinh không  bị thiếu máu}
		{\True Chỉ có $3$ học sinh thiếu  máu dưới mức trung bình}
		{Có $12$ học sinh thiếu máu ở  mức nhẹ}
		{\True Mẫu số liệu ghép nhóm  theo mức độ thiếu máu là}
		\begin{center}
			\begin{tabular}{|c|c|c|c|}
				 \hline
				Chỉ số Hb (gam/lít)  & $[80;110)$ & $[110;130)$ & Từ $130$ trở lên \\
				\hline 
				Số học sinh & $25$ & $7$ & $3$ \\ 
				\hline 
		\end{tabular}
	\end{center}
	\loigiai{Có $25$ học sinh có chỉ số Hb từ $130$ trở lên (không bị thiếu máu), $7$ học sinh có chỉ số Hb từ $110$ đến dưới $130$ (thiếu máu mức nhẹ), $3$ học sinh có  chỉ số Hb từ $80$ đến dưới $110$ (thiếu máu mức trung bình). Ta có mẫu số liệu ghép nhóm
		\begin{center}
			\begin{tabular}{|c|c|c|c|}
				\hline 
				Chỉ số Hb (gam/lít) & $[80;110)$ & $[110;130)$ &  Từ $130$ trở lên \\
				\hline 
				Số học sinh & $25$ & $7$ & $3$  \\
				\hline 
			\end{tabular}
		\end{center}
		\begin{itemchoice}
			\itemch Có $25$ học sinh  không bị thiếu máu.
			\itemch Chỉ có  $3$ học sinh thiếu máu dưới mức trung bình.
			\itemch  Có $7$ học sinh thiếu máu ở mức nhẹ.
			 \itemch Mẫu số liệu ghép nhóm theo mức độ thiếu máu là
			\begin{center}
				\begin{tabular}{|c|c|c|c|}
					\hline 
					Chỉ  số Hb (gam/lít) & $[80;110)$ & $[110;130)$ & Từ $130$ trở lên \\
					\hline 
					Số học sinh & $25$ & $7$ & $3$  \\
					\hline 
				\end{tabular}
			\end{center}
		\end{itemchoice}
	} 
\end{ex}

% nguồn: Ok_dungsai_lop11_Lan2
\begin{ex}%[1D5H1-2]%[ma:fc7fd145]
	Cho bảng số liệu khảo sát về tuổi thọ (đơn vị: nghìn giờ) của  một loại bóng đèn. Xét tính đúng sai của các khẳng định sau
	\begin{center}
		\begin{tabular}{|c|c|c|c|c|c|}
			\hline 
			Tuổi thọ & $[3;  5)$ & $[5; 7)$ & $[7; 9)$ & $[9;11)$ & $[11; 13)$ \\
			 \hline
			Số bóng đèn & $4$ & $20$ &  $26$ & $42$ & $8$ \\
			 \hline
		\end{tabular}
	\end{center}
	\choiceTF 
		{\True Mẫu số liệu đã cho có  là mẫu số liệu ghép nhóm}
		{\True Có $24$ bóng đèn có tuổi thọ ít hơn $7$ nghìn  giờ}
		{\True Có tất cả $100$ bóng đèn  được khảo sát}
		{\True Có $50$ bóng đèn có tuổi thọ từ $9$ nghìn giờ trở lên} 
	 \loigiai{
		\begin{itemchoice}
			\itemch Mẫu số liệu đã cho  là mẫu số liệu ghép nhóm.
			\itemch Số bóng đèn có tuổi thọ ít hơn  $7$ nghìn giờ là $4+20=24$ (bóng đèn).
			 \itemch Số bóng đèn được khảo sát là $4+20+26+42+8=100$ (bóng đèn).
			\itemch Số bóng đèn có tuổi thọ từ $9$ nghìn giờ trở lên là $42+8=50$ (bóng  đèn).
		\end{itemchoice}
	 }
\end{ex}

% nguồn: Ok_dungsai_lop11_Lan2
\begin{ex}%[1D5H1-2]%[ma:e77d08cd]
	Cho mẫu số liệu về  chiều cao của các học sinh lớp $11\text{B}$ (đơn vị: cm)
	\begin{center}
		\begin{tabular}{|c|c|c|c|c|c|c|c|c|c|c|c|c|}
			 \hline
			$156$ & $159$ & $160$ & $161$ & $162$ & $162$ & $163$ & $163$  & $164$ & $164$ & $164$ \\
			\hline 
			$165$ & $165$ & $165$ & $165$ & $165$ & $166$ &  $166$ & $166$ & $167$ & $167$ & $168$ \\
			 \hline
			$168$ & $168$ & $169$ & $169$ & $169$ & $170$ &  $170$ & $170$ & $171$ & $172$ & $173$ \\
			\hline 
		\end{tabular}
	\end{center}
	 \choiceTF
		{\True Khoảng biến thiên của mẫu số liệu là $17$} 
		{Có $12$ em học sinh có chiều cao từ  $166{,}5$ cm trở lên}
		{\True Có $28$ em học sinh có chiều cao  nhỏ hơn $170$ cm}
		{Mẫu số liệu ghép nhóm gồm năm nhóm có độ dài bằng nhau từ mẫu số  liệu trên là}
		\begin{center}
			\begin{tabular}{|c|c|c|c|c|c|}
				\hline 
				Chiều cao & $[156;159{,}5)$ & $[159{,}5;  163)$ & $[163; 166{,}5)$ & $[166{,}5;170)$ & $[170; 173{,}5)$ \\
				\hline 
				Số học sinh & $5$ & $4$ & $11$ &  $8$ & $4$ \\
				 \hline
		\end{tabular}
	\end{center}
	 \loigiai{
		\begin{itemchoice}
			 \itemch Giá trị nhỏ nhất của mẫu là $156$, giá trị lớn nhất là $173$ nên khoảng biến thiên là $173-156=17$.
			\itemch  Có $14$ em học sinh có chiều cao từ $166{,}5$ cm trở lên.
			\itemch  Có $28$ em học sinh có chiều cao nhỏ hơn $170$ cm.
			 \itemch Mẫu số liệu ghép nhóm gồm năm nhóm có độ dài bằng nhau từ mẫu số liệu trên là
			\begin{center}
				\begin{tabular}{|c|c|c|c|c|c|}
					 \hline
					Chiều cao & $[156;159{,}5)$ & $[159{,}5; 163)$  & $[163; 166{,}5)$ & $[166{,}5;170)$ & $[170; 173{,}5)$ \\
					\hline 
					Số học  sinh & $2$ & $4$ & $13$ & $8$ & $6$ \\
					\hline 
				\end{tabular}
			\end{center}
	\end{itemchoice}}
\end{ex}


%%%%%%%%%%%%%%%% TRẢ LỜI NGẮN %%%%%%%%%%%%%%%%

% nguồn: Ok_TraLoiNgan_Lop11_Lan2
\begin{ex}%[1D5N1-2]%[ma:b5b03c3a]
	Khảo sát thời gian xem ti vi trong một ngày của một số học  sinh khối $11$ thu được mẫu số liệu ghép nhóm sau
	\begin{center}
		\begin{tabular}{|c|c|c|c|c|c|}
			 \hline
			Thời gian (phút) & $[0;20)$ & $[20;40)$ & $[40;60)$ & $[60;80)$ & $[80;100)$ \\ 
			 \hline
			Số học sinh & $5$ & $9$ & $12$ & $10$ &  $6$ \\
			 \hline
		\end{tabular}
	\end{center}
	Hỏi có bao nhiêu học sinh xem ti vi từ $40$ phút trở lên trong một ngày? 
	\shortans{28} 
	\loigiai{Từ bảng số liệu ta thấy số học sinh có thời gian xem ti vi  từ $40$ phút trở lên trong một ngày có số liệu ở nhóm thứ $3{,}4$ và $5$ nên có $12+10+6=28$ (học sinh).}
\end{ex}

% nguồn: Ok_TraLoiNgan_Lop11_Lan2
\begin{ex}%[1D5N1-2]%[ma:05417edf]
	Một công ty xây dựng khảo sát khách hàng xem họ có nhu cầu mua  nhà ở mức giá nào. Kết quả khảo sát ghi lại ở bảng sau
	\begin{center}
		\begin{tabular}{|c|c|c|c|c|c|}
			\hline 
			Mức giá (triệu đồng/m$^2$) & $[10;14)$ & $[14;18)$ & $[18;22)$ & $[22;26)$ &  $[26;30)$ \\
			 \hline
			Số khách hàng & $54$ & $78$ & $120$ & $45$ & $12$ \\ 
			\hline 
		\end{tabular}
	\end{center}
	Hỏi có bao nhiêu khách hàng mua được nhà với mức giá nhỏ  hơn $22$ triệu đồng/m$^2$?
	\shortans{242} 
	\loigiai{Số khách hàng mua được nhà với mức giá nhỏ hơn $22$ triệu đồng/m$^2$  là $54+78+120=242$ (khách hàng).}
\end{ex}

% nguồn: Ok_TraLoiNgan_Lop11_Lan2
\begin{ex}%[1D5N1-2]%[ma:aecb64d7]
	Bảng số  liệu ghép nhóm sau cho biết chiều cao (cm) của $50$ học sinh lớp $11\text{A}$
	\begin{center}
		\begin{tabular}{|c|c|c|c|c|c|}
			\hline 
			Khoảng chiều cao (cm) & $[145;150)$ & $[150;155)$  & $[155;160)$ & $[160;165)$ & $[165;170)$ \\
			\hline 
			Số học  sinh & $7$ & $14$ & $10$ & $10$ & $9$ \\
			 \hline
		\end{tabular}
	\end{center}
	 Có bao nhiêu học sinh có chiều cao từ $165$ cm đến dưới $170$ cm?
	 \shortans{9}
	\loigiai{Số học sinh cao từ $165$  cm đến dưới $170$ cm là $9$ học sinh.}
\end{ex}

% nguồn: Ok_TraLoiNgan_Lop11_Lan2
\begin{ex}%[1D5N1-2]%[ma:57f78dde]
	Thời gian hoàn thành quãng đường $100$ m của $40$ học sinh lớp $11$ được cho trong bảng số liệu  dưới đây:
	\begin{center}
		\begin{tabular}{|c|c|c|c|c|}
			\hline 
			Thời gian (s) & $[15;17)$ & $[17;19)$ &  $[19;21)$ & $[21;23)$ \\
			\hline 
			Số học sinh & $8$ & $11$ &  $13$ & $8$ \\
			 \hline
		\end{tabular}
	\end{center}
	Hỏi có bao nhiêu học sinh có thời gian hoàn thành  đường chạy dưới $19$ s?
	\shortans{19} 
	\loigiai{Từ bảng số liệu ta thấy  số học sinh có thời gian hoàn thành dưới $19$ s có số liệu ở nhóm thứ nhất và hai nên có $8+11=19$ (học sinh).}
\end{ex}

% nguồn: Ok_TraLoiNgan_Lop11_Lan2
\begin{ex}%[1D5N1-2]%[ma:0e3dd91b]
	Điều tra về chiều cao của học sinh khối lớp $11$, ta  được mẫu số liệu sau
	\begin{center}
		\begin{tabular}{|c|c|}
			 \hline
			Chiều  cao (cm) & Số học sinh \\
			\hline 
			$[150;152)$ & $5$ \\ 
			\hline 
			 $[152;154)$ & $18$ \\
			\hline 
			$[154;156)$ &  $40$ \\
			 \hline
			$[156;158)$ &  $26$ \\
			 \hline
			$[158;160)$  & $8$ \\
			\hline 
			$[160;162)$ & $3$ \\ 
			\hline 
			Tổng & $N=100$  \\
			\hline 
		\end{tabular}
	\end{center}
	Mẫu số liệu ghép nhóm đã cho có tất cả  bao nhiêu nhóm?
	 \shortans{6}
	\loigiai{Mẫu  số liệu ghép nhóm đã cho có tất cả $6$ nhóm.}
\end{ex}

% nguồn: Ok_TraLoiNgan_Lop11_Lan2
\begin{ex}%[1D5N1-2]%[ma:0355e4f2]
	Doanh thu bán hàng trong $20$ ngày được lựa chọn  ngẫu nhiên của một cửa hàng được ghi lại ở bảng sau (đơn vị: triệu đồng)
	\begin{center}
		\begin{tabular}{|c|c|c|c|c|c|}
			 \hline
			Doanh  thu & $[5;7)$ & $[7;9)$ & $[9;11)$ & $[11;13)$ & $[13;15)$ \\
			 \hline
			Số ngày & $2$ & $7$ & $7$ &  $3$ & $1$ \\
			\hline 
		\end{tabular}
	\end{center}
	Có bao nhiêu  ngày cửa hàng đạt doanh thu từ $7$ triệu trở lên?
	 \shortans{18}
	\loigiai{Số ngày cửa  hàng đạt doanh thu từ $7$ triệu trở lên là $7+7+3+1=18$ (ngày).}
\end{ex}

% nguồn: Ok_TraLoiNgan_Lop11_Lan1
\begin{ex}%[1D5V1-2]%[ma:db70e44a]
	Một công ty may quần áo đồng phục học sinh cho biết  cỡ áo theo chiều cao của học sinh được tính như sau:
	\begin{center}
		\begin{tabular}{|>{\centering\arraybackslash}p{3cm}|>{\centering\arraybackslash}p{2cm}|>{\centering\arraybackslash}p{2cm}|>{\centering\arraybackslash}p{2cm}|>{\centering\arraybackslash}p{2cm}|>{\centering\arraybackslash}p{2cm}|}
			 \hline
			Chiều cao  $(\mathrm{cm})$&$[150; 160)$&$[160; 167)$&$[167; 170)$&$[170; 175)$&$[175; 180)$\\
			\hline 
			Cỡ  áo &$\mathrm{S}$&$\mathrm{M}$&$\mathrm{L}$&$\mathrm{XL}$&$\mathrm{XXL}$\\
			\hline 
		\end{tabular}
	\end{center}
	Công ty muốn ước lượng tỉ lệ các cỡ áo khi may cho học sinh lớp 11 đã đo chiều cao của 36 học sinh nam khối 11 của một trường và thu được mầu số  liệu sau (đơn vị là centimét):
	\begin{center}
		$\begin{array}{lllllllllllll}
		&160 &  161 & 161 & 162 & 162 & 162 & 163 & 163 & 163 & 164 & 164 & 164 \\
		& 164 & 165 & 165 & 165 & 165 & 165 & 166 &  166 & 166 & 166 & 167 & 167\\
		& 168 & 168 & 168 & 168 & 169 & 169 & 170 & 171 & 171 &  172 & 172 & 174
	\end{array}$
\end{center}
Công ty may $500$ áo đồng phục cho học sinh lớp 11 thì cỡ áo được may nhiều nhất là bao nhiêu  chiếc áo?
\shortans{139} 
 \loigiai{
	Bảng tần số ghép nhóm 
	\begin{center}
		\begin{tabular}{|>{\centering\arraybackslash}p{4cm}|>{\centering\arraybackslash}p{2cm}|>{\centering\arraybackslash}p{2cm}|>{\centering\arraybackslash}p{2cm}|>{\centering\arraybackslash}p{2cm}|>{\centering\arraybackslash}p{2cm}|}
			 \hline
			Nhóm Chiều cao &$[150; 160)$&$[160; 167)$&$[167;  170)$&$[170; 175)$&$[175; 180)$\\
			 \hline
			Số lượng  học sinh&$4$&$10$&$7$&$10$&$5$\\
			 \hline
		\end{tabular}
	\end{center}
	Tỉ lệ học  sinh trong mỗi nhóm và số lượng áo cần may với tổng số học sinh là: $36$
	 \renewcommand\arraystretch{1.6}
	\begin{center}
		\begin{tabular}{|>{\centering\arraybackslash}p{4cm}|>{\centering\arraybackslash}p{2cm}|>{\centering\arraybackslash}p{2cm}|>{\centering\arraybackslash}p{2cm}|>{\centering\arraybackslash}p{2cm}|>{\centering\arraybackslash}p{2cm}|}
			 \hline
			 Nhóm Chiều cao &$[150; 160)$&$[160; 167)$&$[167; 170)$&$[170; 175)$&$[175; 180)$\\
			 \hline
			Số lượng  học sinh&$4$&$10$&$7$&$10$&$5$\\
			 \hline
			 Tỉ lệ &$\dfrac{4}{36}$&$\dfrac{10}{36}$&$\dfrac{7}{36}$&$\dfrac{10}{36}$&$\dfrac{5}{36}$\\
			\hline 
		\end{tabular}
	\end{center}
	Với tổng số áo cần may là  $500$, nên các cỡ áo cần may, làm tròn.\\
	$-$ Số lượng áo cần may  cỡ $S$ là $\dfrac{4}{36} \times 500 \approx 56$ chiếc.\\
	$-$ Số lượng áo cần may cỡ $M$ là $\dfrac{10}{36}  \times 500 \approx 139$ chiếc.\\
	$-$ Số lượng áo cần may cỡ $L$ là  $\dfrac{7}{36} \times 500 \approx 97$ chiếc.\\
	$-$ Số  lượng áo cần may cỡ $XL$ là $\dfrac{10}{36} \times 500 \approx 139$ chiếc.\\
	$-$ Số lượng áo cần may cỡ $XXL$ là $\dfrac{5}{36} \times 500 \approx 69$  chiếc.\\
	Vậy cỡ áo được may nhiều nhất là $M$ và $XL$, mỗi loại $139$  chiếc.
} 
\end{ex}

% nguồn: Ok_TraLoiNgan_Lop11_Lan1
\begin{ex}%[1D5V1-2]%[ma:b894f071]
	 Điểm thi THPTQG $2023$ Môn Toán được xếp chia theo bảng số liệu ghép nhóm sau
	\begin{center}
		\begin{tabular}{|c|c|c|c|c|c|}
			 \hline
			Điểm & $[0;2)$ & $[2;4)$  & $[4;6)$ & $[6;8)$ & $[8;10]$\\
			\hline 
			Số thí sinh (ngàn người) & $4$  & $104$ & $262$ & $479$ & $151$\\
			\hline 
		\end{tabular}
	\end{center}
	Gọi $a$ là số trung bình của mẫu số liệu ghép nhóm, $b$ là số trung vi, $c$ là mốt. Tính $P=abc$  (làm tròn các số $a,b,c$ đến hàng phần nghìn, $P$ làm tròn đến hàng đơn vị).
	 \shortans{282}
	\loigiai{ 
		Số trung bình của mẫu số  liệu ghép nhóm là\\
		 $\overline{x}=\dfrac{{{m}_{1}}{{x}_{1}}+\ldots +{{m}_{5}}{{x}_{5}}}{n}=\dfrac{4\cdot1+104\cdot3+262\cdot5+479\cdot7+151\cdot9}{1000}=6{,}338.$\\
		Số trung vị của mẫu  số liệu ghép nhóm nằm ở nhóm $[6;8)$, do đó số trung vị sẽ là
		\begin{eqnarray*}
			{{M}_{e}}&=&{{a}_{p}}+\dfrac{\dfrac{n}{2}-\left( {{m}_{1}}+\ldots {{m}_{p-1}} \right)}{{{m}_{p}}}\cdot\left(  {{a}_{p+1}}-{{a}_{p}} \right) \\
			&=&6+\frac{\dfrac{1000}{2}-\left( 4+104+262  \right)}{479}\cdot\left( 8-6 \right)\approx 6{,}543.
		\end{eqnarray*}
		Mốt của mẫu số liệu  ghép nhóm
		\begin{eqnarray*}
			{{M}_{o}}&=&{{a}_{j}}+\dfrac{\left( {{m}_{j}}-{{m}_{j-1}} \right)}{\left( {{m}_{j}}-{{m}_{j-1}} \right)+\left( {{m}_{j}}-{{m}_{j+1}} \right)}\cdot h \\ 
			&=&6+\dfrac{479-262}{\left( 479-262 \right)+\left( 479-151 \right)}\cdot2\approx  6{,}796. \\
		\end{eqnarray*}
		$P=abc=6{,}338\cdot6,543\cdot6{,}796\approx  282$.
	 }
\end{ex}

% nguồn: Ok_TraLoiNgan_Lop11_Lan1
\begin{ex}%[1D5V1-2]%[ma:5b08bbec]
	Mẫu số liệu sau cho  biết kết quả kiểm tra môn Toán của lớp $11$A năm 2022.
	\begin{center}
		\begin{tabular}{|c|c|c|c|c|}
			 \hline
			Điểm số &  $[3;5)$&$[5;7)$&$[7;9)$&$[9;11)$\\
			 \hline
			Số  học sinh &5&18&10&7\\
			\hline 
		\end{tabular}
	\end{center}
	Số  học sinh của lớp $11A$ là bao nhiêu?
	\shortans{40} 
	 \loigiai{
		Mẫu số liệu  đã cho là mẫu số liệu ghép nhóm.\\
		Có  $5$ em từ $3$ đến dưới $5$ điểm; có 18 em từ $5$ đến dưới $7$ điểm; có $10$ em từ $7$ đến dưới $9$ điểm; có $7$ em từ $9$ đến dưới $11$ điểm.\\
		Số học sinh của lớp $11A$ là $5+18+10+7=40$ em. 
	 }
\end{ex}

% nguồn: Ok_TraLoiNgan_Lop11_Lan1
\begin{ex}%[1D5V1-2]%[ma:0d2904e9]
	Một  công ty may quần áo đồng phục học sinh cho biết cỡ áo theo chiều cao của học sinh được tính như sau:
	\begin{center}
		\begin{tabular}{|c|c|c|c|c|c|}
			 \hline
			\textbf{Chiều  cao (cm)}& $[150;160)$ & $[160;167)$ & $[167;170)$ & $[170;175)$ & $[175;180)$\\
			 \hline
			\textbf{Cỡ áo}& S & M  & L & XL & XXL\\
			\hline 
		\end{tabular}
	\end{center}
	Công ty muốn ước lượng tỉ lệ các cỡ áo khi may cho học sinh lớp 11 đã đo chiều cao của 36 học sinh nam khối 11 của một trường và thu được mẫu số liệu sau (đơn vị là centimét): 
	\begin{center}
		\begin{tabular}{ccccccccccccc}
			$160$ & $161$ & $161$ & $162$ & $162$ & $162$  & $163$ & $163$ & $163$ & $164$ & $164$ & $164$ & $164$ \\
			$165$ & $165$ & $165$ & $165$ & $165$ & $166$ & $166$ & $166$ & $166$ & $167$ & $167$ & $168$ & $168$ \\ 
			$168$ & $168$ & $169$ & $169$ & $170$ & $171$ & $171$ & $172$ & $172$ &  $174$ & & & \\
		\end{tabular}
	\end{center}
	Công ty may $500$ áo đồng phục cho học sinh lớp $11$ thì nên may số lượng áo cỡ XL là bao  nhiêu chiếc?
	 \shortans{83}
	 \loigiai{
		Đếm số giá trị thuộc mỗi nhóm, ta lập được bảng tần số ghép nhóm của mẫu số liệu với các nhóm  đã cho ở bảng trên như sau:
		\begin{center}
			\begin{tabular}{|c|c|c|c|c|c|}
				 \hline
				\textbf{Chiều cao (cm)}& $[150;160)$ & $[160;167)$ & $[167;170)$ &  $[170;175)$ & $[175;180)$ \\
				\hline 
				\textbf{Số học sinh}& $0$  & $22$ & $8$ & $6$ & $0$ \\
				 \hline
			\end{tabular}
		\end{center}
		Công ty may $500$ áo đồng phục cho học sinh lớp $11$ thì nên may số  lượng áo cỡ XL là \[\dfrac{6}{36}\cdot 500 \approx 83\,\text{(chiếc).}\]
	 }
\end{ex}


%%%%%%%%%%%%%%%% TỰ LUẬN %%%%%%%%%%%%%%%%

% nguồn: Ok_TuanLuan_Toan11_Lan1
\begin{bt}%[1D5N1-2]%[ma:e4b868ca]
	Tìm cân nặng trung bình của học sinh lớp $11D$ cho trong bảng sau: 
	\begin{center}
		\begin{tabular}{|c|c|c|c|c|c|c|}
			 \hline
			Cân nặng (kg)& $[40{,}5;45{,}5)$ & $[45{,}5;50{,}5)$ & $[50{,}5;55{,}5)$ & $[55{,}5;60{,}5)$ & $[60{,}5;65{,}5)$ & $[65{,}5;70{,}5)$  \\
			\hline 
			Số họ sinh& $10$ & $7$ & $16$  & $4$ & $2$ & $3$ \\
			 \hline
		\end{tabular}
	\end{center}
	 \loigiai{
		Trong mỗi khoảng cân nặng, giá trị đại diện trung bình  cộng của giá trị hai đầu mút nên ta có bảng sau:
		\begin{center}
			\begin{tabular}{|c|c|c|c|c|c|c|}
				\hline 
				Cân nặng (kg)& $43$ & $48$ &  $53$ & $58$ & $63$ & $65$ \\
				 \hline
				Số họ sinh& $10$  & $7$ & $16$ & $4$ & $2$ & $3$ \\
				 \hline
			\end{tabular}
		\end{center}
		Tổng số học sinh là  $n=42$.\\
		Cân nặng trung bình của học sinh lớp  11D là
		$$\overline{x}=\dfrac{10\cdot43+7\cdot48+16\cdot43+4\cdot58+2\cdot63+3\cdot68}{42}\approx 51{,}81 \text{  (kg).}$$
	} 
\end{bt}

% nguồn: Ok_TuanLuan_Toan11_Lan2
\begin{ex}%[1D5N1-2]%[ma:bf111843]
	Tìm cân nặng trung bình  của học sinh lớp $11D$ cho trong bảng sau:
	\begin{center}
		\begin{tabular}{|c|c|c|c|c|c|c|}
			\hline 
			Cân  nặng (kg)& $[40{,}5;45{,}5)$ & $[45{,}5;50{,}5)$ & $[50{,}5;55{,}5)$ & $[55{,}5;60{,}5)$ & $[60{,}5;65{,}5)$ & $[65{,}5;70{,}5)$ \\
			 \hline
			Số họ sinh& $10$ & $7$ & $16$ & $4$ & $2$ &  $3$ \\
			 \hline
		\end{tabular}
	\end{center}
	\loigiai{ 
		Trong mỗi khoảng cân nặng, giá trị đại diện trung bình  cộng của giá trị hai đầu mút nên ta có bảng sau:
		\begin{center}
			\begin{tabular}{|c|c|c|c|c|c|c|}
				\hline 
				 Cân nặng (kg)& $43$ & $48$ & $53$ & $58$ & $63$ & $65$ \\
				 \hline
				Số họ sinh& $10$ & $7$ & $16$  & $4$ & $2$ & $3$ \\
				\hline 
			\end{tabular}
		\end{center}
		Tổng số học sinh là  $n=42$.\\
		Cân nặng trung bình của học  sinh lớp 11D là
		$$\overline{x}=\dfrac{10\cdot43+7\cdot48+16\cdot43+4\cdot58+2\cdot63+3\cdot68}{42}\approx 51{,}81 \text{ (kg).}$$ 
	} 
\end{ex}

% nguồn: Ok_TuanLuan_Toan11_Lan2
\begin{ex}%[1D5N1-2]%[ma:436c8e71]
	Cho mẫu số liệu ghép nhóm về thời gian (phút) đi từ  nhà đến nơi làm việc của các nhân viên một công ty như sau
	\begin{center}
		\begin{tabular}{|c|c|c|c|c|c|c|c|}
			 \hline
			 Thời gian & $[15;20)$ & $[20;25)$ & $[25;30)$ & $[30;35)$ & $[35;40)$ & $[40;45)$ & $[45;50)$ \\
			\hline 
			Số nhân viên & $6$ & $14$ &  $25$ & $37$ & $21$ & $13$ & $9$ \\
			 \hline
		\end{tabular}
	\end{center}
	Đọc và giải thích mẫu  số liệu này.
	\loigiai{Mẫu số liệu được chia  thành $7$ nhóm:\\
		 Có $6$ nhân viên có thời gian đi từ nhà đến nơi làm việc là từ $15$ đến dưới $20$ phút\\
		Có $14$  nhân viên có thời gian đi từ nhà đến nơi làm việc là từ $20$ đến dưới $25$ phút\\
		Có $25$ nhân viên có thời gian đi từ nhà đến nơi làm việc là từ $25$ đến  dưới $30$ phút\\
		Có $37$ nhân viên có thời gian đi từ nhà đến nơi làm  việc là từ $30$ đến dưới $35$ phút\\
		 Có $21$ nhân viên có thời gian đi tư nhà đên nơi làm việc là từ $35$ đên dưới $40$ phút\\
		Có $13$ nhân viên có thời gian đi từ  nhà đến nơi làm việc là từ $40$ đến dưới $45$ phút\\
	Có $9$ nhân viên có thời gian đi từ nhà đến nơi làm việc là  từ $45$ đến dưới $50$ phút.}
\end{ex}

% nguồn: Ok_TuanLuan_Toan11_Lan2
\begin{ex}%[1D5H1-2]%[ma:53d14efa]
	Độ bão hoà oxygen trong máu (còn được gọi là chỉ số $SpO_2$) biểu thị cho tỉ lệ hemoglobin có oxygen trên tổng lượng hemoglobintrong máu. Chỉ số $SpO_2$ (đơn vị đo là $\%$) từ $97-99$ là oxygen trong máu tốt, $94-96$ là oxygen trong máu trung bình, $90-93$ là oxygen trong máu thấp, dưới $90$ là trường hợp cấp cứu trên lâm sàng (Theo Vinmec.com). Đo chỉ  số $SpO_2$ ở một số bệnh nhân Covid-$19$ người ta thu được kết quả sau:
	\begin{center}
		\begin{tabular}{|c|c|c|c|}
			 \hline
			$SpO_2 (\%)$ & $90-93$ & $94-96$ & $97-99$ \\ 
			 \hline
			Số bệnh nhân & $12$ & $31$ & $7$ \\ 
			 \hline
		\end{tabular}
	\end{center}
	\begin{enumerate}
		\item Cho biết  các nhóm số liệu và tần số tương ứng.
		\item Tính số  trung bình, trung vị và giải thích ý nghĩa của các giá trị thu được.
	\end{enumerate}
	 \loigiai{
		\begin{enumerate}
			\item Có $3$ nhóm số liệu gồm $90-93$, $94-96$, $97-99$ với tần số tương ứng là $12$, $31$, $7$. 
			\item Trước hết ta hiệu chỉnh các nhóm số liệu và  thu được bảng thống kê sau
			\begin{center}
				\begin{tabular}{|c|c|c|c|}
					 \hline
					$SpO_2 (\%)$ & $[89{,}5;93{,}5)$ & $[93{,}5;96,5)$ &  $[96,5;99,5)$ \\
					\hline 
					Số bệnh nhân  & $12$ & $31$ & $7$ \\
					\hline 
				\end{tabular}
			\end{center}
			Các giá trị đại diện cho các nhóm số liệu tương ứng là $91{,}5$; $95$ và  $98$.\\
			Cỡ mẫu $n=12+31+7=50$. Do đó số trung bình là $\overline{x}=\dfrac{1}{50} (12\cdot91,5+31\cdot95+7\cdot98)=94{,}58$.\\ 
			Do có $\dfrac{n}{2}=25$  giá trị nhỏ hơn trung vị nên trung vị thuộc nhóm $[93{,}5;96,5)$.\\
			Ta có $a_p=93{,}5$; $a_{p+1}=96,5$; $m_1+\cdots+m_{p-1}=12$; $m_p=31$. Do đó, trung vị là  $$M_e=93{,}5+\dfrac{25-12}{31} \cdot 3\approx 94{,}76.$$
			Như vậy, chỉ số $SpO_2$ trung  bình của $50$ bệnh nhân là $94{,}58$; có $25$ bệnh nhân có chỉ số $SpO_2$ nhỏ hơn $94{,}76$ và $25$ bệnh nhân có chỉ số $SpO_2$ lớn hơn $94{,}76$.\\
			Nhận xét: nếu chỉ tính số trung bình của mẫu số liệu ghép nhóm thì ta không cần  hiệu chỉnh nhóm rời rạc $k_1-k_2$, chọn giá trị đại diện là $\dfrac{k_1+k_2}{2}$.
		\end{enumerate}
	 }
\end{ex}

% nguồn: Ok_TuanLuan_Toan11_Lan2
\begin{ex}%[1D5H1-2]%[ma:b5b1c078]
	Quãng đường (km) từ nhà đến nơi làm việc của $40$ công nhân  một nhà máy được ghi lại như
	Sau 
	\begin{center}
		\begin{tabular}{|c|c|c|c|c|c|c|c|c|c|}
			\hline  $5$&$3$&$10$&$20$&$25$&$11$&$13$&$7$&$12$&$31$\\
			\hline $19$&$10$&$12$&$17$&$18$&$11$&$32$&$17$&$16$&$2$\\ 
			\hline  $7$&$9$&$7$&$8$&$3$&$5$&$12$&$15$&$18$&$3$\\
			 \hline $12$&$14$&$2$&$9$&$6$&$15$&$15$&$7$&$6$&$12$\\
			 \hline
		\end{tabular}
	\end{center}
	Ghép nhóm mẫu số liệu trên thành các khoảng có  độ rộng bằng nhau, khoảng đầu tiên $[0;5)$. Tần suất nhóm $[10;15)$ bằng bao nhiêu?
	 \loigiai{
		\textbf{\textit{Phương pháp giải:}} Lập bảng phân bố tần số ghép nhóm từ đó tính được tần suất của mỗi nhóm Ta chia thành các nhóm có độ dài bằng $5$. Ta sẽ chọn  đầu mút phải của nhóm cuối cùng là $35$.\\
		Ta có bảng ghép nhóm  sau
		\begin{center}
			\begin{tabular}{|c|c|c|c|c|c|c|c|}
				\hline\begin{tabular}{c}
				Quãng  đường\\
				 $(km)$.
			\end{tabular}&{$[0;5)$}&{$[5;10)$}&{$[10;15)$}&{$[15;20)$}&{$[20;25)$}&{$[25;30)$}&{$[30;35)$}\\
			\hline Số công  nhân&$6$&$10$&$11$&$9$&$1$&$1$&$2$\\
			\hline 
		\end{tabular}
	\end{center}
	Tần suất của  nhóm $[10;15)$ là $\dfrac{11}{40}\cdot 100=27{,}5\%$.
} 
\end{ex}

% nguồn: Ok_TuanLuan_Toan11_Lan2
\begin{ex}%[1D5H1-2]%[ma:b1abad55]
	Một công ty may quần áo đồng phục học  sinh cho biết cỡ áo theo chiều cao của học sinh được tính như sau
	\begin{center}
		\begin{tabular}{|c|c|c|c|c|c|}
			 \hline
			 Chiều cao (cm)& $[150;160)$& $[160;167)$ & $[167;170$& $[170;175)$& $[175;180)$\\
			 \hline
			Cỡ áo  & S & M& L & XL&XXL\\
			\hline 
		\end{tabular}
	\end{center}
	Công ty muốn ước lượng tỉ lệ các cỡ áo  khi may cho học sinh lớp $11$ đã đo chiều cao của $36$ học sinh nam khối $11$ của một trường và thu được mẫu số liệu sau (đơn vị là centimét)
	\begin{center}
		\begin{tabular}{|c|c|c|c|c|c|c|c|c|c|c|c|c|}
			 \hline
			$160$& $161$& $161$&$162$& $162$& $162$&$163$&$163$&$163$&$164$&$164$&$164$&$164$\\ 
			 \hline
			$165$&  $165$& $165$& $165$& $165$&$166$&$166$&$166$&$166$&$167$&$167$&$168$&$168$\\
			 \hline
			$168$  & $168$& $169$& $169$&$170$&$171$&$171$&$172$&$172$&$174$&&&\\
			\hline 
		\end{tabular}
	\end{center}
	\begin{enumerate}[1)]
		\item Lập bảng tần số ghép nhóm của mẫu số liệu với các nhóm đã cho  ở bảng trên.
		\item Công ty may  $500$ áo đồng phục cho học sinh lớp $11$ thì nên may số lượng áo theo mỗi cỡ là bao nhiêu chiếc?
	\end{enumerate}
	\loigiai{ 
		\begin{enumerate}[1)]
			\item Bảng tần số ghép nhóm của mẫu số liệu với các nhóm  đã cho ở bảng trên
			\begin{center}
				\begin{tabular}{|c|c|c|c|}
					 \hline
					 Chiều cao (cm)& $[160;167)$ & $[167;170$& $[170;175)$\\
					\hline 
					Cỡ áo &  $22$ & $8$& $6$\\
					\hline 
				\end{tabular}
			\end{center}
			\item Tỉ lệ học sinh mặc vừa cỡ M  là $(22: 36)\times 100=61{,}11\%$.
			\begin{itemize}
				\item Tỉ lệ học sinh mặc vừa cỡ L là $(8:  36) \times100=22{,}22\%$.
				\item Tỉ lệ học  sinh mặc vừa cỡ XL là $(6: 36) \times 100=16{,}67\%$.
				\item Số lượng áo cỡ M  nên may là $61{,}11\%\times 500=306$.
				\item Số lượng áo cỡ L  nên may là $22{,}22\%\times 500=111$.
				\item Số  lượng áo cỡ XL nên may là $16{,}67\% \times 500=83$.
			\end{itemize}
		\end{enumerate}
	} 
\end{ex}

% nguồn: Ok_TuanLuan_Toan11_Lan2
\begin{ex}%[1D5H1-2]%[ma:54c4a2a3]
	Mẫu số liệu dưới đây ghi lại tốc độ của $42$ ô tô  khi đi qua một trạm đo tốc độ (đơn vị: km/h):
	\begin{center}
		\begin{tabular}{ccccccc}
			47{,}5  & 49{,}5 & 46 & 51 & 52{,}5 & 45 & 61 \\
			42 & 67 & 48 & 63 & 65 &  62{,}5 & 49{,}5 \\
			43{,}5 &  41 & 57{,}5 & 63{,}5 & 56{,}5 & 53 & 48 \\
			61{,}5 & 46 & 57 & 69 & 44{,}5 & 52 & 50  \\
			 45 & 55 & 47 & 60 & 67{,}5 & 62 & 58 \\
			56 & 51{,}5 & 57{,}5 &  59 & 52 & 43 & 56. \\
		\end{tabular}
	\end{center}
	Lập bảng tần số  ghép nhóm bao gồm cả tần số tích luỹ cho mẫu số liệu trên có sáu nhóm ứng với sáu nửa khoảng:
	 \[[40;45),[45;50),[50;55),[55;60),[60;65),[65;70).\]
	\loigiai{ 
		Bảng tần  số ghép nhóm bao gồm cả tần số tích luỹ như ở Bảng trên:
		\begin{center}
			\begin{tabular}{|c|c|c|}
				\hline 
				Nhóm & Tần số & Tần  số tích luỹ \\
				 \hline
				$[40; 45)$ & $5$  & $5$ \\
				\hline 
				$[45; 50)$ & $10$  & $15$ \\
				\hline 
				$[50; 55)$  & $7$ & $22$ \\
				\hline 
				$[55; 60)$ &  $9$ & $31$ \\
				 \hline
				$[60; 65)$ & $7$ & $38$ \\ 
				\hline 
				$[65; 70)$ & $4$ & $42$ \\ 
				 \hline
				 & $n=42$ & \\
				 \hline
			\end{tabular}
		\end{center}
	} 
\end{ex}

% nguồn: Ok_TuanLuan_Toan11_Lan2
\begin{ex}%[1D5H1-2]%[ma:5dc8d470]
	Thời gian ra sân (giờ) của một cựu cầu thủ ở  giải ngoại hạng Anh qua các thời kì được cho như sau\\
	 \[
	653\quad632\quad609\quad572\quad565\quad535\quad516\quad514\quad508\quad505\quad504\quad504\quad503\quad499\quad496\quad492 
	\] 
	\begin{flushright}
		(Theo: https://www.premierleague.com/) 
	\end{flushright}
	Hãy chuyển mẫu số liệu sang dạng ghép nhóm với  bảy nhóm có độ dài bằng nhau.
	 \loigiai{
		Giá trị lớn  nhất của mẫu số liệu là $653$, giá trị nhỏ nhất của mẫu số liệu là $492$.\\
		Khoảng biến thiên là $653-492=161$.\\ 
		Ta chia thành $7$ nhóm có độ  dài bằng nhau. \\
		Ví dụ chọn độ dài mỗi nhóm là $25$, đầu mút trái là $485$, khi đó  ta có mẫu số liệu ghép nhóm như sau
		\begin{center}
			\begin{tabular}{|c|c|c|c|c|c|c|c|}
				\hline Thời gian&{$[485;510)$}&{$[510;535)$}&{$[535;560)$}&{$[560;585)$}&{$[585;610)$}&{$[610;635)$}&{$[635;660)$}\\ 
				\hline  Sồ cầu thủ&$8$&$2$&$1$&$2$&$1$&$1$&$1$\\
				 \hline
			\end{tabular}
		\end{center}
	} 
\end{ex}

% nguồn: Ok_TuanLuan_Toan11_Lan2
\begin{ex}%[1D5H1-2]%[ma:da3d8deb]
	Tính giá trị đại điện và độ dài của mỗi nhóm trong mẫu số liệu ở bảng  sau:
	\begin{center}
		\begin{tabular}{|c|c|c|c|c|c|}
			\hline 
			Khoảng tuổi & $[20;30)$ & $[30;40)$ & $[40;50)$ & $[50;60)$& $[60;70)$\\ 
			 \hline
			Số khách hàng nữ & $3$ & $?$ & $?$ & $?$  & $?$ \\
			 \hline
		\end{tabular}
	\end{center}
	 \loigiai{
		\begin{center}
			\begin{tabular}{|c|c|c|c|c|c|}
				 \hline
				Khoảng tuổi &  $[20;30)$ & $[30;40)$ & $[40;50)$ & $[50;60)$& $[60;70)$\\
				 \hline
				Giá trị đại diện& $25$ & $35$  & $45$& $55$& $65$\\
				 \hline
				Độ dài của nhóm & $10$  & $10$ & $10$ & $10$ & $10$ \\
				 \hline
			\end{tabular}
		\end{center}
	} 
\end{ex}

% nguồn: Ok_TuanLuan_Toan11_Lan2
\begin{ex}%[1D5H1-2]%[ma:f6e4ba42]
	Cân nặng của $28$ học sinh nam lớp $11$ được cho  như sau:
	\[\begin{array}{cccccccccccccc}
	 55{,}4&62{,}6&54{,}2&56{,}8&58{,}8&59{,}4&60{,}7&58&59{,}5&63{,}6&61{,}8&52{,}3&63{,}4&57{,}9\\
	 49{,}7&45{,}1&56{,}2&63{,}2&46{,}1&49{,}6&59{,}1&55{,}3&55{,}8&45{,}5&46{,}8&54&49{,}2&52{,}6.
\end{array}	\]
Hãy chia mẫu dữ liệu trên thành  $5$ nhóm, lập bảng tần số ghép nhóm và xác định giá trị đại diện cho mỗi nhóm.
 \loigiai{
	Khoảng  biến thiên của mẫu số liệu trên là $R=63{,}6-45{,}1=18{,}5$.\\
	Độ dài mỗi nhóm $L>\dfrac{R}{k}=\dfrac{18{,}5}{5}=3{,}7$.\\ 
	Ta chọn $L=4$ và chia dữ liệu thành các nhóm $[45;49)$, $[49;53)$, $[53;57)$, $[57;61)$,  $[61;65)$.\\
	Khi đó ta có bảng tần số  ghép nhóm sau:
	\begin{center}
		\begin{tabular}{|c|c|c|c|c|c|}
			 \hline
			Cân nặng & $[45; 49)$ & $[49; 53)$ & $[53; 57)$ & $[57; 61)$ & $[61;  65)$ \\
			 \hline
			Giá trị đại diện & $47$ & $51$ & $55$  & $59$ & $63$ \\
			\hline 
			Số học sinh & $4$  & $5$ & $7$ & $7$ & $5$ \\
			 \hline
		\end{tabular}
	\end{center}
 }
\end{ex}

% nguồn: Ok_TuanLuan_Toan11_Lan2
\begin{ex}%[1D5H1-2]%[ma:19a28bef]
	Cho mẫu số liệu ghép nhóm về thời gian (phút) đi từ nhà  đến nơi làm việc của các nhân viên một công ty như sau:
	\begin{center}
		\begin{tabular}{|c|c|c|c|c|c|c|c|}
			 \hline
			Thời gian & $[15; 20)$ & $[20; 25)$ & $[25; 30)$ & $[30; 35)$ & $[35; 40)$ & $[40; 45)$ & $[45;  50)$ \\
			 \hline
			Số nhân viên & $6$ & $14$ & $25$ & $37$ & $21$ & $13$ & $9$ \\ 
			\hline 
		\end{tabular}
	\end{center}
	 Đọc và giải thích mẫu số liệu này.
	\loigiai{Mẫu số liệu đã cho là mẫu số liệu ghép nhóm gồm $7$ nhóm mô tả về thời gian đi từ nhà đến nơi làm việc của các nhân  viên một công ty:
		\begin{itemize}
			\item Nhóm $1$: Thời gian đi từ $15$ phút đến dưới $20$  phút, có $6$ nhân viên;
			\item Nhóm $2$: Thời gian đi từ $20$ phút đến dưới $25$ phút, có $14$ nhân viên; 
			\item Nhóm $3$: Thời gian đi từ  $25$ phút đến dưới $30$ phút, có $25$ nhân viên;
			 \item Nhóm $4$: Thời gian đi từ $30$ phút đến dưới $35$ phút, có $37$ nhân viên;
			\item Nhóm $5$: Thời gian đi từ $35$ phút đến dưới $40$ phút, có $21$  nhân viên;
			\item  Nhóm $6$: Thời gian đi từ $40$ phút đến dưới $45$ phút, có $13$ nhân viên;
			\item Nhóm  $7$: Thời gian đi từ $45$ phút đến dưới $50$ phút, có $9$ nhân viên.
	\end{itemize}}
\end{ex}

% nguồn: Ok_TuanLuan_Toan11_Lan2
\begin{ex}%[1D5H1-2]%[ma:c6b2792f]
	Cho mẫu số liệu về khối lượng của $30$ củ khoai tây thu hoạch ở một nông trường (dơn vị: gam).\\ 
	\begin{center}
		\begin{tabular}{|c|c|c|c|c|c|}
			\hline  Giá trị&{$[70;80)$}&{$[80;90)$}&{$[90;100)$}&{$[100;110)$}&{$[110;120)$}\\
			\hline Số lượng  củ khoai&$3$&$6$&$12$&$6$&$3$\\
			 \hline
		\end{tabular}
	\end{center}
	 Đọc và giải thích mẫu số liệu này.\\
	 \loigiai{
		Mẫu số liệu đã cho là  mẫu số liệu ghép nhóm.\\
		Có tất cả $5$ nhóm  là
		\begin{enumEX}[-]{1}
			\item từ $70$ gam đến dưới $80$ gam  có $3$ củ.
			\item từ $80$ gam đến dưới $90$ gam có $6$ củ. 
			\item từ $90$ gam đến dưới $100$ gam có $12$ củ. \item từ $100$ gam đến dưới $110$ gam có  $6$ củ.
			\item từ $110$ gam  đến dưới $120$ gam có $3$ củ.
		\end{enumEX}
	} 
\end{ex}

% nguồn: Ok_TuanLuan_Toan11_Lan2
\begin{ex}%[1D5H1-2]%[ma:341463b5]
	Chỉ số BMI (đo bằng $w/h^{2}$, trong đó $w$ là cân nặng đơn vị là kilôgam, $h$ là chiều cao  đơn vị là mét) của các học sinh trong một tổ được cho như sau:
	\[19{,}2 \quad 21{,}1 \quad 16{,}8 \quad 23{,}5 \quad 20{,}6  \quad 25{,}2 \quad 18{,}7 \quad 19{,}1.\]
	Một người có chỉ số BMI nhỏ hơn $18{,}5$ được xem là thiếu cân; từ $18{,}5$ đến dưới $23$ là có cân nặng  lí tưởng so với chiều cao; từ $23$ trở lên là thừa cân. Hãy lập mẫu số liệu ghép nhóm cho mẫu số liệu trên để biểu diễn tình trạng cân nặng so với chiều cao của các học sinh trong tổ.
	\loigiai{ 
		Quan sát mẫu  số liệu đã cho, ta thấy:
		\begin{itemize}
			\item  Có $1$ giá trị BMI của học sinh trong tổ nhỏ hơn $18{,}5$ đó là $16{,}8$;
			\item Có $5$ giá trị BMI của học sinh trong tổ  thuộc nửa khoảng từ $18{,}5$ đến $23$, đó là $19{,}2$; $21{,}1$; $20{,}6$; $ 18{,}7$; $19{,}1$.
			\item Có $2$ giá trị BMI của học sinh trong tổ từ $23$ trở lên, đó là $23{,}5$; $25{,}2$. 
		\end{itemize}
		Từ đó ta lập được mẫu số liệu ghép nhóm về tình trạng cân nặng so với chiều  cao của các học sinh trong tổ như sau:
		\begin{center}
			\begin{tabular}{|c|c|c|c|}
				\hline 
				Chỉ số BMI (w/h$^{2}$)  & Dưới $18{,}5$ & Từ $18{,}5$ đến dưới $23$ & Từ $23$ trở lên \\
				\hline 
				Số học sinh & $1$  & $5$ & $2$ \\
				 \hline
			\end{tabular}
		\end{center}
	 }
\end{ex}

% nguồn: Ok_TuanLuan_Toan11_Lan2
\begin{ex}%[1D5H1-2]%[ma:acf0c64f]
	Một trường Trung học cơ sở  chọn $36$ học sinh nam của khối $9$ để đo chiều cao của các bạn học sinh đó và thu được mẫu số liệu như bảng sau (đơn vị là cm)\\
	\begin{center}
		\begin{tabular}{|l|l|l|l|l|l|l|l|l|l|}
			\hline  $160$&$161$&$161$&$162$&$162$&$162$&$163$&$163$&$163$&$164$\\
			 \hline $164$&$164$&$164$&$165$&$165$&$165$&$165$&$165$&$166$&$166$\\
			\hline $166$&$166$&$167$&$167$&$168$&$168$&$168$&$168$&$169$&$169$\\ 
			\hline  $170$&$171$&$171$&$172$&$172$&$174$&&&&\\
			 \hline
		\end{tabular}
	\end{center}
	Từ mẫu số liệu không ghép nhóm trên, hãy ghép các số liệu thành  5 nhóm theo các nửa khoảng có độ dài bằng nhau.\\
	 \loigiai{
		 Để chuyển mẫu số liệu không ghép nhóm sang mẫu số liệu ghép nhóm ta làm như sau
		\begin{enumEX}[-]{1}
			\item Bước $1$: Chia miền giá trị của mẫu số liệu thành một  số nhóm theo tiêu chí cho trước.
			\item Bước $2$: Đếm số giá trị của mẫu số liệu  thuộc mỗi nhóm và lập bảng thống kê cho mẫu số liệu ghép nhóm.
		\end{enumEX}
		Ta thực hiện bài toán như  sau\\
		Khoảng biến thiên là $174-160=14$.\\ 
		Để  số liệu thành $5$ nhóm theo các nửa khoảng có độ dài bằng nhau, ta chia thành các nhóm có độ dài bằng $3$. \\
		Ta sẽ chọn đầu mút phải của nhóm  cuối cùng là $175$.\\
		Ta được các nhóm  $[160;163)$, $[163;166)$, $[166;169)$, $[169;172)$, $[172;175)$.
	} 
\end{ex}

% nguồn: Ok_TuanLuan_Toan11_Lan2
\begin{ex}%[1D5H1-2]%[ma:1cfe7a6f]
	Số sản phẩm một công nhân làm được trong một ngày được cho như  sau:
	\begin{center}
		\begin{tabular}{llllllllllll}
			18 & 25 & 39 & 12 & 54 & 27 & 46 &  25 & 19 & 8 & 36 & 22 \\
			20 & 19 & 17 & 44 & 5 & 18 & 23 & 28 & 25  & 34 & 46 & 27. \\
		\end{tabular}
	\end{center}
	Hãy chuyển mẫu số liệu sang dạng ghép nhóm với sáu nhóm có độ dài bằng  nhau.
	 \loigiai{
		Giá trị nhỏ nhất của mẫu số liệu là $5$, giá trị lớn nhất của mẫu số liệu là $54$, do đó khoảng biến thiên của mẫu số liệu là $R=54-5=49$. Ta cần chia thành sáu nhóm với độ dài bằng nhau. Để cho thuận tiện, ta  chọn đầu mút trái của nhóm đầu tiên là 3 và đầu mút phải của nhóm cuối cùng bằng $57$ và độ dài của mỗi nhóm bằng $9$ ta được các nhóm là $[3; 12),[12; 21),[21; 30),[30; 39),[39; 48),[48; 57)$. Đếm số giá trị thuộc mỗi nhóm, ta có mẫu số liệu ghép nhóm như sau:
		\begin{center}
			\begin{tabular}{|c|c|c|c|c|c|c|}
				 \hline
				Số sản phẩm & $[3; 12)$ & $[12; 21)$ & $[21; 30)$ & $[30; 39)$ & $[39; 48)$ &  $[48; 57)$ \\
				 \hline
				Số công nhân & $2$ & $8$  & $8$ & $2$ & $4$ & $1$ \\
				\hline 
			\end{tabular}
		\end{center}
	} 
\end{ex}

% nguồn: Ok_TuanLuan_Toan11_Lan2
\begin{ex}%[1D5H1-2]%[ma:44fb0901]
	Từ $1/7/2019$, dựa trên thu nhập bình quân đầu người (kí hiệu là GNIPC, tính theo đô la Mĩ), Ngân hàng Thế giới xác định một nền kinh tế ở mức thu nhập thấp nếu GNIPC nhỏ hơn $1026$, ở mức thu nhập dưới trung bình nếu GNIPC từ $1026$ đến  dưới $3996$, ở mức thu nhập trên trung bình nếu GNIPC từ $3996$ đến dưới $12376$ và ở mức thu nhập cao nếu GNIPC từ $12376$ trở lên (Theo Ngân hàng Thế giới). Thu nhập bình quân đầu người của một số nền kinh tế thuộc khu vực châu Á Thái Bình Dương năm $2021$ được cho như sau
	\begin{center}
		$\begin{array}{|c|c|c|c|c|c|c|c|c|}
		 \hline
		102450&70700&67580&55290&47490&45440&44570&28730&19170 \\ 
		\hline 
		18530&16520&13790&12904&11090&11040&10440&9450&8150\\ 
		 \hline
		7220 &6960&5800&4430&4340&4280&4230&2100&\\ 
		\hline 
	\end{array}$
	\begin{flushright}
		(Theo  statistica.com)
	\end{flushright}
\end{center}
\begin{enumerate}[1)]
	\item Ghép nhóm mẫu số liệu trên theo mức thu  nhập của nền kinh tế.
	\item GNIPC của Việt Nam năm $2021$ là $11040$. Nền kinh  tế Việt Nam được xếp ở mức nào?
\end{enumerate}
\loigiai{ 
	\begin{enumerate}[1)]
		\item  Mẫu số liệu ghép nhóm
		\begin{center}
			\begin{tabular}{|c|c|c|c|}
				 \hline
				GNIPC & $[1026;3996)$ & $[3996;12376)$ & Từ 12376 trở lên \\ 
				 \hline
				Số nền kinh tế  & $1$ & $12$ & $13$ \\
				\hline 
			\end{tabular}
		\end{center}
		\item Nền  kinh tế Việt Nam được xếp ở mức thu nhập trên trung bình.
	\end{enumerate}
} 
\end{ex}

% nguồn: Ok_TuanLuan_Toan11_Lan2
\begin{ex}%[1D5H1-2]%[ma:54d672cb]
	Cho bảng khảo sát về chiều  cao học sinh trong lớp
	\begin{center}
		\begin{tabular}{|c|c|c|c|c|c|}
			\hline Chiều cao  cm&{$[150;160)$}&{$[160;167)$}&{$[167;170)$}&{$[170;175)$}&{$[175;180)$}\\
			\hline Số học  sinh&$12$&$18$&$8$&$3$&$1$\\
			\hline 
		\end{tabular}
	\end{center}
	Xác định độ dài của các nhóm số  liệu?
	\loigiai{ 
		Độ dài của nhóm $[a;b)$  là $b-a$ nên
		\begin{enumEX}[-]{1}
			\item Độ dài của nhóm $[150;160)$ là  $160-150=10$.
			\item Độ dài của nhóm  $[160;167)$ là $167-160=7$.
			\item Độ dài  của nhóm $[167;170)$ là $170-167=3$.
			\item Độ dài của nhóm $[170;175)$ là $175-170=5$. 
			\item Độ  dài của nhóm $[175;180)$ là $180-175=5$.
		\end{enumEX}
	} 
\end{ex}

% nguồn: Ok_TuanLuan_Toan11_Lan2
\begin{ex}%[1D5H1-2]%[ma:a4ae505f]
	Cho bảng khảo sát về tiền điện của một  số gia đình
	\begin{center}
		\begin{tabular}{|c|c|c|c|c|c|}
			\hline\begin{tabular}{c}
			Số tiền  \\(nghìn
			đồng) 
		\end{tabular}&{$[350;400)$}&{$[400;450)$}&{$[450;500)$}&{$[500;550)$}&{$[550;600)$}\\
		\hline Số  hộ gia đình&$12$&$18$&$8$&$3$&$1$\\
		\hline 
	\end{tabular}
\end{center}
Các nhóm số liệu ở trên có độ dài  bằng bao nhiêu?
 \loigiai{
	\begin{enumEX}[-]{1}
		 \item Độ dài của nhóm $[350;400)$ là $400-350=50$.
		\item  Độ dài của nhóm $[400;450)$ là $450-400=50$.
		\item Độ  dài của nhóm $[450;500)$ là $500-450=50$.
		\item Độ dài của nhóm $[500;550)$  là $550-500=50$.
		\item Độ dài của  nhóm $[550;600)$ là $600-550=50$.
	\end{enumEX}
	Vậy các nhóm số liệu ở trên có độ dài bằng  $50$.
 }
\end{ex}

% nguồn: Ok_TuanLuan_Toan11_Lan2
\begin{ex}%[1D5H1-2]%[ma:243d108a]
	Thành tích chạy $50$ m của học sinh lớp $10$A của trường THPT B  được thống kê trong bảng sau
	\begin{center}
		\begin{tabular}{|c|c|c|c|c|c|c|}
			 \hline
			$6{,}3$& $6{,}2$& $6{,}5$&$6{,}8$&  $6{,}9$& $8{,}2$&$8{,}6$\\
			 \hline
			$6{,}6$& $6{,}7$& $7{,}0$& $7{,}1$&  $7{,}2$ &$8{,}3$&$8{,}5$\\
			 \hline
			$7{,}4$ & $7{,}3$& $7{,}2$& $7{,}1$&$7{,}0$&$8{,}4$&$8{,}1$\\ 
			 \hline
			$7{,}1$&$7{,}3$&$7{,}5$&  $7{,}5$& $7{,}6$& $8{,}7$&$7{,}6$ \\
			 \hline
			 $7{,}7$& $7{,}8$& $7{,}5$& $7{,}7$ & $7{,}8$& &\\
			 \hline
		\end{tabular}
	\end{center}
	\begin{enumerate}[1)]
		\item Từ mẫu số liệu không ghép nhóm trên, hãy ghép các số liệu thành $6$ nhóm theo các nửa khoảng  có độ dài bằng nhau.
		\item Trong lớp $10$A, có bao nhiêu học sinh chạy $50$ m mà khoảng thời gian  dưới $8$ giây.
	\end{enumerate}
	 \loigiai{
		\begin{enumerate}[1)]
			 \item Lập bảng phân bố tần số ghép nhóm\\
			Ta chia thành các nhóm  có độ dài bằng $0{,}5$.\\
			Ta sẽ chọn  đầu mút phải của nhóm cuối cùng là $90$. \\
			Ta có  bảng ghép nhóm sau
			\begin{center}
				\begin{tabular}{|c|c|c|c|c|c|c|}
					 \hline
					Thời gian chạy(giây)& $[6{,}0;6{,}5)$&  $[6{,}5;7{,}0)$ & $[7{,}0;7{,}5)$& $[7{,}5;8{,}0)$& $[8{,}0;8{,}5)$& $[8{,}5;9{,}0)$\\
					\hline 
					Số học sinh& $2$ & $5$& $10$  & $9$&$4$& $3$\\
					 \hline
				\end{tabular}
			\end{center}
			\item Trong lớp $10$A, số học sinh chạy  $50$ m mà khoảng thời gian dưới $8$ giây $2+5+10+9=26$ học sinh.
		\end{enumerate}
	 }
\end{ex}

% nguồn: Ok_TuanLuan_Toan11_Lan2
\begin{ex}%[1D5H1-2]%[ma:64ac6031]
	Bảng thống kê sau cho biết thời gian chạy (phút) của  $30$ vận động viên (VĐV) trong một giải chạy Marathon.
	\begin{center}
		\begin{tabular}{|c|c|c|c|c|c|c|c|c|c|c|c|c|}
			 \hline
			Thời gian & $129$ & $130$ & $133$ & $134$  & $135$ & $136$ & $138$ & $141$ & $142$ & $143$ & $144$ & $145$ \\
			 \hline
			SỐ VĐV & $1$ & $2$ & $1$ & $1$ & $1$ &  $2$ & $3$ & $3$ & $4$ & $5$ & $2$ & $5$ \\
			\hline 
		\end{tabular}
	\end{center}
	Hãy chuyển mẫu số liệu trên sang mẫu  số liệu ghép nhóm gồm sáu nhóm có độ dài bằng nhau và bằng $3$.
	\loigiai{Giá trị nhỏ  nhất là $129$, giá trị lớn nhất là $145$ nên khoảng biến thiên là $145-129=16$.\\
		Theo yêu cầu của bài toán ta có tổng độ  dài của sáu nhóm là $18$.\\
		 Để cho đối xứng, ta chọn đầu mút trái của nhóm đầu tiên là $127{,}5$ và đầu mút phải của nhóm cuối cùng là $145{,}5$ ta được các nhóm là $[127{,}5;130{,}5)$, $[130{,}5;133{,}5)$,$\cdots$,$[142{,}5;145{,}5)$.\\
		Đếm số giá trị thuộc  mỗi nhóm, ta có mẫu số liệu ghép nhóm như sau
		\begin{center}
			\begin{tabular}{|c|c|c|c|c|c|c|}
				\hline 
				Thời gian & $[127{,}5;130{,}5)$ & $[130{,}5;133{,}5)$ & $[133{,}5;136{,}5)$ & $[136{,}5;139{,}5)$ &  $[139{,}5;142{,}5)$ & $[142{,}5;145{,}5)$ \\
				 \hline
				Số VĐV  & $3$ & $1$ & $4$ & $3$ & $7$ & $12$ \\
				\hline 
			\end{tabular}
		\end{center}
	 }
\end{ex}

% nguồn: Ok_TuanLuan_Toan11_Lan2
\begin{ex}%[1D5H1-2]%[ma:853ea2be]
	Thời gian ra sân (giờ) của một  số cựu cầu thủ ở giải ngoại hạng Anh qua các thời kì được cho như sau
	\begin{center}
		\begin{tabular}{|c|c|c|c|c|c|c|c|c|c|c|c|c|c|c|c|c|}
			 \hline
			$653$ & $632$ & $609$ & $572$ & $565$ & $535$ & $516$ & $514$ & $508$ & $505$ & $504$ & $504$ & $503$ & $499$ & $496$ & $492$  \\
			 \hline
		\end{tabular}
	\end{center}
	Hãy chuyển mẫu số liệu  trên sang dạng ghép nhóm với bảy nhóm có độ dài bằng nhau.
	 \loigiai{
		Giá  trị lớn nhất là $653$, giá trị bé nhất là $492$. Khoảng biến thiên là $653-492=161$.\\
		Để chia thành $7$ nhóm có độ dài bằng nhau, ta lấy điểm đầu mút trái của nhóm đầu tiên là $492$, điểm đầu mút phải của nhóm  cuối là $653$ với độ dài mỗi nhóm là $23$.\\
		Ta có mẫu số liệu  ghép nhóm sau
		\begin{center}
			\begin{tabular}{|c|c|}
				 \hline
				Thời gian ra sân  (giờ) & Số cầu thủ\\
				 \hline
				$[492;515)$& $9$\\ 
				\hline 
				$[515;538)$&  $2$\\
				\hline 
				$[538;561)$& $0$\\ 
				\hline 
				$[461;584)$&  $2$\\
				 \hline
				$[584;584)$& $0$\\ 
				\hline 
				$[607;630)$& $1$\\ 
				\hline 
				 $[630;653)$& $2$ \\
				 \hline
			\end{tabular}
		\end{center}
	 }
\end{ex}

% nguồn: Ok_TuanLuan_Toan11_Lan2
\begin{ex}%[1D5H1-2]%[ma:33a122fe]
	Cho mẫu số liệu về số tiền điện phải trả của  $50$ gia đình trong một tháng ở một khu phố (đơn vị: nghìn đồng).\\
	\begin{center}
		\begin{tabular}{|c|c|c|c|c|c|c|}
			\hline Giá trị&{$[375;450)$}&{$[450;525)$}&{$[525;600)$}&{$[600;675)$}&{$[675;750)$}&{$[750;825)$}\\ 
			\hline  Số lượng gia đình&$6$&$15$&$10$&$6$&$9$&$4$\\
			 \hline
		\end{tabular}
	\end{center}
	Đọc và giải thích mẫu số  liệu này.\\
	\loigiai{ 
		Mẫu số liệu đã cho là mẫu  số liệu ghép nhóm.\\
		Có tất cả  $6$ nhóm là
		\begin{enumEX}[-]{1}
			\item từ $375$ nghìn đồng đến dưới $450$ nghìn đồng  có $6$ gia đình.
			\item từ  $450$ nghìn đồng đến dưới $525$ nghìn đồng có $15$ gia đình.
			\item từ $525$ nghìn đồng đến dưới $600$ nghìn đồng có  $10$ gia đình.
			\item từ  $600$ nghìn đồng đến dưới $675$ nghìn đồng có $6$ gia đình.
			\item từ $675$ nghìn đồng đến dưới $750$  nghìn đồng có $9$ gia đình.
			\item từ $750$ nghìn đồng đến dưới $825$ nghìn đồng có $4$ gia  đình.
		\end{enumEX}
	 }
\end{ex}

% nguồn: Ok_TuanLuan_Toan11_Lan2
\begin{ex}%[1D5H1-2]%[ma:ec585537]
	Cho bảng khảo sát về cân nặng của học sinh trong lớp như sau \begin{center}
	\begin{tabular}{|c|c|c|c|c|c|c|}
		\hline 
		Cân nặng (kg)& $[40;45)$& $[45;50)$& $[50;55)$& $[55;60)$& $[60;65)$&  $[65;70)$\\
		 \hline
		Số học sinh &  $7$& $10$& $11$& $10$& $3$& $1$\\
		 \hline
	\end{tabular}
\end{center}
\begin{enumerate}[1)]
	\item Mẫu số liệu đã  cho có phải là mẫu số nhiệu ghép nhóm hay không?
	\item Nêu các nhóm và tần số tương  ứng?
\end{enumerate}
\loigiai{ 
	\begin{enumerate}[1)]
		\item Mẫu số liệu đã cho  là mẫu số nhiệu ghép nhóm.
		\item Nêu các  nhóm và tần số tương ứng
		\begin{itemize}
			\item Có $7$ học sinh có cân nặng từ $40$ kg đến dưới $45$ kg. 
			\item Có $10$ học sinh có cân  nặng từ $45$ kg đến dưới $50$ kg.
			\item Có $11$  học sinh có cân nặng từ $50$ kg đến dưới $55$ kg.
			 \item Có $10$ học sinh có cân nặng từ $55$ kg đến dưới $60$ kg.
			 \item Có $3$ học sinh có cân nặng từ $60$ kg đến dưới $65$ kg.
			\item Có $1$ học sinh có cân nặng từ  $65$ kg đến dưới $70$ kg.
		\end{itemize}
	\end{enumerate}
 }
\end{ex}

% nguồn: Ok_TuanLuan_Toan11_Lan2
\begin{ex}%[1D5H1-2]%[ma:6a22bda8]
	Bảng thống kê sau cho biết điện năng tiêu thụ của 30 hộ ở một khu dân cư trong một tháng như sau (đơn vị:  kW):\\
	\begin{center}
		\begin{tabular}{|l|l|l|l|l|l|l|l|l|l|l|l|l|l|l|}
			 \hline $50$&$47$&$30$&$65$&$63$&$70$&$38$&$34$&$48$&$53$&$33$&$39$&$32$&$40$&$50$\\
			\hline $55$&$50$&$61$&$37$&$37$&$43$&$35$&$65$&$60$&$31$&$33$&$41$&$45$&$55$&$59$\\ 
			 \hline
		\end{tabular}
	\end{center}
	Hãy chuyển mẫu số liệu trên sang mẫu số liệu ghép nhóm gồm $8$ nhóm có độ dài bằng  nhau và bằng $5$.\\
	 \loigiai{
		 Giá trị nhỏ nhất là $30$, giá trị lớn nhất là $70$ nên khoảng biến thiên là $70-30=40$.\\
		Tổng độ dài của $8$ nhóm là $40$ nên ta được các nhóm  như sau
		$$[30;35),  [35;40), [40;45), [45;50), [50;55), [55;60), [60;65), [65;70]$$
		Đếm số giá trị  thuộc mỗi nhóm ta có mẫu số liệu ghép nhóm như sau\\
		\begin{center}
			\begin{tabular}{|c|c|c|c|c|c|c|c|c|}
				\hline  Giá trị&{$[30;35)$}&{$[35;40)$}&{$[40;45)$}&{$[45;50)$}&{$[50;55)$}&{$[55;60)$}&{$[60;65)$}&{$[65;70]$}\\
				\hline Số  lượng&$6$&$5$&$3$&$3$&$4$&$3$&$3$&$3$\\
				 \hline
			\end{tabular}
		\end{center}
	} 
\end{ex}

% nguồn: Ok_TuanLuan_Toan11_Lan2
\begin{ex}%[1D5H1-2]%[ma:60ef972c]
	Cho mẫu số liệu về thời gian (phút) đi từ nhà đến nơi làm việc của các nhân viên của một công ty  như sau
	\begin{center}
		\begin{tabular}{|c|c|c|c|c|c|c|c|}
			\hline  Thời gian&{$[15;20)$}&{$[20;25)$}&{$[25;30)$}&{$[30;35)$}&{$[35;40)$}&{$[40;45)$}&{$[45;50)$}\\
			\hline Số nhân viên&$6$&$14$&$25$&$37$&$21$&$13$&$9$\\ 
			 \hline
		\end{tabular}
	\end{center}
	\begin{enumEX}[a)]{1}
		\item Mẫu số  liệu đã cho có phải là mẫu số nhiệu ghép nhóm hay không?
		\item Nêu các nhóm và tần số tương ứng?  Tổng số nhân viên công ty là bao nhiêu?
	\end{enumEX}
	 \loigiai{
		\begin{enumEX}[a)]{1}
			\item Mẫu  số liệu đã cho là mẫu số liệu ghép nhóm.
			\item Có $7$ nhóm là 
			\begin{enumEX}[-]{1}
				\item Từ $15$ đến dưới $20$ phút. 
				\item Từ $20$ đến dưới $25$  phút.
				\item Từ $25$ đến dưới $30$ phút. 
				\item  Từ $30$ đến dưới $35$ phút.
				\item Từ $35$  đến dưới $40$ phút.
				\item Từ $40$ đến dưới $45$ phút. 
				\item Từ $45$ đến dưới $50$  phút.
			\end{enumEX}
			Các tần  số
			\begin{enumEX}[-]{1}
				\item Có $6$ nhân viên có thời gian di chuyển từ nhà đến nơi làm việc từ  $15$ đến dưới $20$ phút
				\item Có $14$ nhân viên có thời  gian di chuyển từ nhà đến nơi làm việc từ $20$ đến dưới $25$ phút.
				\item Có $25$ nhân viên có thời gian di chuyển từ  nhà đến nơi làm việc từ $25$ đến dưới $30$ phút.
				\item Có $37$ nhân viên có thời gian di chuyển từ nhà đến nơi làm việc từ $30$  đến dưới $35$ phút.
				\item Có $21$  nhân viên có thời gian di chuyển từ nhà đến nơi làm việc từ $35$ đến dưới $40$ phút.
				\item Có  $13$ nhân viên có thời gian di chuyển từ nhà đến nơi làm việc từ $40$ đến dưới $45$ phút.
				\item Có $9$ nhân viên có thời gian di chuyển từ nhà đến nơi làm  việc từ $45$ đến dưới $50$ phút.
			\end{enumEX}
			Tổng số nhân viên của công ty là $6+14+25+37+21+13+9=125$ người. 
		\end{enumEX}
	} 
\end{ex}

% nguồn: Ok_TuanLuan_Toan11_Lan2
\begin{ex}%[1D5H1-2]%[ma:a3d7ec44]
	Số tiền mà sinh viên chi cho thanh toán cước điện thoại  trong tháng.
	\begin{center}
		\begin{tabular}{|c|c|c|c|c|c|}
			 \hline
			Khoảng  tuổi & $[0;50)$ & $[50;100)$ & $[100;150)$ & $[150;200)$& $[200;250)$\\
			 \hline
			Số khách hàng nữ & $5$ &  $12$ & $23$ & $17$ & $3$ \\
			\hline 
		\end{tabular}
	\end{center}
	Có bao nhiêu sinh viên  được khảo sát trong bảng số liệu trên?
	 \loigiai{
		Số sinh viên được  khảo sát trong bảng số liệu trên $5+12+23+17+3=60$.
	 }
\end{ex}

% nguồn: Ok_TuanLuan_Toan11_Lan2
\begin{ex}%[1D5H1-2]%[ma:0ebfa392]
	Cho bảng khảo sát  về cân nặng của học sinh trong lớp như sau
	\begin{center}
		\begin{tabular}{|c|c|c|c|c|c|c|}
			\hline 
			Cân nặng (kg)& $[40;45)$& $[45;50)$& $[50;55)$& $[55;60)$& $[60;65)$  & $[65;70)$ \\
			 \hline
			Số học sinh& $7$ & $10$ & $11$ & $10$ & $3$  & $1$\\
			\hline 
		\end{tabular}
	\end{center}
	\textbf{}Độ dài của các nhóm là  bao nhiêu?
	\loigiai{ 
		Độ dài của  các nhóm là $5$.
	} 
\end{ex}

% nguồn: Ok_TuanLuan_Toan11_Lan2
\begin{ex}%[1D5H1-2]%[ma:0428dbf9]
	Trong các mẫu  số liệu sau, mẫu nào là mẫu số liệu ghép nhóm? Đọc và giải thích mẫu số liệu ghép nhóm đó
	\begin{enumEX}[a)]{1}
		\item Số tiền mà sinh viên chi cho  thanh toán cước điện thoại trong tháng:
		\begin{center}
			\begin{tabular}{|c|c|c|c|c|c|}
				\hline Số tiền  (nghìn đồng)&{$[0;50)$}&{$[50;100)$}&{$[100;150)$}&{$[150;200)$}&{$[200;250)$}\\
				\hline Số  sinh viên&$5$&$12$&$23$&$17$&$3$\\
				 \hline
			\end{tabular}
		\end{center}
		\item Thống kê nhiệt độ tại một địa điểm  trong $40$ ngày, ta có bảng số liệu sau
		\begin{center}
			\begin{tabular}{|c|c|c|c|c|}
				 \hline Nhiệt độ ($^\circ$)C&{$[19;22)$}&{$[22;25)$}&{$[25;28)$}&{$[28;31)$}\\
				 \hline Số ngày&$7$&$15$&$12$&$6$\\
				 \hline
			\end{tabular}
		\end{center}
	\end{enumEX}
	\loigiai{ 
		\begin{enumEX}[a)]{1}
			\item Mẫu số liệu này mô tả về số tiền mà sinh viên chi cho thanh toán cước điện thoại trong tháng, gồm có $5$  nhóm. Cụ thể:
			\begin{enumEX}[-]{1}
				 \item Nhóm thanh toán với số tiền từ $0$ đến dưới $50$ nghìn đồng, có $5$ sinh viên
				\item Nhóm thanh toán với số tiền từ $50$ nghìn đến dưới $100$ nghìn đồng,  có $12$ sinh viên
				\item Nhóm thanh toán với số tiền từ $100$ nghìn đến dưới $150$ nghìn đồng, có $23$  sinh viên
				\item Nhóm thanh toán với số tiền từ $150$ nghìn đến dưới $200$ nghìn đồng,  có $17$ sinh viên
				\item Nhóm thanh toán với số tiền từ $200$  nghìn đến dưới $250$ nghìn đồng, có $3$ sinh viên
			\end{enumEX}
			\item  Mẫu số liệu này mô tả về nhiệt độ tại một địa điểm trong $40$ ngày, gồm $4$ nhóm nhiệt độ
			\begin{enumEX}[-]{1}
				\item Từ $19^\circ$C  đến dưới $22^\circ$C.
				\item  Từ $22^\circ$C đến dưới $25^\circ$C.
				 \item Từ $25^\circ$C đến dưới $28^\circ$C.
				\item Từ  $28^\circ$C đến dưới $31^\circ$C.
			\end{enumEX}
			Cụ thể:\\ 
			\begin{enumEX}[-]{1}
				\item Có $7$ ngày  có nhiệt độ từ $19^\circ$C đến dưới $22^\circ$C.
				\item Có  $15$ ngày có nhiệt độ từ $22^\circ$C đến dưới $25^\circ$C.
				\item Có $12$ ngày có nhiệt  độ từ từ $25^\circ$C đến dưới $28^\circ$C.
				\item Có $6$ ngày có nhiệt độ từ $28^\circ$C  đến dưới $31^\circ$C.
			\end{enumEX}
		\end{enumEX}
	} 
\end{ex}

% nguồn: Ok_TuanLuan_Toan11_Lan2
\begin{ex}%[1D5H1-2]%[1D5H1-2]%[ma:b7fce0aa]
	Thời gian hoàn thành bài kiểm tra Toán $45$ phút của các bạn trong  lớp được cho như sau:
	\begin{center}
		\begin{tabular}{|c|c|c|c|c|}
			\hline 
			Thời gian (phút) & $[25; 30)$ & $[30;  35)$ & $[35; 40)$ & $[40; 45]$ \\
			 \hline
			Só học  sinh & $2$ & $7$ & $10$ & $25$ \\
			 \hline
		\end{tabular}
	\end{center}
	\begin{enumEX}{2}
		\item! Nêu các nhóm số liệu và tần  số tương ứng.
		\item! Có bao nhiêu học sinh  hoàn thành bài kiểm tra trước khi hết giờ trên $5$ phút?
	\end{enumEX}
	\loigiai{ 
		\begin{enumEX}{2}
			\item! Các nhóm số liệu $[25; 30)$, $[30; 35)$, $[35; 40)$, $[40; 45]$ với  tần số tương ứng là $2$, $7$, $10$, $25$.
			\item! Số học sinh hoàn thành bài  kiểm tra trước khi hết giờ ít nhất $5$ phút là $2+7+10=19$.
		\end{enumEX}
	 }
\end{ex}

% nguồn: Ok_TuanLuan_Toan11_Lan2
\begin{ex}%[1D5H1-2]%[ma:ad94effe]
	Thời gian ra sân (giờ) của một số cựu cầu thủ ở giải ngoại hạng Anh qua các thời kì được cho như sau: \[\begin{array}{cccccccc}
	653 & 632 & 609 & 572 & 565  & 535 & 516 & 514\\
	508 & 505&504&504&503&499 & 496 & 492. 
\end{array}\]
Hãy chuyển mẫu số liệu trên sang dạng ghép nhóm với bảy  nhóm có độ dài bằng nhau.
\loigiai{ 
	Giá trị nhỏ nhất của mẫu số liệu là $492$, giá trị lớn nhất của mẫu số liệu là $653$, do đó khoảng biến  thiên của mẫu số liệu là $R=653-492=161$.\\
	Ta cần chia thành bảy nhóm có độ dài bằng nhau. Để cho thuận tiện,  ta chọn đầu mút trái của nhóm đầu tiên là $485$ và đầu mút phải của nhóm cuối cùng bằng $660$ và độ dài của mỗi nhóm bằng $25$ ta được các nhóm là \[[485;510),[510;535),[535;560),[560;585),[585;610),[610;635),[635;660].\]
	Đếm số giá trị  thuộc mỗi nhóm, ta có mẫu số liệu ghép nhóm như sau:
	\begin{center}
		\begin{tabular}{|c|c|c|c|c|c|c|c|}
			 \hline
			\begin{tabular}{l}
				Thời  \\
				gian  \\
			\end{tabular}& $[485; 510)$ & $[510; 535)$ & $[535; 560)$ & $[560; 585)$ & $[585; 610)$ & $[610; 635)$ & $[635; 660]$ \\
			\hline 
			\begin{tabular}{c}
				Số \\ 
				 cầu \\
				thủ  \\
			\end{tabular}& $8$ & $2$ & $1$ & $2$ & $1$ & $1$ & $1$ \\
			 \hline
		\end{tabular}
	\end{center}
 }
\end{ex}

% nguồn: Ok_TuanLuan_Toan11_Lan2
\begin{ex}%[1D5H1-2]%[ma:ec616bcb]
	Thống kê chỉ số chất lượng không khí (AQI) tại một địa điểm vào các ngày  trong tháng $6/2022$ được cho trong bảng sau:
	\begin{center}
		\begin{tabular}{|c|c|c|c|c|c|}
			 \hline
			Chỉ số  $AQI$ & $[0; 50)$ & $[50; 100)$ & $[100; 150)$ & $[150; 200)$ & Trên $200$ \\
			\hline 
			Số ngày & $5$ &  $11$ & $7$ & $4$ & $3$ \\
			\hline 
		\end{tabular}
	\end{center}
	\begin{enumEX}{2}
		\item! Đọc và giải thích mẫu số  liệu ghép nhóm.
		\item! Chất lượng không khí được xem là tốt nếu $AQI$  nhỏ hơn $50$, là trung bình nếu $AQI$ từ $50$ đến dưới $100$. Trong tháng $6 / 2022$ tại địa điểm này có bao nhiêu ngày chất lượng không khí dưới mức trung bình?
	\end{enumEX}
	 \loigiai{
		\begin{enumerate}[a)]
			\item Trong tháng $6 / 2022$ có $5$ ngày chỉ số $A Q I$ dưới $50; 11$ ngày chỉ số $A Q I$ từ $50$ đến dưới $100$; $7$ ngày chỉ số $A Q I$ từ $100$ đến dưới $150; 4$ ngày chỉ số $A Q I$ từ $150$ đến dưới $200$;  $3$ ngày chỉ số $A Q I$ trên $200$.
			\item Số ngày chất lượng không khí dưới mức trung bình  là $7+4+3=14$.
		\end{enumerate}
	} 
\end{ex}

% nguồn: Ok_TuanLuan_Toan11_Lan2
\begin{ex}%[1D5H1-2]%[ma:b6ead170]
	Cân nặng (kg) của $35$ người  trưởng thành tại một khu dân cư được cho như sau:
	\begin{center}
		\begin{tabular}{ccccccccccccc}
			$43$ & $51$ & $47$ & $62$ & $48$ & $40$ & $50$ & $62$ & $53$ & $56$ &  $40$ & $48$ & $56$ \\
			$53$ &  $50$ & $42$ & $55$ & $52$ & $48$ & $46$ & $45$ & $54$ & $52$ & $50$ & $47$ & $44$ \\
			$54$ & $55$ &  $60$ & $63$ & $58$ & $55$ & $60$ & $58$ & $53$. & & & & \\
		\end{tabular}
	\end{center}
	Chuyển mẫu số liệu trên thành dạng ghép nhóm, các  nhóm có độ dài bằng nhau, trong đó có nhóm $[40; 45)$.
	 \loigiai{
		Độ dài  của mỗi nhóm là $45-40=5$.\\
		Giá trị nhỏ nhất là $40$, giá trị lớn nhất là  $63$ , do đó, khoảng biến thiên của mẫu số liệu là $63-40=23$.\\
		Để cho  thuận tiện, ta chọn đầu mút trái của nhóm đầu tiên là $40$ và đầu mút phải của nhóm cuối cùng là $65$ , ta được các nhóm là $[40; 45),[45; 50),[50; 55),[55; 60),[60; 65)$.\\
		Đếm số giá trị thuộc mỗi nhóm, ta  có mẫu số liệu ghép nhóm như sau:
		\begin{center}
			\begin{tabular}{|c|c|c|c|c|c|}
				 \hline
				Cân nặng (kg) & $[40;  45)$ & $[45; 50)$ & $[50; 55)$ & $[55; 60)$ & $[60; 65)$ \\
				\hline 
				Số người & $5$ &  $7$ & $11$ & $7$ & $5$ \\
				 \hline
			\end{tabular}
		\end{center}
	 }
\end{ex}

% nguồn: Ok_TuanLuan_Toan11_Lan2
\begin{ex}%[1D5H1-2]%[ma:7a15d63d]
	$100$ người thực hiện bài trắc nghiệm để đo chỉ số IQ, kết quả thu được  như sau
	\begin{center}
		\begin{tabular}{|l|c|c|c|c|c|c|}
			\hline  Chỉ số IQ&Dưới $70$&{$[70;85)$}&{$[85;115)$}&{$[115;130)$}&{$[130;145)$}&Từ $145$ trở lên\\
			\hline Số  người&$2$&$15$&$45$&$20$&$15$&$3$\\
			\hline 
		\end{tabular}
	\end{center}
	\begin{enumEX}[a)]{1}
		\item Nêu các nhóm số liệu và tần số tương  ứng.
		\item Người có chỉ số $IQ$ từ $85$ đến dưới $115$ là ở mức trung bình. Xác định tỉ lệ người có $IQ$ cao hơn mức trung bình.\\ 
	\end{enumEX}
	 \loigiai{
		\begin{enumEX}[a)]{1}
			\item Các nhóm số liệu gồm dưới  $70$; $[70;85)$; $[85;115)$; $[115;130$); $[130;145)$.\\
			Từ $145$ trở lên với tần số tương ứng là $2$, $15$, $45$,  $20$, $15$, $3$.
			\item Số người có chỉ số IQ cao hơn mức trung bình  là $20+15+3=38$. \\
			Vậy tỉ lệ người  có chỉ số IQ cao hơn mức trung bình là $38/100=38\%$.
		\end{enumEX}
	} 
\end{ex}

% nguồn: Ok_TuanLuan_Toan11_Lan2
\begin{ex}%[1D5H1-2]%[ma:7963809f]
	Số sản phẩm một công nhân làm được trong  một ngày được cho như sau
	\begin{center}
		\begin{tabular}{|c|c|c|c|c|c|c|c|c|c|c|c|c|}
			 \hline
			$18$ & $25$ & $39$ & $12$  & $54$ & $27$ & $46$ & $25$ & $19$ & $8$ & $36$ & $22$& \\
			\hline 
			$20$ & $19$ & $17$  & $44$ & $5$ & $18$ & $23$ & $28$ & $25$ & $34$ & $46$ & $27$ & $16$ \\
			 \hline
		\end{tabular}
	\end{center}
	Hãy chuyển mẫu số liệu sang dạng ghép nhóm với sáu  nhóm có độ dài bằng nhau.
	\loigiai{ 
		Giá trị  nhỏ nhất là $5$, giá trị lớn nhất là $54$. Do đó khoảng biến thiên là $54-5=49$.\\
		Để chia thành $6$ nhóm với độ dài bằng nhau ta lấy điểm đầu mút phải trái của nhóm đầu tiên là $3$ và đầu mút phải của  nhóm cuối cùng là $57$ với độ dài mỗi nhóm là $9$.\\
		Ta được  mẫu số liệu ghép nhóm như sau
		\begin{center}
			\begin{tabular}{|c|c|c|c|c|c|c|}
				\hline 
				Số sản phẩm  & $[3;12)$ & $[12;21)$ & $[21;30)$ & $[30;39)$ & $[39;48)$ & $[48;57)$ \\
				\hline 
				Số công nhân & $2$ & $8$ & $8$ & $2$ & $4$ &  $1$ \\
				 \hline
			\end{tabular}
		\end{center}
	} 
\end{ex}

% nguồn: Ok_TuanLuan_Toan11_Lan2
\begin{ex}%[1D5H1-2]%[ma:6ff18d59]
	Cân nặng của 28 học sinh nam lớp $11$ được cho như sau 
	\begin{center}
		$\begin{array}{llllllllllllll}
		{55{,}4} &  {62{,}6} & {54{,}2} & {56{,}8} & {58{,}8} & {59{,}4} & {60{,}7} & {58} & {59{,}5} & {63{,}6} & {61{,}8} & {52{,}3} & {63{,}4} & {57{,}9} \\
		{49{,}7} & {45{,}1} & {56{,}2} & {63{,}2} & {46{,}1} & {49{,}6} & {59{,}1} & {55{,}3} & {55{,}8} & {45{,}5} &  {46{,}8} & {54} & {49{,}2} & {52{,}6}
	\end{array}$
\end{center}
Hãy chia mẫu dữ  liệu trên thành $5$ nhóm, lập bảng tần số ghép nhóm và xác định giá trị đại diện cho mỗi nhóm.
 \loigiai{
	 Khoảng biến thiên của mẫu số liệu trên là $R=63{,}6-45{,}1=18{,}5$.\\
	Độ dài mỗi nhóm  $L> \dfrac{R}{k}=\dfrac{18{,}5}{5}=3{,}7$.\\
	Ta chọn $L=4$ và chia  dữ liệu thành các nhóm $\left[45;49\right),\left[49;53\right),\left[53;57\right),\left[57;61\right),\left[61;65\right)$.\\
	Khi đó ta có bảng  tần số ghép nhóm sau
	\begin{center}
		\begin{tabular}{|c|c|c|c|c|c|}
			\hline 
			Cân nặng & $[45;49)$ &  $[49;53)$ & $[53;57)$ & $[57;61)$& $[61;65)$\\
			 \hline
			 Giá trị đại diện& $47$ & $51$ & $55$ & $59$ & $63$\\
			\hline 
			Số học sinh& $4$ & $5$  & $7$ & $7$ & $5$ \\
			\hline 
		\end{tabular}
	\end{center}
 }
\end{ex}

% nguồn: Ok_TuanLuan_Toan11_Lan2
\begin{ex}%[1D5H1-2]%[ma:4b814ae0]
	Bảng thống kê sau cho biết điện năng tiêu thụ của $30$ hộ ở một khu dân cư  trong một tháng như sau (đơn vị: kW):
	\begin{center}
		\begin{tabular}{|c|c|c|c|c|c|c|c|c|c|}
			\hline 
			$50$ & $47$ & $30$ & $65$ & $63$ &  $70$ & $38$ & $34$ & $48$ & $53$ \\
			 \hline
			$40$ & $50$ & $55$  & $50$ & $61$ & $37$ & $37$ & $43$ & $35$ & $65$ \\
			\hline 
			$41$ & $45$ & $55$ & $59$ & $39$ &  $32$& $31$ & $33$ & $33$ & $60$ \\
			\hline 
		\end{tabular}
	\end{center}
	Hãy chuyển mẫu số liệu trên sang mẫu số  liệu ghép nhóm gồm $8$ nhóm có độ dài bằng nhau và bằng $5$.
	 \loigiai{
		Giá trị nhỏ nhất là $30$, giá trị lớn nhất  là $70$ nên khoảng biến thiên là $70-30=40$.\\
		Tổng độ dài của $8$ nhóm là $40$ nên ta  được các nhóm như sau
		$[30;35)$,  $[35;40)$, $[40;45)$, $[45;50)$, $[50;55)$, $[55;60)$, $[60;65)$, $[65;70]$.\\
		Đếm số giá trị thuộc mỗi nhóm ta có mẫu số liệu ghép nhóm như sau 
		\begin{center}
			\begin{tabular}{|c|c|c|c|c|c|c|c|c|}
				 \hline
				Giá  trị & $[30;35)$ & $[35;40)$ & $[40;45)$ & $[45:50)$ & $[50;55)$ & $[55;60)$ & $[60;65)$ & $[65;70]$ \\
				 \hline
				Số lượng & $6$ & $5$ & $3$ & $3$  & $4$ & $3$ & $3$ & $3$ \\
				 \hline
			\end{tabular}
		\end{center}
	} 
\end{ex}

% nguồn: Ok_TuanLuan_Toan11_Lan2
\begin{ex}%[1D5H1-2]%[ma:0591b6cf]
	Cho mẫu  số liệu về số tiền điện phải trả của $50$ gia đình trong một tháng ở một khu phố (đơn vị: nghìn đồng).
	\begin{center}
		\begin{tabular}{|c|c|c|c|c|c|c|}
			\hline 
			Giá trị & $[375;450)$ &  $[450;525)$ & $[525;600)$ & $[600;675)$ & $[675;750)$ & $[750;825)$ \\
			 \hline
			Số lượng gia đình & $6$ & $15$ & $10$  & $6$ & $9$ & $4$ \\
			\hline 
		\end{tabular}
	\end{center}
	Đọc và giải thích  mẫu số liệu này.
	\loigiai{ 
		Mẫu số liệu đã cho là mẫu  số liệu ghép nhóm.\\
		Có tất  cả $6$ nhóm:
		\begin{itemize}
			\item Từ $375$  nghìn đồng đến dưới $450$ nghìn đồng có $6$ gia đình;
			\item Từ $450$ nghìn đồng  đến dưới $525$ nghìn đồng có $15$ gia đình;
			\item Từ $525$ nghìn đồng đến dưới  $600$ nghìn đồng có $10$ gia đình;
			\item Từ $600$ nghìn đồng  đến dưới $675$ nghìn đồng có $6$ gia đình;
			\item Từ $675$ nghìn đồng đến  dưới $750$ nghìn đồng có $9$ gia đình;
			\item Từ $750$ nghìn đồng đến dưới $825$  nghìn đồng có $4$ gia đình.
		\end{itemize}
	} 
\end{ex}

% nguồn: Ok_TuanLuan_Toan11_Lan2
\begin{ex}%[1D5H1-2]%[ma:3db76b2d]
	Mẫu số liệu dưới đây ghi lại tốc độ của $42$ ô tô khi đi qua một trạm đo tốc độ (đơn vị: km/h). Lập bảng tần số ghép nhóm bao gồm cả tần số tích luỹ cho mẫu số liệu trên có sáu nóm ứng với sáu  nửa khoảng
	\begin{center}
		$\begin{array}{cccccc}
		[40;45)&[45;50)&[50;55)&[55;60)&[60;65)&[65;70)\\ 
	\end{array}$
\end{center}
 \loigiai{
	Bảng tần số ghép nhóm bao gồm cả tần  số tích luỹ như ở Bảng 5
	\begin{center}
		\begin{tabular}{|c|c|c|}
			\hline 
			 Nhóm &Tần số & Tần số tích luỹ\\
			 \hline
			$[40;45)$ & $5$  & $5$ \\
			\hline 
			$[45;50)$ & $10$ & $15$  \\
			\hline 
			$[50;55)$& $7$ & $22$  \\
			 \hline
			$[55;60)$ &  $9$ & $31$ \\
			 \hline
			$[60;65)$&  $7$ & $38$ \\
			 \hline
			$[65;70)$& $4$ & $42$ \\ 
			\hline 
			Tổng& $n=42$ & \\ 
			 \hline
		\end{tabular}\\
		Bảng  5
	\end{center}
} 
\end{ex}

% nguồn: Ok_TuanLuan_Toan11_Lan2
\begin{ex}%[1D5H1-2]%[ma:f6d362e0]
	Một trường Trung học cơ sở chọn $36$ học sinh nam của khối $9$ để đo chiều cao của các bạn học sinh đó và thu được mẫu số liệu như bảng sau (đơn vị là $cm$) \begin{center}
	\begin{tabular}{|c|c|c|c|c|c|c|c|c|c|}
		\hline 
		$160$  & $161$ & $161$ & $162$ & $162$ & $162$ & $163$ & $163$ & $163$ & $164$ \\
		 \hline
		$164$ & $164$ & $164$ & $165$ & $165$  & $165$ & $165$ & $165$ & $166$ & $166$\\
		\hline 
		$166$  & $166$ & $167$ & $167$ & $168$ & $168$ & $168$ & $168$ & $169$ & $169$ \\
		\hline 
		$170$& $171$ & $171$ &  $172$ & $172$ & $174$ & & & & \\
		\hline 
	\end{tabular}
\end{center}
Tìm khoảng biến thiên của mẫu  số liệu trên.
 \loigiai{
	\begin{itemize}
		\item  Giá trị nhỏ nhất là $160$.
		\item Giá trị lớn nhất là  $174$.
	\end{itemize}
	Nên  khoảng biến thiên là $174-160=14$.
 }
\end{ex}

% nguồn: Ok_TuanLuan_Toan11_Lan2
\begin{ex}%[1D5H1-2]%[ma:e13817f0]
	Một cửa hàng đã thống kê số ba lô bán được mỗi ngày trong tháng 9 với kết quả cho như  sau
	\begin{center}
		$\begin{array}{lllllllllllllll}
		{12} & {29} & {12} & {19} & {15} & {21} & {19} & {29} & {28} & {12} & {15} & {25} & {16} & {20} & {29} \\ {21} & {12} & {24} &  {14} & {10} & {12} & {10} & {23} & {27} & {28} & {18} & {16} & {10} & {20} & {21}
	\end{array}$
\end{center}
Hãy chia mẫu số liệu trên thành  $5$ nhóm, lập bảng tần số ghép nhóm, hiệu chỉnh bảng tần số ghép nhóm và xác định giá trị đại diện cho mỗi nhóm.
 \loigiai{
	Khoảng biến thiên của mẫu  số liệu trên là $R=29-10=19$.\\
	Độ dài mỗi nhóm $L>  \dfrac{R}{k}=\dfrac{19}{5}=3{,}8$.\\
	Ta chọn $L=4$ và  chia dữ liệu thành các nhóm
	 \[\left[10;14\right), \; \left[14;18\right), \; \left[18;22\right), \; \left[22;26\right), \; \left[26;30\right)
	\] 
	Khi  đó ta có bảng tần số ghép nhóm sau
	\begin{center}
		\begin{tabular}{|c|c|c|c|c|c|}
			\hline 
			Số ba lô& $[10;14)$& $[14;18)$ & $[18;22)$&  $[22;26)$ & $[26;30)$\\
			 \hline
			Giá trị đại diện & $12$ & $16$& $20$&  $24$ &$28$ \\
			 \hline
			Số ngày & $8$ & $5$&  $8$& $3$& $6$\\
			\hline 
		\end{tabular}
	\end{center}
} 
\end{ex}

% nguồn: Ok_TuanLuan_Toan11_Lan2
\begin{ex}%[1D5H1-2]%[ma:d48aa8b1]
	Số  sản phẩm một công nhân làm được trong một ngày được cho như sau\\
	\begin{center}
		\begin{tabular}{lllllllllllll}
			$18$&$25$&$39$&$12$&$54$&$27$&$46$&$25$&$19$&$8$&$36$&$22$&\\ 
			$20$&$19$&$17$&$44$&$5$&$18$&$23$&$28$&$25$&$34$&$46$&$27$&$16$. 
		\end{tabular}
	\end{center}
	Hãy chuyển mẫu số liệu sang dạng ghép nhóm với sáu nhóm có  độ dài bằng nhau.
	\loigiai{ 
		Giá trị lớn nhất của mẫu số  liệu là $54$, giá trị nhỏ nhất của mẫu số liệu là $5$. \\
		Khoảng  biến thiên là $54-5=49$.\\
		Ta chia thành $6$ nhóm có độ dài bằng nhau. \\ 
		 Ví dụ chọn độ dài mỗi nhóm là $8{,}5$, đầu mút trái là $4$, khi đó ta có mẫu số liệu ghép nhóm như sau
		\begin{center}
			\begin{tabular}{|c|c|c|c|c|c|c|}
				 \hline Số sản phẩm&{$[4;12,5)$}&{$[12,5;21)$}&{$[21;29,5)$}&{$[29,5;38)$}&{$[38;46,5)$}&{$[46,5;55)$}\\
				\hline Số công  nhân&$3$&$7$&$8$&$2$&$4$&$1$\\
				 \hline
			\end{tabular}
		\end{center}
	} 
\end{ex}

% nguồn: Ok_TuanLuan_Toan11_Lan2
\begin{ex}%[1D5H1-2]%[ma:561c4659]
	Theo Tổ chức Y tế Thế giới (WHO), thiếu  máu là tình trạng giảm lượng huyết sắc tố (Hb) dẫn tới sự thiếu cung cấp oxygen cho các mô trong cơ thể. Đối với nam giới trên $15$ tuổi, chỉ số Hb (đơn vị tính là g/l) lớn hơn hoặc bằng $130$ được xem là không bị thiếu máu, từ $110$ đến dưới $130$ là thiếu máu mức nhẹ, từ $80$ đến dưới $110$ là thiếu máu mức vừa, dưới $80$ là mức nặng. Đo chỉ số Hb của một số học sinh nam lớp $12$ cho kết quả như sau:
	\begin{center}
		\begin{tabular}{cccccccccccc}
			132&135&137&131&129&125&140&147&138&137&128&112\\ 
			127&129&125&98&139&138&139&141&140&105&136&133\\ 
			137&138&108&133&136&141&144&134&136&137&142. 
		\end{tabular}
	\end{center}
	Ghép nhóm  cho mẫu số liệu này theo mức độ thiếu máu.
	\loigiai{ 
		Có $25$ học sinh có chỉ số Hb từ $130$ trở lên (không bị thiếu máu), $7$ học sinh có chỉ số Hb từ $110$ đến dưới $130$ (thiếu máu mức nhẹ), $3$ học sinh có chỉ  số Hb từ $80$ đến dưới $110$ (thiếu máu mức trung bình). Ta có mẫu số liệu ghép nhóm:
		\begin{center}
			\begin{tabular}{|c|c|c|c|}
				 \hline
				 Chỉ số Hb(g/l)& Từ $130$ trở lên & $[110; 130)$ & $[80; 110)$ \\
				\hline 
				Số học sinh & $25$ & $7$  & $3$ \\
				\hline 
			\end{tabular}
		\end{center}
	} 
\end{ex}

% nguồn: Ok_TuanLuan_Toan11_Lan2
\begin{ex}%[1D5H1-2]%[ma:743264bd]
	Theo Tổ chức Y tế Thế giới (WHO), thiếu máu là tình trạng giảm lượng huyết sắc tố (Hb) dẫn tới sự thiếu cung cấp oxygen cho các mô trong cơ thể. Đối với nam giới trên $15$ tuổi, chỉ số Hb (đơn vị tính là g/l) lớn hơn hoặc bằng $130$ được xem là không bị thiếu máu, từ $110$ đến dưới $130$ là thiếu máu mức nhẹ, từ $80$  đến dưới $110$ là thiếu máu mức vừa, dưới $80$ là mức nặng. Đo chỉ số Hb của một số học sinh nam lớp $12$ cho kết quả như sau
	$$\begin{aligned}
	&132,135,137,131,129,125,140,147,138,137,128,112,\\ 
	 &127,129,125,98,139,138,139,141,140,105,136,133,\\
	 &137,138,108,133,136,141,144,134,136,137,142.
\end{aligned}$$
Ghép nhóm cho mẫu số liệu này  theo mức độ thiếu máu.
\loigiai{ 
	Có $25$ học sinh có chỉ số Hb từ $130$ trở lên (không bị thiếu máu), $7$ học sinh có chỉ số Hb từ $110$ đến dưới $130$ (thiếu máu mức nhẹ), $3$ học sinh có chỉ số Hb từ $80$ đến dưới $110$ (thiếu máu mức trung bình).\\  Ta có mẫu số liệu ghép nhóm:
	\begin{center}
		\begin{tabular}{|l|c|c|c|}
			\hline  Chỉ số Hb (g/l)&Từ $130$ trở lên&{$[110;130)$}&{$[80;110)$}\\
			 \hline Số học sinh&$25$&$7$&$3$\\
			 \hline
		\end{tabular}
	\end{center}
 }
\end{ex}

% nguồn: Ok_TuanLuan_Toan11_Lan2
\begin{ex}%[1D5H1-2]%[ma:21171459]
	Cho bảng số  liệu thống kê sau
	Tiền lãi (nghìn đồng)  của mỗi ngày trong $14$ ngày được khảo sát ở một quầy bán báo
	\begin{center}
		\begin{tabular}{|l|l|l|l|l|l|l|}
			\hline  $69$&$37$&$39$&$65$&$31$&$33$&$63$\\
			\hline  $51$&$44$&$62$&$33$&$47$&$55$&$42$\\
			 \hline
		\end{tabular}
	\end{center}
	Bảng số liệu  ghép nhóm nào sau đây là \textbf{đúng}?
	\begin{center}
		\begin{tabular}{|c|c|c|c|c|c|}
			\hline\multirow{2}{*}{Bảng $1$}&\begin{tabular}{c}
			Số tiền  lãi\\
			(nghìn đồng) 
		\end{tabular}&{$[30;40)$}&{$[40;50)$}&{$[50;60)$}&{$[60;70)$}\\
		\cline {2-6}&Số  ngày&5&3&2&4\\
		\hline\multirow{2}{*}{Bảng $2$}&\begin{tabular}{c}
		Số tiền lãi\\ 
		(nghìn  đồng)
	\end{tabular}&{$[30;40)$}&{$[40;50)$}&{$[50;60)$}&{$[60;70)$}\\
	 \cline {2-6}&Số ngày&$5$&$3$&$4$&$2$\\
	\hline\multirow{2}{*}{Bảng 3}&\begin{tabular}{c}
	Số  tiền lãi\\
	 (nghìn đồng)
\end{tabular}&{$[30;40)$}&{$[40;50)$}&{$[50;60)$}&{$[60;70)$}\\
\cline {2-6}&Số  ngày&$5$&$2$&$3$&$4$\\
\hline\multirow{2}{*}{Bảng $4$}&\begin{tabular}{c}
Số  tiền lãi\\
(nghìn  đồng)
\end{tabular}&{$[30;40)$}&{$[40;50)$}&{$[50;60)$}&{$[60;70)$}\\
\cline {2-6}&Số  ngày&$3$&$5$&$2$&$4$\\
 \hline
\end{tabular}
\end{center}
 \loigiai{
	Để chuyển mẫu số liệu không ghép nhóm sang  mẫu số liệu ghép nhóm ta làm như sau
	\begin{enumEX}[-]{1}
		\item Bước $1$: Chia miền giá trị của mẫu số liệu thành một số nhóm theo tiêu chí cho  trước.
		\item Bước $2$: Đếm số giá trị của mẫu số liệu thuộc mỗi nhóm và lập  bảng thống kê cho mẫu số liệu ghép nhóm.
	\end{enumEX}
	Ta thực hiện bài  toán như sau\\
	 Khoảng biến thiên là $69-31=38$.\\
	Ta chia thành các  nhóm sau $[30;40)$, $[40;50)$, $[50;60)$, $[60;70)$.\\
	Đếm số giá trị của mỗi nhóm, ta có bảng ghép  nhóm sau
	\begin{center}
		\begin{tabular}{|c|c|c|c|c|}
			\hline\begin{tabular}{c}
			Số tiền lãi\\ 
			(nghìn  đồng)
		\end{tabular}&{$[30;40)$}&{$[40;50)$}&{$[50;60)$}&{$[60;70)$}\\
		\hline Số ngày&$5$&$3$&$2$&$4$\\ 
		\hline 
	\end{tabular}
	\end{center}
	Vậy bảng $1$ là bảng số liệu ghép nhóm của bảng  dữ liệu đã cho.\\
} 
\end{ex}

% nguồn: Ok_TuanLuan_Toan11_Lan2
\begin{ex}%[1D5H1-2]%[ma:fe17f35d]
	Cho bảng khảo sát về cân nặng  của học sinh trong lớp như sau
	\begin{center}
		\begin{tabular}{|c|c|c|c|c|c|}
			\hline Cân nặng (kg)&{$[45;50)$}&{$[50;55)$}&{$[55;60)$}&{$[60;65)$}&{$[65;70)$}\\ 
			\hline Số  học sinh&$2$&$14$&$11$&$10$&$3$\\
			 \hline
		\end{tabular}
	\end{center}
	 Khoảng cân nặng nào mà học sinh chiếm nhiều nhất?
	\loigiai{ 
		\begin{enumEX}[-]{1}
			\item Có $2$ học sinh có cân nặng từ $45$ kg đến dưới $50$ kg. 
			\item Có $14$ học sinh có cân nặng từ  $50$ kg đến dưới $55$ kg.
			\item Có $11$ học sinh có cân nặng từ $55$ kg đến dưới  $60$ kg.
			\item Có $10$ học sinh có cân nặng từ $60$  kg đến dưới $65$ kg.
			\item Có $3$ học sinh có cân nặng từ $65$ kg đến dưới  $70$ kg.
		\end{enumEX}
		Vậy  khoảng cân nặng từ $50$ kg đến dưới $55$ kg chiếm nhiều học sinh nhất.
	 }
\end{ex}

% nguồn: Ok_TuanLuan_Toan11_Lan2
\begin{ex}%[1D5H1-2]%[ma:d991e394]
	\immini[thm]{Bảng biểu diễn mẫu số liệu ghép nhóm được cho dưới dạng bảng tần  số ghép nhóm. Hãy cho biết
	\begin{enumerate}[1)]
		\item Mẫu số liệu đó có bao nhiêu số  liệu; bao nhiêu nhóm;
		 \item Tần số của mỗi nhóm.
	\end{enumerate}}
	{\begin{tabular}{|c|c|}
	\hline 
	Nhóm & Tần số  \\
	\hline 
	\textit{$[0;10)$} &  $8$ \\
	\hline 
	 \textit{$[10;20)$} & $7$\\
	 \hline
	\textit{$[20;30)$}  & $9$ \\
	\hline 
	\textit{$[30;40)$}  & $6$\\
	 \hline
	& \textit{$n=30$}  \\
	\hline 
\end{tabular}}
 \loigiai{
	 Từ Bảng 4, ta thấy
	\begin{enumerate}[1)]
		\item  Mẫu số liệu đó gồm $30$ số liệu và $4$ nhóm.
		\item Tần số của các nhóm $1$, $2$,  $3$, $4$ lần lượt là $8$, $7$, $9$, $6$.
	\end{enumerate}
 }
\end{ex}

% nguồn: Ok_TuanLuan_Toan11_Lan2
\begin{ex}%[1D5H1-2]%[ma:6b47e69a]
	Thành tích chạy 50 m của học sinh lớp 10A của trường THPT  B được thống kê trong bảng sau\\
	\begin{center}
		\begin{tabular}{|l|l|l|l|l|l|l|}
			\hline  $6{,}3$&$6{,}2$&$6{,}5$&$6{,}8$&$6{,}9$&$8{,}2$&$8{,}6$\\
			 \hline $6{,}6$&$6{,}7$&$7{,}0$&$7{,}1$&$7{,}2$&$8{,}3$&$8{,}5$\\
			\hline $7{,}4$&$7{,}3$&$7{,}2$&$7{,}1$&$7{,}0$&$8{,}4$&$8{,}1$\\ 
			 \hline $7{,}1$&$7{,}3$&$7{,}5$&$7{,}5$&$7{,}6$&$8{,}7$&$7{,}6$\\
			\hline $7{,}7$&$7{,}8$&$7{,}5$&$7{,}7$&$7{,}8$&&\\ 
			\hline 
		\end{tabular}
	\end{center}
	Trong lớp $10$A, số học sinh chạy $50$ $m$ hết $7$ giây đến dưới  $8{,}5$ giây chiếm bao nhiêu phần trăm.
	 \loigiai{
		\textbf{\textit{Phương pháp giải:}} Lập bảng phân bố tần số ghép  nhóm từ đó tính được tần suất của mỗi nhóm. Ta chia thành các nhóm có độ dài bằng $5$. Ta sẽ chọn đầu mút phải của nhóm cuối cùng là $90$.\\
		Ta  có bảng ghép nhóm sau
		\begin{center}
			\begin{tabular}{|c|c|c|c|c|c|c|}
				\hline\begin{tabular}{c}
				Thời  gian chạy\\
				 (giây)
			\end{tabular}&{$[6,0;6,5)$}&{$[6,5;7,0)$}&{$[7,0;7,5)$}&{$[7,5;8,0)$}&{$[8,0;8,5)$}&{$[8,5;9,0)$}\\
			\hline Số  học sinh&2&5&10&9&4&3\\
			 \hline
		\end{tabular}
	\end{center}
	\begin{enumEX}[-]{1}
		\item  Tần suất của nhóm $[6{,}0;6{,}5)$ là $f_1=\dfrac{2}{33}\cdot 100\approx 6{,}06\%$.
		\item Tần  suất của nhóm $[6{,}5;7{,}0)$ là $f_2=\dfrac{5}{33}\cdot 100\approx 15{,}15\%$.
		 \item Tần suất của nhóm $[7{,}0;7{,}5)$ là $f_3=\dfrac{10}{33}\cdot 100\approx 30{,}3\%$.
		\item Tần suất của nhóm $[7{,}5;8{,}0)$  là $f_4=\dfrac{9}{33}\cdot 100\approx 27{,}27\%$.
		\item Tần suất  của nhóm $[8{,}0;8{,}5)$ là $f_5=\dfrac{4}{33}\cdot 100\approx 12{,}12\%$.
		\item Tần suất của nhóm $[8{,}5;9{,}0)$ là $f_6=\dfrac{3}{33}\cdot 100\approx  9{,}1\%$.
	\end{enumEX}
	Suy ra số học sinh chạy $50$ m hết $7$  giây đến dưới $8{,}5$ giây chiếm
	 \[30{,}3 \%+27{,}27 \%+12{,}12 \%=69{,}69 \%.\]
 }
\end{ex}

% nguồn: Ok_TuanLuan_Toan11_Lan2
\begin{ex}%[1D5H1-2]%[ma:6e655bcb]
	Một công ty may quần áo đồng phục học sinh cho biết cỡ áo theo chiều cao của học sinh được tính  như sau:
	\begin{center}
		\begin{tabular}{|c|c|c|c|c|c|}
			 \hline
			Chiều cao $(\mathrm{cm})$ & $[150;  160)$ & $[160; 167)$ & $[167; 170)$ & $[170; 175)$ & $[175; 180)$ \\
			 \hline
			Cỡ  áo & $S$ & $M$ & $L$ & $X L$ & $X X L$ \\
			 \hline
		\end{tabular}
	\end{center}
	Công ty muốn ước lượng tỉ lệ các cỡ áo khi may cho học sinh lớp 11 đã đo chiều cao của 36 học sinh nam khối  11 của một trường và thu được mẩu số liệu sau (đơn vị là centimét):
	\begin{center}
		\begin{tabular}{ccccccccccccc}
			160 & 161 & 161 & 162 & 162 & 162 & 163 &  163 & 163 & 164 & 164 & 164 & 164 \\
			165 & 165 &  165 & 165 & 165 & 166 & 166 & 166 & 166 & 167 & 167 & 168 & 168 \\
			168 & 168 & 169 & 169 & 170 &  171 & 171 & 172 & 172 & 174. & & & \\
		\end{tabular}
	\end{center}
	\begin{enumEX}{2}
		\item! Lập  bảng tần số ghép nhóm của mẫu số liệu với các nhóm đã cho ở bảng trên.\\
		\item! Công ty may $500$ áo đồng phục cho học sinh lớp $11$ thì nên may số lượng áo  theo mỗi cỡ là bao nhiêu chiếc?
	\end{enumEX}
	\loigiai{ 
		\begin{enumerate}[a)]
			\item Đếm số giá trị thuộc mỗi nhóm, ta lập được bảng tần số ghép nhóm của mẫu số liệu với các nhóm đã cho ở bảng  trên như sau:
			\begin{center}
				\begin{tabular}{|c|c|c|c|c|c|}
					\hline 
					Chiều cao (cm) & $[150; 160)$ & $[160; 167)$ & $[167; 170)$ & $[170; 175)$  & $[175; 180)$ \\
					 \hline
					Số  học sinh & $0$ & $22$ & $8$ & $6$ & $0$ \\
					 \hline
				\end{tabular}
			\end{center}
			\item Công ty may $500$ áo đồng phục cho học sinh lớp $11$ thì nên may số lượng  áo theo mỗi cỡ như sau:
			\begin{itemize}
				\item Không nên may áo  cỡ $S$ và cỡ $XXL$;
				\item Số lượng áo cỡ $M$ nên may là $\dfrac{22}{36}\cdot 500 \approx  306$ (chiếc);
				\item Số lượng áo cỡ $L$ nên may là $\dfrac{8}{36}\cdot 500 \approx 111$ (chiếc); 
				 \item Số lượng áo cỡ $XL$ nên may là $500-306-111=83$ (chiếc).
			\end{itemize}
		\end{enumerate}
	} 
\end{ex}

% nguồn: Ok_TuanLuan_Toan11_Lan2
\begin{ex}%[1D5H1-2]%[ma:79e4bdaf]
	Bảng dưới đây biểu diễn mẫu số liệu ghép  nhóm được cho dưới dạng bảng tần số ghép nhóm. Hãy cho biết:
	\begin{enumEX}{2}
		\item! Mẫu số  liệu đó có bao nhiêu số liệu; bao nhiêu nhóm;
		\item! Tần  số của mỗi nhóm.
	\end{enumEX}
	\begin{center}
		\begin{tabular}{|c|c|}
			 \hline
			 Nhóm & Tần số \\
			 \hline
			$[0; 10)$  & $8$ \\
			 \hline
			 $[10; 20)$ & $7$ \\
			\hline 
			$[20; 30)$ &  $9$ \\
			 \hline
			$[30; 40)$  & $6$ \\
			\hline 
			& $n=30$  \\
			 \hline
		\end{tabular}
	\end{center}
	\loigiai{ 
		Từ bảng  trên, ta thấy:
		\begin{enumEX}{2}
			\item! Mẫu số liệu đó gồm $30$  số liệu và $4$ nhóm.
			\item! Tần số của các nhóm $1$, $2$, $3$, $4$ lần  lượt là $8$, $7$, $9$, $6$.
		\end{enumEX}
	 }
\end{ex}

% nguồn: Ok_TuanLuan_Toan11_Lan2
\begin{ex}%[1D5H1-2]%[ma:323d64fc]
	Trẻ sơ sinh được xem là nhẹ cân nếu cân nặng khi sinh dưới $2$ kg, là thừa cân nếu cân nặng khi sinh trên $4$ kg, là có cân nặng trung bình nếu cân nặng khi sinh từ $2$ kg đến $4$ kg. Thống kê cân nặng (tính theo kg) của $15$ trẻ sơ sinh tại một bệnh viện cho kết quả như sau: \[\begin{array}{ccccccccccccccc}3{,}4 & 2{,}7 & 1{,}9 & 3{,}5 & 3{,}3 & 2{,}8 & 4{,}2 & 2{,}6 & 2{,}8 & 3{,}0 & 3{,}7 & 3{,}9 & 4{,}1 & 2{,}7 & 2{,}5.
\end{array}\]
\begin{enumEX}{2}
	\item! Tìm số trẻ nhẹ cân, thừa cân, có cân nặng trung bình  trong $15$ trẻ sơ sinh trên.
	\item! Xây dựng  mẫu số liệu ghép nhóm cho mẫu số liệu trên.
\end{enumEX}
 \loigiai{
	\begin{enumerate}[a)]
		\item Số trẻ nhẹ cân, cân nặng  trung bình, thừa cân tương ứng là $1$, $12$, $2$.
		\item Mẫu số liệu  ghép nhóm:
		\begin{center}
			\begin{tabular}{|c|c|c|c|}
				 \hline
				Cân nặng (kg)& Dưới $2$ & $[2;  4]$ & Trên $4$ \\
				\hline 
				Số trẻ & $1$ &  $12$ & $2$ \\
				 \hline
			\end{tabular}
		\end{center}
	\end{enumerate}
 }
\end{ex}

% nguồn: Ok_TuanLuan_Toan11_Lan2
\begin{ex}%[1D5H1-2]%[ma:819a8f62]
	Từ $1 / 7 / 2019$, dựa trên thu nhập bình quân đầu người (kí hiệu là $GNIPC$, tính theo đô la Mĩ), Ngân hàng Thế giới xác định một nền kinh tế ở mức thu nhập thấp nếu GNIPC nhỏ hơn $1\,026$, ở mức thu nhập dưới trung bình nếu GNIPC từ $1\,026$ đến dưới $3\,996$, ở mức thu nhập trên trung bình nếu GNIPC từ $3\,996$ đến dưới $12\,376$ và ở mức thu nhập cao nếu GNIPC từ $12\,376$ trở lên (Theo Ngân hàng Thế giới). Thu nhập bình quân đầu người của một số nền kinh tế thuộc khu  vực châu Á Thái Bình Dương năm $2021$ được cho như sau:
	\begin{center}
		\begin{tabular}{cccccccccc}
			102\,450&70\,700&67\,580&55\,290& 47\,490& 45\,440& 44\,570& 28\,730&19\,170\\ 
			18\,530&16\,520&13\,790&12\,904&11\,090&11\,040&10\,440&9\,450&8\,150\\ 
			 7\,220& 6\,960&5\,800&4\,430&4\,340&4\,280&4\,230& 2\,100.
		\end{tabular}
	\end{center}
	\begin{enumEX}{2}
		\item! Ghép nhóm mẫu  số liệu trên theo mức thu nhập của nền kinh tế.
		\item! $GNIPC$ của Việt Nam năm $2021$ là  $11\,040$. Nền kinh tế Việt Nam được xếp Ở mức nào?
	\end{enumEX}
	 \loigiai{
		\begin{enumerate}[a)]
			\item Mẫu số  liệu ghép nhóm:
			\begin{center}
				\begin{tabular}{|c|c|c|c|}
					\hline 
					GNIPC & $[1\,026; 3\,996)$ & $[3\,996; 12\,376)$ & Từ $12\,376$ trở lên  \\
					 \hline
					Số nền kinh tế  & $1$ & $12$ & $13$ \\
					\hline 
				\end{tabular}
			\end{center}
			\item Nền kinh tế Việt Nam được xếp  ở mức thu nhập trên trung bình.
		\end{enumerate}
	} 
\end{ex}

% nguồn: Ok_TuanLuan_Toan11_Lan2
\begin{ex}%[1D5H1-2]%[ma:a64cbfb6]
	Trong các mẫu số liệu sau, mẫu nào là mẫu số liệu ghép nhóm? Đọc và giải thích mẫu số liệu ghép nhóm  đó.
	\begin{enumerate}[a)]
		\item Số tiền mà sinh viên  chi cho thanh toán cước điện thoại trong tháng.
		\begin{center}
			\begin{tabular}{|c|c|c|c|c|c|}
				\hline 
				\begin{tabular}{c}
					Số tiền (nghìn  đồng) \\
				\end{tabular}& $[0; 50)$ & $[50; 100)$ & $[100; 150)$ & $[150; 200)$ & $[200; 250)$ \\
				\hline 
				Số  sinh viên & $5$ & $12$ & $23$ & $17$ & $3$ \\
				 \hline
			\end{tabular}
		\end{center}
		\item Thống kê nhiệt độ tại  một địa điểm trong $40$ ngày, ta có bảng số liệu sau:
		\begin{center}
			\begin{tabular}{|c|c|c|c|c|}
				 \hline
				Nhiệt độ $\left({}^{\circ}\mathrm{C}\right)$ & $[19; 22)$ & $[22; 25)$ & $[25; 28)$ &  $[28; 31)$ \\
				 \hline
				Số ngày & $7$ & $15$  & $12$ & $6$ \\
				 \hline
			\end{tabular}
		\end{center}
	\end{enumerate}
	 \loigiai{
		\begin{enumerate}[a)]
			\item Mẫu số  liệu đã cho là mẫu số liệu ghép nhóm.
			Mẫu số liệu này mô tả về số tiền mà sinh viên chi cho thanh toán cước điện thoại trong  tháng, gồm có $5$ nhóm. Cụ thể:
			\begin{itemize}
				\item Nhóm thanh toán với số tiền từ $0$ đến dưới $50$  nghìn đồng, có $5$ sinh viên;
				 \item Nhóm thanh toán với số tiền từ $50$ đến dưới $100$ nghìn đồng, có $12$ sinh viên;
				\item Nhóm thanh toán với số tiền từ $100$  đến dưới $150$ nghìn đồng, có $23$ sinh viên;
				\item Nhóm thanh toán với số tiền từ $150$ đến dưới $200$ nghìn đồng, có $17$  sinh viên;
				\item Nhóm thanh toán với số tiền từ $200$ đến dưới $250$ nghìn đồng, có $3$ sinh viên; 
			\end{itemize}
			\item Mẫu số liệu đã cho  là mẫu số liệu ghép nhóm.
			Mẫu số liệu này mô tả về nhiệt độ tại một địa điểm trong $40$ ngày, gồm $4$ nhóm nhiệt độ: từ $19^{\circ}\mathrm{C}$ đến dưới $22^{\circ}\mathrm{C}$; từ $22^{\circ}\mathrm{C}$ đến dưới $25^{\circ}\mathrm{C}$; từ $25^{\circ}\mathrm{C}$ đến dưới $28^{\circ}\mathrm{C}$;  từ $28^{\circ}\mathrm{C}$ đến dưới $31^{\circ}\mathrm{C}$. Cụ thể:
			\begin{itemize}
				\item Có $7$ ngày có nhiệt độ từ $19^{\circ}\mathrm{C}$ đến dưới $22^{\circ}\mathrm{C}$; 
				 \item Có $15$ ngày có nhiệt độ từ $22^{\circ}\mathrm{C}$ đến dưới $25^{\circ}\mathrm{C}$;
				\item Có $12$ ngày có nhiệt độ từ $25^{\circ}\mathrm{C}$  đến dưới $28^{\circ}\mathrm{C}$;
				\item Có $6$ ngày có  nhiệt độ từ $28^{\circ}\mathrm{C}$ đến dưới $31^{\circ}\mathrm{C}$.
			\end{itemize}
		\end{enumerate}
	 }
\end{ex}

% nguồn: Ok_TuanLuan_Toan11_Lan2
\begin{ex}%[1D5H1-2]%[ma:e1641fb4]
	Thời gian hoàn thành bài kiểm tra Toán  $45$ phút của các bạn trong lớp được cho như sau
	\begin{center}
		\begin{tabular}{|c|c|c|c|c|}
			 \hline
			Thời gian (phút) & $[25;30)$ & $[30;35)$ & $[35;40)$ & $[40;45]$ \\ 
			 \hline
			Số học sinh & $2$  & $7$ & $10$ & $25$ \\
			 \hline
		\end{tabular}
	\end{center}
	\begin{enumerate}[1)]
		\item  Nêu các nhóm số liệu và tần số tương ứng.
		 \item Có bao nhiêu học sinh hoàn thành bài kiểm tra trước khi hết giờ trên $5$ phút?.
	\end{enumerate}
	 \loigiai{
		\begin{enumerate}[1)]
			\item Các nhóm số liệu $[25;30)$, $[30;35)$, $[35;40)$, $[40;45)$ với tần  số tương ứng là $2$, $7$, $10$, $25$.
			 \item Số học sinh hoàn thành bài kiểm tra trước khi hết giờ ít nhất $5$ phút là $2+7+10=19$.
		\end{enumerate}
	} 
\end{ex}

% nguồn: Ok_TuanLuan_Toan11_Lan2
\begin{ex}%[1D5H1-2]%[ma:d0039fb8]
	Một công ty may quần áo đồng phục học sinh cho  biết cỡ áo theo chiều cao của học sinh được tính như sau
	\begin{center}
		\begin{tabular}{|c|c|c|c|c|c|}
			\hline 
			Chiều cao (cm) & $[150;160)$ & $[160;167)$  & $[167;170)$ & $[170;175)$ & $[175;180)$ \\
			 \hline
			Cỡ áo & S & M & L & XL & XXL  \\
			 \hline
		\end{tabular}
	\end{center}
	Công ty muốn ước lượng  tỉ lệ các cỡ áo khi may cho học sinh lớp $11$ đã đo chiều cao của $36$ học sinh nam khối $11$ của một trường và thu được mẫu số liệu sau (đơn vị là centimét):
	\begin{center}
		\begin{tabular}{|c|c|c|c|c|c|c|c|c|c|c|c|}
			 \hline
			$160$ & $161$ & $161$ & $162$ & $162$ & $162$ & $163$ & $163$ & $163$ & $164$ & $164$ &  $164$ \\
			\hline 
			$165$ & $165$ & $165$  & $165$ & $165$ & $166$ & $166$ & $166$ & $166$ & $167$ & $167$ & $168$ \\
			 \hline
			$168$ & $168$ & $169$ & $169$ & $170$ & $171$ & $171$ & $172$ & $172$ &  $174$ & $164$ & $168$\\
			\hline 
		\end{tabular}
	\end{center}
	\begin{enumerate}
		\item Lập bảng tần số ghép nhóm của mẫu số liệu với các nhóm  đã cho ở bảng trên.
		\item Công ty may $500$ áo đồng phục cho học sinh lớp $11$ thì nên may số lượng  áo theo mỗi cở là bao nhiêu chiếc?
	\end{enumerate}
	 \loigiai{
		\begin{enumerate}
			\item Bảng tần số ghép nhóm của mẫu số liệu với các nhóm đã  cho là
			\begin{center}
				\begin{tabular}{|c|c|c|c|}
					\hline 
					Chiều cao (cm)  & $[160;167)$ & $[167;170)$ & $[170;175)$ \\
					 \hline
					Số học sinh & $22$ & $8$ & $6$ \\ 
					\hline 
				\end{tabular}
			\end{center}
			\item Tỉ lệ  học sinh mặc vừa cỡ M: $\dfrac{22}{36} \cdot 100=61{,}11\%$;\\
			Tỉ lệ học sinh  mặc vừa cỡ L: $\dfrac{8}{36} \cdot 100=22{,}22\%$;\\
			Tỉ lệ học sinh  mặc vừa cỡ XL: $\dfrac{6}{36} \cdot 100=16{,}67\%$;\\
			Số lượng áo cỡ M nên may: $61{,}11\% \cdot 500=306$;\\ 
			Số  lượng áo cỡ L nên may: $22{,}22\% \cdot 500=111$;\\
			Số lượng  áo cỡ XL nên may: $16{,}67\% \cdot 500=83$.
	\end{enumerate}}
\end{ex}

% nguồn: Ok_TuanLuan_Toan11_Lan2
\begin{ex}%[1D5H1-2]%[ma:90cfae7c]
	Trong các mẫu số liệu sau, mẫu nào là  mẫu số liệu ghép nhóm? Đọc và giải thích mẫu số liệu ghép nhóm đó.
	\begin{enumerate}
		\item Thống kê số  tiền mà sinh viên chi cho thanh toán cước điện thoại trong tháng
		\begin{center}
			\begin{tabular}{|c|c|c|c|c|c|}
				\hline 
				Số tiền (nghìn đồng) & $[0;50)$ & $[50;100)$ &  $[100;150)$ & $[150;200)$ & $[200;250)$ \\
				 \hline
				Số sinh viên  & $5$ & $12$ & $23$ & $17$ & $3$ \\
				\hline 
			\end{tabular}
		\end{center}
		\item Thống kê nhiệt độ tại một địa điểm trong $40$ ngày, ta có  bảng số liệu sau
		\begin{center}
			\begin{tabular}{|c|c|c|c|c|}
				 \hline
				Nhiệt độ ($^\circ\text{C}$) & $[19;22)$ & $[22;25)$ & $[25;28)$ & $[28;31)$ \\ 
				 \hline
				Số ngày &  $5$ & $12$ & $23$ & $17$ \\
				 \hline
			\end{tabular}
		\end{center}
	\end{enumerate}
	\loigiai{Cả hai mẫu số liệu đều là mẫu số liệu ghép  nhóm.
		\begin{itemize}
			\item Với mẫu số liệu của  câu a)
			\begin{itemize}
				\item Có $5$ sinh viên chi dưới $50$ nghìn đồng cho việc thanh toán cước điện  thoại trong tháng;
				\item Có $12$ sinh viên  chi từ $50$ đến dưới $100$ nghìn đồng cho việc thanh toán cước điện thoại trong tháng;
				\item Có $23$ sinh viên chỉ từ $100$ đến dưới  $150$ nghìn đồng cho việc thanh toán cước điện thoại trong tháng;
				\item Có $17$ sinh viên chỉ từ $150$ đến dưới $200$ nghìn đồng cho việc thanh toán cước điện thoại trong tháng; 
				\item Có $3$  sinh viên chi từ $200$ đến dưới $250$ nghìn đồng cho việc thanh toán cước điện thoại trong tháng.
			\end{itemize}
			Kết luận: đa số sinh viên chi từ $100$ đến dưới $150$ nghìn  đồng mỗi tháng cho cước điện thoại và có ít sinh viên chi trên $200$ nghìn đồng cho cước điện thoại mỗi tháng.
			\item Với mẫu số  liệu của câu b)
			\begin{itemize}
				\item Có $7$ ngày có nhiệt  độ từ $19^\circ \text{C}$ đến dưới $22^\circ \text{C}$;
				\item Có $15$ ngày có nhiệt độ từ $22^\circ \text{C}$  đến dưới $25^\circ \text{C}$;
				\item Có $12$ ngày có nhiệt độ từ $25^\circ  \text{C}$ đến dưới $28^\circ \text{C}$;
				 \item Có $6$ ngày có nhiệt độ từ $28^\circ \text{C}$ đến dưới $31^\circ \text{C}$.
			\end{itemize}
			Kết luận: nhiệt  độ tại địa điểm này chủ yếu dao động từ $25^\circ\text{C}$ đến $31^\circ\text{C}$. Những ngày mát mẻ dưới $22^\circ\text{C}$ là khá ít.
		\end{itemize}
	} 
\end{ex}

% nguồn: Ok_TuanLuan_Toan11_Lan2
\begin{ex}%[1D5H1-2]%[ma:1e991e9f]
	Một  cửa hàng đã thống kê số ba lô bán được mỗi ngày trong tháng $9$ với kết quả cho như sau:
	$$ 
	\begin{array}{ccccccccccccccc}
		12 & 29 & 12 & 19 & 15 & 21 & 19 &  29 & 28 & 12 & 15 & 25 & 16 & 20 & 29 \\
		21 & 12 & 24 & 14 & 10 & 12 & 10 & 23 & 27 & 28 & 18  & 16 & 10 & 20 & 21.
	\end{array}
	 $$
	Hãy chia mẫu số liệu trên thành $5$ nhóm, lập bảng tần số ghép nhóm, hiệu chỉnh bảng tần số  ghép nhóm và xác định giá trị đại diện cho mỗi nhóm.
	 \loigiai{
		Khoảng biến thiên của mẫu  số liệu trên là $R=29-10=19$.\\
		Độ  dài mỗi nhóm $L>\dfrac{R}{k}=\dfrac{19}{5}=3,8$.\\
		Ta chọn $L=4$ và chia dữ liệu thành các nhóm: 
		 $[10; 14)$, $[14; 18)$, $[18; 22)$, $[22; 26)$, $[26; 30)$.\\
		Khi đó  ta có bảng tần số ghép nhóm sau:
		\begin{center}
			\begin{tabular}{|c|c|c|c|c|c|}
				 \hline
				Số ba lô & $[10; 14)$ & $[14; 18)$ & $[18; 22)$ & $[22; 26)$ &  $[26; 30)$ \\
				\hline 
				Giá trị đại diện & $12$ & $16$  & $20$ & $24$ & $28$ \\
				 \hline
				Số ngày & $8$ & $5$ & $8$ & $3$  & $6$ \\
				 \hline
			\end{tabular}
		\end{center}
	 }
\end{ex}

% nguồn: Ok_TuanLuan_Toan11_Lan1
\begin{bt}%[1D5V1-2]%[ma:e149d5d1]
	Một cửa hàng bán $3$ loại hoa quả nhập khẩu: Lê, Dưa vàng và Bưởi với số liệu tính toán được cho bởi bảng (trong một quý) sau khi giảm giá mỗi  loại lần lượt là $x$, $y$, $z$ trên $1$ kg
	\begin{center}
		\begin{tabular}{|c|c|c|c|}
			\hline  Loại quả&Lê&Dưa vàng&Bưởi\\
			\hline Giá bán  (nghìn/kg)&$200-x$&$300-y$&$400-z$\\
			\hline Số lượng bán (kg)&$200+x$&$300+y$&$400+z$\\ 
			\hline 
		\end{tabular}
	\end{center}
	Biết rằng $x+y+z=90$ (nghìn). Tính giá trị $x$, $y$, $z$ để lợi nhuận bình  quân của $1$ kg hoa quả đạt được cao nhất.
	 \loigiai{
		Do khối lượng hoa  quả bán được là $200+x+300+y+400+z=990$ là cố định, vì thế bình quân mỗi kg hoa quả có giá cao nhất khi tổng số tiền thu được là cao nhất.\\
		Tổng số tiền thu được là 
		 $$P=(200-x)(200+x)+(300-y)(300+y)+(400-z)(400+z)=290000-({x^2}+{y^2}+{z^2}).$$
		Ta có bất đẳng thức ${x^2}+{y^2}+{z^2}\ge\dfrac{1}{3}{(x+y+z)^2}=2700$, từ đó $P\le 287\,300$.\\ 
		Vậy $P$ lớn nhất khi dấu bằng xảy ra tức  là $x=y=z=30$ (nghìn).
	 }
\end{bt}

% nguồn: Ok_TuanLuan_Toan11_Lan2
\begin{ex}%[1D5V1-2]%[ma:3fb47fbf]
	Một cửa hàng bán $3$ loại hoa quả nhập khẩu: Lê, Dưa vàng và Bưởi với số liệu tính toán được cho bởi bảng (trong một quý)  sau khi giảm giá mỗi loại lần lượt là $x$, $y$, $z$ trên $1$ kg
	\begin{center}
		\begin{tabular}{|c|c|c|c|}
			\hline Loại quả&Lê&Dưa  vàng&Bưởi\\
			 \hline Giá bán (nghìn/kg)&$200-x$&$300-y$&$400-z$\\
			\hline Số lượng  bán (kg)&$200+x$&$300+y$&$400+z$\\
			\hline 
		\end{tabular}
	\end{center}
	Biết rằng $x+y+z=90$ (nghìn). Tính giá trị  $x$, $y$, $z$ để lợi nhuận bình quân của $1$ kg hoa quả đạt được cao nhất.
	 \loigiai{
		Do khối lượng hoa quả bán được là $200+x+300+y+400+z=990$ là cố định, vì thế bình quân mỗi kg hoa quả có giá cao nhất khi tổng số tiền thu được là  cao nhất.\\
		Tổng số tiền thu được là 
		$$P=(200-x)(200+x)+(300-y)(300+y)+(400-z)(400+z)=290000-({x^2}+{y^2}+{z^2}).$$ 
		Ta có bất đẳng thức ${x^2}+{y^2}+{z^2}\ge\dfrac{1}{3}{(x+y+z)^2}=2700$, từ  đó $P\le 287\,300$.\\
		Vậy $P$ lớn nhất khi dấu bằng xảy  ra tức là $x=y=z=30$ (nghìn).
	 }
\end{ex}
